% Fonctions numériques d’une variable réelle — Mathématiques Tle Tech — Cours signé OnBuch+
% Programme officiel MINESEC. Compiler avec : tectonic N08-fonctions-numeriques-variable-reelle.tex
\newcommand{\DOCMATIERE}{Mathématiques}
\newcommand{\DOCNIVEAU}{Tle Tech}
\newcommand{\DOCMODULE}{Mathématiques – classe de Terminale technique}
\newcommand{\DOCLECON}{N08}
\newcommand{\DOCTITRE}{Fonctions numériques d’une variable réelle}
% OnBuch+ — préambule « cours signés » ENRICHI (compilé avec tectonic / XeLaTeX)
% Système visuel « Pop » d'OnBuch+ : fond crème (jamais blanc), bordures encre
% épaisses, ombres « dures » décalées, titres Archivo Black en capitales.
% Version MATHÉMATIQUES : graphes (pgfplots/tikz), boîtes pédagogiques
% (définition, propriété, méthode, attention, le savais-tu, exercice résolu,
% exercice, TP, bilan), page de garde et signature OnBuch+ ancrée en bas.
%
% Macros attendues dans main.tex AVANT \input{preamble} :
%   \DOCMATIERE  \DOCNIVEAU  \DOCMODULE  \DOCLECON (numéro)  \DOCTITRE
\documentclass[11pt]{article}

\usepackage{fontspec}
\usepackage[french]{babel}
\usepackage[a4paper,margin=2cm,top=3.4cm,bottom=2.8cm,headheight=1.4cm,footskip=1.3cm]{geometry}
\usepackage{amsmath,amssymb,amsfonts}
\usepackage{mathtools} % \xmapsto, \coloneqq... souvent utilisés par les modèles
\usepackage{cancel} % simplifications barrées
% \abs{z} (module, valeur absolue) et \norm{u} (norme), utilisés par les modèles
\providecommand{\abs}{}\let\abs\relax\DeclarePairedDelimiter{\abs}{\lvert}{\rvert}
\providecommand{\norm}{}\let\norm\relax\DeclarePairedDelimiter{\norm}{\lVert}{\rVert}
\usepackage[table]{xcolor} % [table] : fournit aussi \rowcolors (alternance de lignes)
\usepackage{pagecolor}
\usepackage{tikz}
\usetikzlibrary{arrows.meta,positioning,calc,shapes.geometric,shapes.misc,shapes.arrows,decorations.pathmorphing,decorations.pathreplacing,decorations.markings,patterns,shadows,mindmap,trees,fit,backgrounds,matrix,intersections,angles,quotes}
\usepackage{pgfplots}
\usepackage{pgf-pie} % diagrammes circulaires \pie (statistiques)
\pgfplotsset{compat=1.18}
\usepgfplotslibrary{fillbetween,groupplots} % aires entre courbes (calcul intégral)
\usepackage{enumitem}
\usepackage{fancyhdr}
% [most] charge toutes les bibliothèques tcolorbox (listings, minted, raster,
% documentation...) et gonfle inutilement la mémoire TeX sur un document long
% (~300 boîtes) : on ne charge que ce qui est réellement utilisé.
\usepackage[skins,breakable]{tcolorbox}
\usepackage{graphicx}
\usepackage{colortbl}
\usepackage{booktabs}
\usepackage{tabularx}
\usepackage{diagbox} % case d'en-tête diagonale des tableaux à double entrée
% Colonnes extensibles centrée / gauche / droite (C, L, R), souvent utilisées par les modèles
\newcolumntype{C}{>{\centering\arraybackslash}X}
\newcolumntype{L}{>{\raggedright\arraybackslash}X}
\newcolumntype{R}{>{\raggedleft\arraybackslash}X}
\usepackage{array}
\usepackage{multicol}
\usepackage{siunitx}
% parse-numbers=false : accepte aussi \SI{2,5 \times 10^{-3}}{...}
\sisetup{locale=FR,output-decimal-marker={,},group-minimum-digits=4,parse-numbers=false}
\DeclareSIUnit\ohm{\ensuremath{\symup{\Omega}}} % U+2126 absent de la police
\usepackage{qrcode}
% \degree : commande fréquemment utilisée par les modèles pour un angle ou une
% température en dehors de \SI{}{} (non fournie par les paquets chargés).
\providecommand{\degree}{^{\circ}}
% \qed (fin de démonstration), utilisé par les modèles sans amsthm
\providecommand{\qed}{\hfill$\square$}
% \barre : double barre des valeurs interdites dans un tableau de variations (tkz-tab)
\providecommand{\barre}{\ensuremath{\|}}
% environnement proof (amsthm non chargé) utilisé par les modèles
\ifdefined\proof\else\newenvironment{proof}[1][Démonstration]{\par\noindent\textit{#1.}\ }{\qed\par}\fi
\usepackage{lastpage}
\usepackage{hyperref}
\hypersetup{hidelinks}

% ============================================================
% Palette « Pop » OnBuch+ — mêmes hex que la classe P (plus_ui.dart)
% ============================================================
\definecolor{poporange}{HTML}{F59321}
\definecolor{popdark}{HTML}{E07A0C}
\definecolor{popcream}{HTML}{FFF6E8}
\definecolor{popink}{HTML}{1C1714}
\definecolor{poppurple}{HTML}{7A5AE0}
\definecolor{popgreen}{HTML}{2E9E6A}
\definecolor{poppink}{HTML}{E0457B}
\definecolor{popblue}{HTML}{2F7DE1}
\definecolor{popgold}{HTML}{C98A12}
\definecolor{popmuted}{HTML}{8A7F76}
\definecolor{popline}{HTML}{EADFD2}
\definecolor{popgray}{HTML}{8A7F76}
\colorlet{popgrey}{popgray}
% Alias tolérants (le modèle nomme parfois les couleurs différemment)
\colorlet{popwhite}{white}
\colorlet{popblack}{popink}
\colorlet{popred}{poppink}
\colorlet{popyellow}{popgold}
% Teintes claires pour fonds de tableaux / zones
\colorlet{popblueL}{popblue!10!white}
\colorlet{poporangeL}{poporange!14!white}
\colorlet{popgreenL}{popgreen!12!white}
\colorlet{poppurpleL}{poppurple!12!white}
\colorlet{poppinkL}{poppink!12!white}
\colorlet{popgoldL}{popgold!14!white}

\pagecolor{popcream}

% ============================================================
% Typographie
% ============================================================
\defaultfontfeatures{Ligatures=TeX}
\setmainfont{PlusJakartaSans-Medium.ttf}[Path=../../../fonts/,BoldFont=PlusJakartaSans-Bold.ttf]
% Maths en sans-serif (Fira Math) avec chiffres et lettres droites de Plus
% Jakarta, pour que les formules se fondent dans le texte.
\usepackage[math-style=ISO,bold-style=ISO]{unicode-math}
\setmathfont{FiraMath-Regular.otf}[Path=../../../fonts/]
\setmathfont{PlusJakartaSans-Medium.ttf}[Path=../../../fonts/,range={up/num,up/{Latin,latin},\mathup}]
\usepackage{icomma} % virgule décimale sans espace en mode math
% Caractères mathématiques Unicode absents des polices de texte (Plus Jakarta,
% Archivo Black) : rendus par leur équivalent mathématique, en texte comme en
% formule — sinon ils s'affichent en carré vide (ex. « Arithmétique dans ℤ »).
\usepackage{newunicodechar}
\newunicodechar{ℝ}{\ensuremath{\mathbb{R}}}
\newunicodechar{ℤ}{\ensuremath{\mathbb{Z}}}
\newunicodechar{ℂ}{\ensuremath{\mathbb{C}}}
\newunicodechar{ℕ}{\ensuremath{\mathbb{N}}}
\newunicodechar{ℚ}{\ensuremath{\mathbb{Q}}}
\newunicodechar{𝔻}{\ensuremath{\mathbb{D}}}
\newunicodechar{α}{\ensuremath{\alpha}}
\newunicodechar{β}{\ensuremath{\beta}}
\newunicodechar{γ}{\ensuremath{\gamma}}
\newunicodechar{δ}{\ensuremath{\delta}}
\newunicodechar{ε}{\ensuremath{\varepsilon}}
\newunicodechar{θ}{\ensuremath{\theta}}
\newunicodechar{λ}{\ensuremath{\lambda}}
\newunicodechar{μ}{\ensuremath{\mu}}
\newunicodechar{σ}{\ensuremath{\sigma}}
\newunicodechar{τ}{\ensuremath{\tau}}
\newunicodechar{φ}{\ensuremath{\varphi}}
\newunicodechar{ψ}{\ensuremath{\psi}}
\newunicodechar{ω}{\ensuremath{\omega}}
\newunicodechar{Σ}{\ensuremath{\Sigma}}
% majuscules produites par \MakeUppercase dans les titres (α-aminés -> Α)
\newunicodechar{Α}{\ensuremath{\alpha}}
\newunicodechar{Β}{\ensuremath{\beta}}
\newunicodechar{Γ}{\ensuremath{\Gamma}}
\newunicodechar{Δ}{\ensuremath{\Delta}}
\newunicodechar{Ε}{\ensuremath{\varepsilon}}
\newunicodechar{Θ}{\ensuremath{\Theta}}
\newunicodechar{Λ}{\ensuremath{\Lambda}}
\newunicodechar{Μ}{\ensuremath{\mu}}
\newunicodechar{Τ}{\ensuremath{\tau}}
\newunicodechar{Φ}{\ensuremath{\Phi}}
\newunicodechar{Ψ}{\ensuremath{\Psi}}
\newunicodechar{Ω}{\ensuremath{\Omega}}
\newunicodechar{↦}{\ensuremath{\mapsto}}
\newunicodechar{∈}{\ensuremath{\in}}
\newunicodechar{∉}{\ensuremath{\notin}}
\newunicodechar{∧}{\ensuremath{\wedge}}
\newunicodechar{⇔}{\ensuremath{\Leftrightarrow}}
\newunicodechar{⟺}{\ensuremath{\Longleftrightarrow}}
\newunicodechar{⇒}{\ensuremath{\Rightarrow}}
\newunicodechar{⟹}{\ensuremath{\Longrightarrow}}
\newunicodechar{∘}{\ensuremath{\circ}}
\newunicodechar{∪}{\ensuremath{\cup}}
\newunicodechar{∩}{\ensuremath{\cap}}
\newunicodechar{≡}{\ensuremath{\equiv}}
\newunicodechar{⊂}{\ensuremath{\subset}}
\newunicodechar{⊥}{\ensuremath{\perp}}
\newunicodechar{∃}{\ensuremath{\exists}}
\newunicodechar{∀}{\ensuremath{\forall}}
\newunicodechar{∛}{\ensuremath{\sqrt[3]{\,}}}
\newunicodechar{∜}{\ensuremath{\sqrt[4]{\,}}}
\newunicodechar{⃗}{\ensuremath{{}^{\to}}}
\newunicodechar{ⁿ}{\ifmmode^{n}\else\textsuperscript{\ensuremath{n}}\fi}
\newunicodechar{ˣ}{\ifmmode^{x}\else\textsuperscript{\ensuremath{x}}\fi}
\newunicodechar{ᵇ}{\ifmmode^{b}\else\textsuperscript{\ensuremath{b}}\fi}
\newunicodechar{ʳ}{\ifmmode^{r}\else\textsuperscript{\ensuremath{r}}\fi}
\newunicodechar{ᵘ}{\ifmmode^{u}\else\textsuperscript{\ensuremath{u}}\fi}
\newunicodechar{ʸ}{\ifmmode^{y}\else\textsuperscript{\ensuremath{y}}\fi}
\newunicodechar{ᵃ}{\ifmmode^{a}\else\textsuperscript{\ensuremath{a}}\fi}
\newunicodechar{ᵉ}{\ifmmode^{e}\else\textsuperscript{\ensuremath{e}}\fi}
\newunicodechar{ᵏ}{\ifmmode^{k}\else\textsuperscript{\ensuremath{k}}\fi}
\newunicodechar{ᵖ}{\ifmmode^{p}\else\textsuperscript{\ensuremath{p}}\fi}
\newunicodechar{ᴺ}{\ifmmode^{N}\else\textsuperscript{\ensuremath{N}}\fi}
\newunicodechar{ᶜ}{\ifmmode^{c}\else\textsuperscript{\ensuremath{c}}\fi}
\newunicodechar{ʰ}{\ifmmode^{h}\else\textsuperscript{\ensuremath{h}}\fi}
\newunicodechar{ᵠ}{\ifmmode^{q}\else\textsuperscript{\ensuremath{q}}\fi}
\newunicodechar{ₙ}{\ifmmode_{n}\else\textsubscript{\ensuremath{n}}\fi}
\newunicodechar{ₐ}{\ifmmode_{a}\else\textsubscript{\ensuremath{a}}\fi}
\newunicodechar{ᵢ}{\ifmmode_{i}\else\textsubscript{\ensuremath{i}}\fi}
\newunicodechar{ᵣ}{\ifmmode_{r}\else\textsubscript{\ensuremath{r}}\fi}
\newunicodechar{ₖ}{\ifmmode_{k}\else\textsubscript{\ensuremath{k}}\fi}
\newunicodechar{ᵦ}{\ifmmode_{\beta}\else\textsubscript{\ensuremath{\beta}}\fi}
\newunicodechar{ₓ}{\ifmmode_{x}\else\textsubscript{\ensuremath{x}}\fi}
\newfontfamily\popdisplay{ArchivoBlack-Regular.ttf}[Path=../../../fonts/]
\newfontfamily\poplogo{SpaceGrotesk-Bold.ttf}[Path=../../../fonts/]
\newfontfamily\popsemi{PlusJakartaSans-SemiBold.ttf}[Path=../../../fonts/]
\newfontfamily\popxbold{PlusJakartaSans-ExtraBold.ttf}[Path=../../../fonts/]
\newfontfamily\popxboldls{PlusJakartaSans-ExtraBold.ttf}[Path=../../../fonts/,LetterSpace=3.0]

\setlist[enumerate]{leftmargin=1.7em,itemsep=5pt,topsep=4pt}
\setlist[itemize]{leftmargin=1.7em,itemsep=4pt,topsep=4pt,label={\color{poporange}\textbullet}}
\setlength{\parindent}{0pt}
\setlength{\parskip}{5pt}
\renewcommand{\arraystretch}{1.25}

% pgfplots : style Pop par défaut
\pgfplotsset{
  every axis/.append style={axis line style={popink,line width=1pt},tick style={popink},
    grid style={popline,line width=0.5pt},label style={font=\small},tick label style={font=\footnotesize}},
  popcurve/.style={poporange,line width=1.8pt,no markers,smooth},
}
% Même style disponible dans un \draw TikZ simple (hors pgfplots)
\tikzset{popcurve/.style={poporange,line width=1.8pt}}

% ============================================================
% Pill / squiggle
% ============================================================
\newcommand{\poppill}[3]{%
  \tikz[baseline=-0.55ex]\node[draw=popink,line width=1.2pt,rounded corners=2.6pt,%
    fill=#2,inner xsep=6.5pt,inner ysep=3pt]{%
    \popxboldls\fontsize{7}{7}\selectfont\color{#3}\MakeUppercase{#1}};%
}
\newcommand{\popsquiggle}[1]{%
  \tikz[baseline=-0.4ex]{
    \draw[#1,line width=2.6pt,line cap=round]
      plot [smooth] coordinates {(0,0.09) (0.34,0.01) (0.68,0.21) (1.02,0.01) (1.36,0.10)};
  }%
}
% Mot-clé surligné (marqueur orange sous le texte)
\newcommand{\cle}[1]{{\popxbold\color{popdark}#1}}

% ============================================================
% En-tête / pied
% ============================================================
\pagestyle{fancy}
\fancyhf{}
\renewcommand{\headrulewidth}{0pt}
\renewcommand{\footrulewidth}{0pt}
\lhead{\raisebox{-0.15ex}{\poplogo\fontsize{13}{13}\selectfont\color{popink}OnBuch\color{poporange}+}}
\rhead{\poppill{\DOCMATIERE\ \textbullet\ \DOCNIVEAU\ \textbullet\ Leçon \DOCLECON}{white}{popink}}
\lfoot{\popxbold\fontsize{7.6}{7.6}\selectfont\color{popmuted}\MakeUppercase{Cours signé OnBuch+}\ \textbullet\ {\popsemi\DOCTITRE}}
\rfoot{\tikz[baseline=-0.6ex]\node[rounded corners=8.5pt,fill=popink,minimum height=17pt,minimum width=17pt,inner xsep=5pt,inner sep=0]{\popxbold\fontsize{8}{8}\selectfont\color{popcream}\thepage\,/\,\pageref*{LastPage}};}

% ============================================================
% Titres
% ============================================================
\newcommand{\chaptitre}[2]{%
  \begin{tcolorbox}[enhanced,colback=poporange,colframe=popink,boxrule=2.6pt,arc=11pt,
      drop shadow=popink,left=15pt,right=15pt,top=12pt,bottom=11pt,before skip=3pt,after skip=18pt]
    \poppill{#2}{popink}{popcream}\par\vspace{9pt}
    {\raggedright\hyphenpenalty=10000\exhyphenpenalty=10000\popdisplay\fontsize{22}{24}\selectfont\color{popcream}\MakeUppercase{#1}\par}
  \end{tcolorbox}
}
% Section numérotée : pastille orange avec le numéro + titre Archivo
\newcounter{popsec}
\newcommand{\coursec}[1]{%
  \stepcounter{popsec}\setcounter{popsub}{0}%
  \par\vspace{16pt}\noindent
  \begin{minipage}{\linewidth}
  \tikz[baseline=-0.7ex]\node[draw=popink,line width=1.6pt,fill=poporange,rounded corners=4pt,minimum size=22pt,inner sep=0]{\popdisplay\fontsize{12}{12}\selectfont\color{popcream}\thepopsec};%
  \hspace{8pt}{\popdisplay\fontsize{15}{17}\selectfont\color{popink}\MakeUppercase{#1}}\par
  \vspace{4pt}\noindent{\color{poporange}\rule{36pt}{3.6pt}}
  \end{minipage}\par\nopagebreak\vspace{8pt}
}
\newcounter{popsub}
\newcommand{\courssub}[1]{%
  \stepcounter{popsub}%
  \par\vspace{9pt}\noindent{\popxbold\fontsize{11.5}{13}\selectfont\color{popink}{\color{poporange}\thepopsec.\thepopsub}\hspace{6pt}#1}\par\nopagebreak\vspace{4pt}
}

% ============================================================
% Boîtes pédagogiques — même recette : fond blanc, bordure épaisse colorée,
% ombre dure encre, badge plein. Titre optionnel [#1].
% ============================================================
\tcbset{popbase/.style={breakable,enhanced,colback=white,boxrule=2.3pt,arc=9pt,
  shadow={3pt}{-3pt}{0pt}{popink},left=14pt,right=14pt,top=11pt,bottom=11pt,before skip=11pt,after skip=15pt,
  fontupper=\popsemi\fontsize{10.6}{14.5}\selectfont\color{popink},
  fontlower=\popsemi\fontsize{10.6}{14.5}\selectfont\color{popink}}}
% Coche dessinée (le glyphe ✓ n'existe pas dans les polices du document)
\newcommand{\popcheck}[1]{\tikz[baseline=-0.5ex]\draw[#1,line width=1.6pt,line cap=round,line join=round](-0.13,0.0)--(-0.04,-0.09)--(0.13,0.1);}
\newcommand{\popboxhead}[3]{\poppill{#1}{#2}{white}\ifx\relax#3\relax\else\hspace{6pt}{\popxbold\fontsize{10.6}{12}\selectfont\color{popink}#3}\fi\par\vspace{7pt}}

\newtcolorbox{aretenir}[1][]{popbase,colframe=popgreen,before upper={\popboxhead{À retenir}{popgreen}{#1}}}
\newtcolorbox{exemplebox}[1][]{popbase,colframe=poporange,before upper={\popboxhead{Exemple}{poporange}{#1}}}
\newtcolorbox{definition}[1][]{popbase,colframe=popblue,before upper={\popboxhead{Définition}{popblue}{#1}}}
\newtcolorbox{propriete}[1][]{popbase,colframe=poppurple,before upper={\popboxhead{Propriété}{poppurple}{#1}}}
\newtcolorbox{methode}[1][]{popbase,colframe=popgold,before upper={\popboxhead{Méthode}{popgold}{#1}}}
\newtcolorbox{attention}[1][]{popbase,colframe=poppink,before upper={\popboxhead{Attention}{poppink}{#1}}}
\newtcolorbox{experience}[1][]{popbase,colframe=popblue,colback=popblueL,before upper={\popboxhead{Expérience}{popblue}{#1}}}
% « Le savais-tu ? » : carte violette pleine (contexte, culture, Cameroun)
\newtcolorbox{savaistu}[1][]{popbase,colback=poppurple,colframe=popink,
  fontupper=\popsemi\fontsize{10.4}{14.2}\selectfont\color{white},
  before upper={\poppill{Le savais-tu\,?}{white}{poppurple}\ifx\relax#1\relax\else\hspace{6pt}{\popxbold\color{white}#1}\fi\par\vspace{7pt}}}
% Exercice résolu : énoncé en haut, \tcblower puis la solution sur fond crème
\newcounter{popexo}\newcounter{popexr}
\newtcolorbox{exoresolu}[1][]{popbase,colframe=poporange,colbacklower=popcream,
  segmentation style={popink,line width=1.2pt,dashed},
  before upper={\stepcounter{popexr}\popboxhead{Exercice résolu \thepopexr}{poporange}{#1}},
  before lower={\poppill{Solution}{popink}{popcream}\par\vspace{6pt}}}
% Exercice (non résolu en place) — la correction est dans la partie Corrigés
\newtcolorbox{exercice}[1][]{popbase,colframe=popink,
  before upper={\stepcounter{popexo}\popboxhead{Exercice \thepopexo}{popink}{#1}}}
% Corrigé
\newtcolorbox{corrige}[1][]{popbase,colframe=popgreen,colback=popgreenL,
  before upper={\popboxhead{Corrigé}{popgreen}{#1}}}
% Objectifs / prérequis / bilan
\newtcolorbox{objectifs}{popbase,colframe=popink,colback=poporangeL,
  before upper={\popboxhead{Objectifs de la leçon}{popink}{}}}
\newtcolorbox{prerequis}{popbase,colframe=popink,colback=popblueL,
  before upper={\popboxhead{Prérequis}{popblue}{}}}
\newtcolorbox{bilan}{popbase,colframe=popink,colback=white,boxrule=2.8pt,
  before upper={\popboxhead{Fiche bilan}{poporange}{L'essentiel à connaître pour l'examen}}}
% Figure encadrée avec légende
\newtcolorbox{popfigure}[1][]{enhanced,breakable=false,colback=white,colframe=popline,boxrule=1.4pt,arc=8pt,
  left=10pt,right=10pt,top=10pt,bottom=8pt,before skip=10pt,after skip=12pt,halign=center,
  lower separated=false,halign lower=center,
  fontlower=\popsemi\fontsize{9}{11.5}\selectfont\color{popmuted},
  #1}
\newcommand{\legende}[1]{\par\vspace{6pt}{\popsemi\fontsize{9}{11.5}\selectfont\color{popmuted}#1}}

% ============================================================
% Signature OnBuch+
% ============================================================
\newcommand{\poplogoinline}[1]{{\poplogo\fontsize{#1}{#1}\selectfont\color{popink}OnBuch\color{poporange}+}}
\newcommand{\signatureonbuch}{%
  \begin{tcolorbox}[enhanced,colback=popink,colframe=popink,boxrule=2.2pt,arc=10pt,
      drop shadow=poporange,left=14pt,right=14pt,top=10pt,bottom=10pt,before skip=0pt,after skip=0pt]
    \tikz[baseline=-0.7ex]\node[circle,fill=popgreen,draw=popcream,line width=1.3pt,minimum size=22pt,inner sep=0]{\popcheck{white}};%
    \hspace{9pt}\begin{minipage}[c]{0.62\linewidth}
      {\popxbold\fontsize{12}{13}\selectfont\color{popcream}Signé\ \textbullet\ {\poplogo OnBuch\color{poporange}+}}\par\vspace{2pt}
      {\raggedright\popsemi\fontsize{8}{10.5}\selectfont\color{popline}\MakeUppercase{Équipe pédagogique OnBuch+}\\\MakeUppercase{Conforme au programme officiel MINESEC}\par}
    \end{minipage}\hfill
    {\poplogo\fontsize{22}{22}\selectfont\color{popcream}OnBuch\color{poporange}+}
  \end{tcolorbox}%
}

% ============================================================
% Page de garde
% #1 titre · #2 matière · #3 niveau · #4 module · #5 n° leçon · #6 durée ·
% #7 description / objectifs (texte ou itemize)
% ============================================================
\newcommand{\pagegarde}[7]{%
  \thispagestyle{empty}
  \newgeometry{a4paper,margin=1.7cm,top=1.6cm,bottom=1.4cm}%
  \begin{tikzpicture}[remember picture,overlay]
    \fill[poporange!18] ([xshift=-3.2cm,yshift=-3.6cm]current page.north east) circle (4.6cm);
    \fill[poppurple!14] ([xshift=2.4cm,yshift=7.6cm]current page.south west) circle (3.8cm);
  \end{tikzpicture}%
  {\poplogo\fontsize{30}{30}\selectfont\color{popink}OnBuch\color{poporange}+}\hfill
  \raisebox{0.6ex}{\poppill{Cours signé \textbullet\ Premium}{popink}{popcream}}\par
  \vspace{1pt}\popsquiggle{poppurple}\par
  \vspace{8pt}
  \begin{tcolorbox}[enhanced,colback=poporange,colframe=popink,boxrule=2.8pt,arc=13pt,
      drop shadow=popink,left=18pt,right=18pt,top=12pt,bottom=13pt,after skip=10pt]
    \poppill{#2 \textbullet\ #3}{popink}{popcream}\hspace{5pt}\poppill{#4}{white}{popink}\par\vspace{8pt}
    {\popdisplay\fontsize{14}{14}\selectfont\color{popink}LEÇON #5}\par\vspace{4pt}
    {\raggedright\hyphenpenalty=10000\exhyphenpenalty=10000\popdisplay\fontsize{25}{27}\selectfont\color{popcream}\MakeUppercase{#1}\par}
    \vspace{8pt}
    {\popxbold\fontsize{9.5}{9.5}\selectfont\color{popink}\MakeUppercase{Durée conseillée : #6}}
  \end{tcolorbox}
  \begin{tcolorbox}[enhanced,colback=white,colframe=popink,boxrule=2.2pt,arc=10pt,
      drop shadow=popink,left=16pt,right=16pt,top=10pt,bottom=10pt,after skip=8pt]
    \poppill{Dans cette leçon}{poppurple}{white}\par\vspace{5pt}
    \popsemi\fontsize{9.3}{12.6}\selectfont\color{popink}#7
  \end{tcolorbox}
  \noindent\begin{minipage}[t]{0.487\linewidth}
    \begin{tcolorbox}[enhanced,colback=popgreen,colframe=popink,boxrule=2.2pt,arc=10pt,
        drop shadow=popink,left=11pt,right=11pt,top=10pt,bottom=10pt,height=3.1cm,valign=center]
      \tcbox[colback=white,colframe=popink,boxrule=1.3pt,arc=5pt,boxsep=0pt,nobeforeafter,
        left=3pt,right=3pt,top=3pt,bottom=3pt,valign=center]{\qrcode[height=1.9cm]{https://play.google.com/store/apps/details?id=com.luvvix.onbuch&hl=fr}}%
      \hspace{8pt}\begin{minipage}[c]{0.5\linewidth}
        {\popdisplay\fontsize{11}{12.5}\selectfont\color{white}\MakeUppercase{Télécharge OnBuch}}\par\vspace{4pt}
        {\raggedright\popsemi\fontsize{8.4}{11}\selectfont\color{white}Gratuit \textbullet\ Google Play\\Tuteur IA, annales, concours, résultats\par}
      \end{minipage}
    \end{tcolorbox}
  \end{minipage}\hfill
  \begin{minipage}[t]{0.487\linewidth}
    \begin{tcolorbox}[enhanced,colback=poppurple,colframe=popink,boxrule=2.2pt,arc=10pt,
        drop shadow=popink,left=11pt,right=11pt,top=10pt,bottom=10pt,height=3.1cm,valign=center]
      \tikz[baseline=-0.7ex]\node[circle,fill=white,draw=popink,line width=1.3pt,minimum size=1.6cm,inner sep=0]{\popdisplay\fontsize{24}{24}\selectfont\color{poppurple}+};%
      \hspace{8pt}\begin{minipage}[c]{0.6\linewidth}
        {\popdisplay\fontsize{11}{12.5}\selectfont\color{white}\MakeUppercase{Passe à OnBuch+}}\par\vspace{4pt}
        {\raggedright\popsemi\fontsize{8.4}{11}\selectfont\color{white}Tous les cours signés, Tuteur IA illimité, fiches PDF — onglet OnBuch+ de l'app\par}
      \end{minipage}
    \end{tcolorbox}
  \end{minipage}
  \vspace{8pt}
  \popcontenu
  \vfill
  \signatureonbuch
  \restoregeometry
  \newpage
}

% Rangée « Ce cours contient » (page de garde)
\newcommand{\poptile}[3]{% #1 couleur #2 gros texte #3 légende
  \begin{tcolorbox}[enhanced,colback=#1,colframe=popink,boxrule=2pt,arc=9pt,shadow={2.5pt}{-2.5pt}{0pt}{popink},
      left=6pt,right=6pt,top=9pt,bottom=9pt,halign=center,valign=center,height=1.75cm,width=\linewidth,nobeforeafter]
    {\popdisplay\fontsize{12.5}{13}\selectfont\color{white}\MakeUppercase{#2}}\par\vspace{3pt}
    {\centering\popsemi\fontsize{7.8}{9.5}\selectfont\color{white}#3\par}
  \end{tcolorbox}}
\newcommand{\popcontenu}{%
  \par\noindent{\popxboldls\fontsize{7.5}{7.5}\selectfont\color{popmuted}CE COURS SIGNÉ CONTIENT}\par\vspace{7pt}
  \noindent\begin{minipage}[t]{0.235\linewidth}\poptile{popblue}{Cours}{complet, illustré, pas à pas}\end{minipage}\hfill
  \begin{minipage}[t]{0.235\linewidth}\poptile{poporange}{Méthodes}{et exercices résolus}\end{minipage}\hfill
  \begin{minipage}[t]{0.235\linewidth}\poptile{poppink}{Exercices}{type examen + corrigés}\end{minipage}\hfill
  \begin{minipage}[t]{0.235\linewidth}\poptile{popgreen}{Fiche bilan}{l'essentiel à retenir}\end{minipage}\par}

% Sommaire
\newtcolorbox{sommaire}{enhanced,colback=white,colframe=popink,boxrule=2.2pt,arc=10pt,
  drop shadow=popink,left=18pt,right=18pt,top=13pt,bottom=13pt,before skip=6pt,after skip=20pt,
  before upper={\poppill{Sommaire}{poppurple}{white}\par\vspace{9pt}\popsemi\fontsize{10.6}{14.5}\selectfont\color{popink}}}

% Signature de fin de cours — poussée tout en bas de la dernière page
\newcommand{\findecours}{%
  \par\vfill\nopagebreak
  \begin{center}\popsquiggle{poporange}\end{center}\vspace{4pt}
  \signatureonbuch
}
\AtBeginDocument{\def\llbracket{\mathopen{[\mkern-3.5mu[}}\def\rrbracket{\mathclose{]\mkern-3.5mu]}}} % intervalles d'entiers (glyphe ⟦ absent de la police)
% Glyphes absents de Fira Math : redéfinis à partir de glyphes disponibles
\AtBeginDocument{%
  \def\setminus{\mathbin{\backslash}}\let\smallsetminus\setminus
  \def\vdots{\mathord{\vbox{\baselineskip3.2pt\lineskiplimit0pt\kern4pt\hbox{.}\hbox{.}\hbox{.}}}}%
  \def\ddots{\mathinner{\mkern1mu\raise6.5pt\hbox{.}\mkern2mu\raise3.6pt\hbox{.}\mkern2mu\raise0.7pt\hbox{.}\mkern1mu}}%
  \def\triangle{\mathord{\scalebox{0.9}{$\symup{\Delta}$}}}%
  \def\nabla{\mathord{\raisebox{\depth}{\scalebox{1}[-1]{$\symup{\Delta}$}}}}%
  \def\checkmark{\popcheck{popdark}}%
  \def\star{\mathbin{\ast}}\def\bigstar{\mathord{\ast}}%
  \def\diamond{\mathbin{\scalebox{0.6}{$\Diamond$}}}%
  \def\bigcup{\mathop{\mathchoice{\raisebox{-0.3em}{\scalebox{1.7}{$\cup$}}}{\scalebox{1.25}{$\cup$}}{\cup}{\cup}}\displaylimits}%
  \def\bigcap{\mathop{\mathchoice{\raisebox{-0.3em}{\scalebox{1.7}{$\cap$}}}{\scalebox{1.25}{$\cap$}}{\cap}{\cap}}\displaylimits}%
  \def\lesseqqgtr{\mathrel{\lessgtr}}\def\bigoplus{\mathop{\oplus}\displaylimits}%
}
\providecommand{\ssi}{si et seulement si} % abréviation fréquente
\AtBeginDocument{\providecommand{\wideparen}{\overparen}} % arc de cercle

%% FIGFIX-BEGIN v4 — correctifs globaux des figures (docs/onbuch-generation/tools/figfix)
\usepackage{adjustbox}\usepackage{etoolbox}
% toute figure TikZ/pgfplots plus large que la ligne est réduite à la largeur disponible
\BeforeBeginEnvironment{tikzpicture}{\begin{adjustbox}{max width=\linewidth}}
\AfterEndEnvironment{tikzpicture}{\end{adjustbox}}
% légendes pgfplots lisibles par-dessus les courbes
\pgfplotsset{every axis/.append style={legend style={fill=white,fill opacity=0.92,text opacity=1,draw=black!15,inner sep=3pt}}}
% étiquettes d'arêtes (syntaxe edge["..."]) avec fond blanc
\tikzset{every edge quotes/.append style={fill=white,inner sep=1.2pt}}
%% FIGFIX-END
\begin{document}

\pagegarde{Fonctions numériques d’une variable réelle}{Mathématiques}{Tle Tech}{Mathématiques – classe de Terminale technique}{N08}{4 séances (fiche de progression harmonisée)}{Cette leçon approfondit l'étude des fonctions numériques (continuité, monotonie, limites trigonométriques et irrationnelles) et leur représentation graphique. Elle outille l'élève pour modéliser et résoudre des problèmes concrets par une analyse rigoureuse et justifiée.
\vspace{4pt}
\begin{multicols}{2}\small
\begin{itemize}
\item À la fin de cette leçon, je sais définir la continuité d'une fonction sur un intervalle et l'appliquer aux fonctions usuelles.
\item Je sais énoncer et utiliser le théorème des valeurs intermédiaires (TVI) et le théorème de la bijection pour résoudre des équations f(x)=k.
\item Je sais déterminer les limites des fonctions trigonométriques (sin, cos, tan) en 0 et en ±∞, et les limites des fonctions irrationnelles par rationalisation.
\item Je sais étudier la monotonie d'une fonction continue sur un intervalle à l'aide de la dérivée ou par composition.
\item Je sais construire un tableau de variation complet et tracer la courbe représentative d'une fonction faisant intervenir des expressions trigonométriques ou irrationnelles.
\item Je sais modéliser une situation concrète (physique, économique, technique) par une fonction, en interpréter les résultats (limites, variations, solutions d'équations) et rédiger une conclusion argumentée.
\end{itemize}\end{multicols}}

\begin{sommaire}
\begin{multicols}{2}
\begin{enumerate}
\item 1. Continuité sur un intervalle : définitions et propriétés fondamentales
\item 2. Fonction continue et strictement monotone : Bijection et fonction réciproque
\item 3. Limites des fonctions trigonométriques : limites remarquables et encadrement
\item 4. Limites des fonctions irrationnelles : rationalisation et composition
\item 5. Représentation graphique : méthodologie complète et lecture graphique
\item 6. Modélisation et résolution de problèmes concrets (Synthèse)
\item Activité d'intégration
\item Méthodes et exercices résolus
\item Exercices
\item Corrigés des exercices
\item Fiche bilan
\end{enumerate}
\end{multicols}
\end{sommaire}

\begin{objectifs}
\begin{itemize}
\item À la fin de cette leçon, je sais définir la continuité d'une fonction sur un intervalle et l'appliquer aux fonctions usuelles.
\item Je sais énoncer et utiliser le théorème des valeurs intermédiaires (TVI) et le théorème de la bijection pour résoudre des équations f(x)=k.
\item Je sais déterminer les limites des fonctions trigonométriques (sin, cos, tan) en 0 et en ±∞, et les limites des fonctions irrationnelles par rationalisation.
\item Je sais étudier la monotonie d'une fonction continue sur un intervalle à l'aide de la dérivée ou par composition.
\item Je sais construire un tableau de variation complet et tracer la courbe représentative d'une fonction faisant intervenir des expressions trigonométriques ou irrationnelles.
\item Je sais modéliser une situation concrète (physique, économique, technique) par une fonction, en interpréter les résultats (limites, variations, solutions d'équations) et rédiger une conclusion argumentée.
\item Je sais identifier et corriger les erreurs fréquentes sur la continuité des fonctions composées et sur l'application du TVI (intervalle non fermé, fonction non continue).
\end{itemize}
\end{objectifs}

\begin{prerequis}
\begin{itemize}
\item Définition de la limite d'une fonction en un point et en ±∞ (classe de 1ère).
\item Opérations sur les limites et formes indéterminées (1ère).
\item Définition de la dérivée, lien avec la monotonie et tableau de variations (1ère).
\item Fonctions usuelles : polynômes, rationnelles, exponentielle, logarithme népérien (1ère).
\item Notions de base sur les fonctions trigonométriques (cercle trigonométrique, formules d'addition) et irrationnelles (domaine de définition).
\end{itemize}
\end{prerequis}

\newpage

\coursec{Continuité sur un intervalle : définitions et propriétés fondamentales}

\textbf{Situation d'accroche.} Monsieur Kamga, gérant d'une petite usine de transformation de cacao à Dschang, observe que la température de fermentation des fèves, cruciale pour la qualité de l'arôme, varie au cours de la journée selon une loi modélisée par une fonction trigonométrique. Il doit garantir que cette température reste dans un intervalle précis $[45\,^\circ\mathrm{C}\,;\,55\,^\circ\mathrm{C}]$ pendant au moins 6 heures continues. Comment l'étude de la continuité, de la monotonie et des limites de cette fonction peut-elle l'aider à valider son processus sans surveillance constante ?

\textbf{Problématique.} Quand peut-on affirmer qu'une grandeur évolue « sans saut » ? Et si une fonction évolue sans saut entre deux valeurs, passe-t-elle forcément par toutes les valeurs intermédiaires ? C'est toute la question de la \cle{continuité}, que nous allons étudier dans cette section.

\courssub{De la continuité en un point à la continuité sur un intervalle}

\textbf{Intuition.} Imagine que tu traces la courbe d'une fonction au crayon, sans jamais lever la main. Si tu peux parcourir toute la courbe d'un bout à l'autre d'un intervalle sans lever le crayon, la fonction est \emph{continue} sur cet intervalle. Si à un endroit tu es obligé de lever le crayon (saut, trou, courbe qui part à l'infini), la fonction y est \emph{discontinue}.

\begin{definition}[Continuité en un point (rappel)]
Soit $f$ une fonction définie sur un intervalle contenant le réel $a$. On dit que $f$ est \cle{continue en $a$} lorsque :
\[ \lim_{x \to a} f(x) = f(a). \]
Autrement dit : la limite de $f$ en $a$ existe, est finie, et est égale à la valeur $f(a)$.
\end{definition}

\begin{definition}[Continuité sur un intervalle]
Soit $I$ un intervalle et $f$ une fonction définie sur $I$.
\begin{itemize}
\item $f$ est \cle{continue sur $I$} si elle est continue en tout point de $I$.
\item Pour une borne \textbf{fermée} de l'intervalle, on exige seulement la continuité « d'un côté » : $f$ est continue sur $[a\,;\,b]$ si elle est continue en tout point de $]a\,;\,b[$, \emph{continue à droite en $a$} ($\lim_{x \to a^+} f(x) = f(a)$) et \emph{continue à gauche en $b$} ($\lim_{x \to b^-} f(x) = f(b)$).
\item Pour un intervalle semi-ouvert comme $[a\,;\,b[$, on demande la continuité à droite en $a$ et la continuité en tout point de $]a\,;\,b[$.
\end{itemize}
\end{definition}

\begin{exemplebox}[La fonction racine carrée]
La fonction $f : x \mapsto \sqrt{x}$ est définie sur $[0\,;\,+\infty[$. Elle est continue en tout point de $]0\,;\,+\infty[$, et continue \emph{à droite} en $0$ car $\lim_{x \to 0^+} \sqrt{x} = 0 = f(0)$. On dit donc que $f$ est continue sur $[0\,;\,+\infty[$, même si elle n'est pas définie à gauche de $0$.
\end{exemplebox}

\courssub{Continuité des fonctions usuelles}

\begin{propriete}[Continuité des fonctions usuelles (admise)]
Les fonctions suivantes sont continues sur \textbf{leur ensemble de définition} :
\begin{itemize}
\item les fonctions \cle{polynômes} (continues sur $\mathbb{R}$ tout entier) ;
\item les fonctions \cle{rationnelles} (quotients de polynômes), continues partout sauf aux valeurs qui annulent le dénominateur ;
\item les fonctions \cle{trigonométriques} : $\sin$ et $\cos$ continues sur $\mathbb{R}$ ; $\tan$ continue sur chaque intervalle de son domaine $\mathbb{R} \setminus \left\{ \frac{\pi}{2} + k\pi,\ k \in \mathbb{Z} \right\}$ ;
\item la fonction \cle{exponentielle} $\exp$, continue sur $\mathbb{R}$ ;
\item la fonction \cle{logarithme népérien} $\ln$, continue sur $]0\,;\,+\infty[$ ;
\item les fonctions \cle{irrationnelles} du type $x \mapsto \sqrt{u(x)}$, continues là où $u(x) > 0$ (et à droite/gauche aux bornes où $u(x) = 0$).
\end{itemize}
\end{propriete}

\begin{attention}[« Continue sur son ensemble de définition » ne veut pas dire « continue partout »]
La fonction $f : x \mapsto \dfrac{1}{x}$ est continue sur son ensemble de définition $\mathbb{R}^* = ]-\infty\,;\,0[ \cup ]0\,;\,+\infty[$. Mais elle n'est \textbf{pas} continue en $0$ : elle n'y est même pas définie, et $\lim_{x \to 0^+} \dfrac{1}{x} = +\infty$. Dire « $f$ est continue » sans préciser l'intervalle est une erreur de rédaction fréquente au Baccalauréat technique.
\end{attention}

\courssub{Opérations sur les fonctions continues}

\begin{propriete}[Opérations algébriques et composition]
Soient $f$ et $g$ deux fonctions continues sur un intervalle $I$, et $\lambda$ un réel.
\begin{itemize}
\item $f + g$, $\lambda f$, $f \times g$ sont continues sur $I$ ;
\item $\dfrac{f}{g}$ est continue sur $I$ \textbf{à condition que $g$ ne s'annule pas sur $I$} ;
\item $|f|$ est continue sur $I$ ;
\item si $f$ est continue sur $I$ et $g$ continue sur un intervalle contenant $f(I)$, alors la composée $g \circ f$ est continue sur $I$.
\end{itemize}
\end{propriete}

\begin{attention}[La partie entière : le contre-exemple à connaître]
La fonction \cle{partie entière} $E(x)$ (plus grand entier inférieur ou égal à $x$) est continue sur chaque intervalle $[n\,;\,n+1[$ à droite en $n$ et sur $]n\,;\,n+1[$ ($n \in \mathbb{Z}$), mais elle est \textbf{discontinue en tout point entier} : par exemple $\lim_{x \to 2^-} E(x) = 1$ alors que $E(2) = 2$. La courbe présente des « marches d'escalier » avec des sauts. Une opération (somme, produit, composition) impliquant la partie entière peut donc créer des \cle{points de rupture}.
\end{attention}

\begin{popfigure}
\begin{center}
\begin{tikzpicture}[scale=1.1]
\draw[->,popmuted] (-0.5,0) -- (4.5,0) node[below right]{$x$};
\draw[->,popmuted] (0,-0.5) -- (0,3.5) node[above left]{$y$};
\foreach \n in {0,1,2,3}{
  \draw[very thick,popblue] (\n,\n) -- (\n+1,\n);
  \fill[popblue] (\n,\n) circle (1.6pt);
  \draw[fill=white,draw=popblue,thick] (\n+1,\n) circle (1.6pt);
}
\foreach \x in {1,2,3}{\draw[popline,dashed] (\x,0) node[below]{$\x$} -- (\x,\x);}
\foreach \y in {1,2,3}{\draw[popline,dashed] (0,\y) node[left]{$\y$} -- (\y,\y);}
\end{tikzpicture}
\end{center}
\legende{Fonction partie entière : points pleins = valeurs prises, points creux = valeurs non prises. Sauts en chaque entier.}
\end{popfigure}

\courssub{Le théorème des valeurs intermédiaires (TVI)}

\textbf{Intuition.} Reprenons l'usine de Monsieur Kamga : si la température passe de $42\,^\circ\mathrm{C}$ à $58\,^\circ\mathrm{C}$ de façon continue (sans saut brutal), alors elle est forcément passée par $45\,^\circ\mathrm{C}$, par $50\,^\circ\mathrm{C}$, et par toute valeur intermédiaire. On ne peut pas chauffer une cuve de $42$ à $58$ degrés sans traverser toutes les températures entre les deux. C'est exactement ce que dit le TVI.

\begin{propriete}[Théorème des valeurs intermédiaires (admis ; intuition géométrique exigible)]
Soit $f$ une fonction \textbf{continue sur un segment} $[a\,;\,b]$. Pour tout réel $k$ compris entre $f(a)$ et $f(b)$, il existe \textbf{au moins un} réel $c \in [a\,;\,b]$ tel que :
\[ f(c) = k. \]
Autrement dit : l'image par $f$ du segment $[a\,;\,b]$ contient tout l'intervalle d'extrémités $f(a)$ et $f(b)$.
\end{propriete}

\begin{experience}[Démonstration guidée du TVI par dichotomie (intuition)]
On suppose $f(a) < k < f(b)$ et on construit une suite d'intervalles emboîtés qui « traquent » une solution.
\begin{enumerate}
\item On coupe $[a\,;\,b]$ en deux : soit $m = \dfrac{a+b}{2}$ le milieu.
\item Si $f(m) = k$, c'est gagné, on pose $c = m$. Sinon :
   \begin{itemize}
   \item si $f(m) < k$, la solution se cache dans $[m\,;\,b]$ (car $f$ passe de $f(m) < k$ à $f(b) > k$) ;
   \item si $f(m) > k$, la solution se cache dans $[a\,;\,m]$.
   \end{itemize}
\item On recommence avec le nouvel intervalle, deux fois plus petit. À chaque étape, $k$ reste compris entre les images des bornes.
\item Après $n$ étapes, on obtient un intervalle de longueur $\dfrac{b-a}{2^n}$ contenant une solution. En faisant tendre $n$ vers $+\infty$, cet intervalle se réduit à un point $c$, et la continuité de $f$ impose $f(c) = k$.
\end{enumerate}
C'est le principe de la \cle{dichotomie}, utilisé par les calculatrices et les ordinateurs pour résoudre les équations numériquement.
\end{experience}

\begin{popfigure}
\begin{center}
\begin{tikzpicture}
\begin{axis}[
  width=0.85\linewidth, height=6cm,
  axis lines=middle,
  xlabel={$x$}, ylabel={$y$},
  xmin=-0.4, xmax=4.4, ymin=-0.6, ymax=4.4,
  xtick={0.5,3.5}, xticklabels={$a$,$b$},
  ytick={1,2.2,3.09}, yticklabels={$f(a)$,$k$,$f(b)$},
  clip=false
]
\addplot[popcurve,domain=0.5:3.5,samples=200]{1 + 0.35*(x-0.5) + 0.18*(x-0.5)^2 - 0.02*(x-0.5)^3 + 0.15*sin(deg(2*(x-0.5)))};
\addplot[poporange,thick,dashed,domain=-0.3:4.2]{2.2};
\addplot[popline,dashed] coordinates {(2.58,0) (2.58,2.2)};
\node[below,popdark] at (axis cs:2.58,0) {$c$};
\addplot[only marks,mark=*,poporange,mark size=2pt] coordinates {(2.58,2.2)};
\end{axis}
\end{tikzpicture}
\end{center}
\legende{Illustration du TVI : $f$ continue sur $[a\,;\,b]$, la droite $y = k$ (avec $k$ entre $f(a)$ et $f(b)$) coupe la courbe au moins une fois (ici en $c$).}
\end{popfigure}

\begin{popfigure}
\begin{center}
\begin{tikzpicture}
\begin{axis}[
  width=0.85\linewidth, height=6cm,
  axis lines=middle,
  xlabel={$x$}, ylabel={$y$},
  xmin=-3.2, xmax=3.2, ymin=-3.2, ymax=3.2,
  xtick={-1,1}, ytick={-1,1},
  clip=false
]
\addplot[popcurve,domain=-3:0.34,samples=100]{1/x};
\addplot[popcurve,domain=0.34:3,samples=100]{1/x};
\addplot[poporange,thick,dashed,domain=-3:3]{0};
\node[poporange] at (axis cs:2.2,0.35) {$y=0$};
\end{axis}
\end{tikzpicture}
\end{center}
\legende{Contre-exemple : $f(x) = \dfrac{1}{x}$ sur $[-1\,;\,1]$. On a $f(-1) = -1 < 0 < 1 = f(1)$, pourtant l'équation $\frac{1}{x} = 0$ n'a aucune solution : $f$ n'est pas continue en $0$, le TVI ne s'applique pas.}
\end{popfigure}

\begin{attention}[Erreurs fréquentes avec le TVI]
\begin{itemize}
\item Appliquer le TVI sur un intervalle \textbf{ouvert} $]a\,;\,b[$ : il faut un \textbf{segment} $[a\,;\,b]$ (bornes fermées incluses).
\item Oublier de vérifier la \textbf{continuité} : pour $f(x) = \dfrac{1}{x}$ sur $[-1\,;\,1]$, on a bien $f(-1) = -1$ et $f(1) = 1$, donc $0$ est « entre » $f(-1)$ et $f(1)$, mais $\dfrac{1}{x} = 0$ n'a \textbf{aucune} solution. La continuité fait défaut en $0$.
\item Conclure à l'\textbf{unicité} de la solution : le TVI donne seulement l'\textbf{existence} d'au moins une solution. Pour l'unicité, il faut ajouter la \textbf{stricte monotonie} (voir section 2).
\end{itemize}
\end{attention}

\begin{methode}[{Montrer qu'une équation $f(x) = k$ admet une solution sur $[a\,;\,b]$}]
\begin{enumerate}
\item \textbf{Vérifier la continuité} : justifier que $f$ est continue sur $[a\,;\,b]$ (fonction usuelle, somme, composée de fonctions continues...).
\item \textbf{Calculer} $f(a)$ et $f(b)$.
\item \textbf{Comparer} : vérifier que $k$ est compris entre $f(a)$ et $f(b)$.
\item \textbf{Conclure} : « D'après le théorème des valeurs intermédiaires, l'équation $f(x) = k$ admet au moins une solution $c$ dans $[a\,;\,b]$. »
\end{enumerate}
\end{methode}

\begin{exoresolu}[{L'équation $\cos x = x$ admet une solution unique dans $\left[0\,;\,\dfrac{\pi}{2}\right]$}]
\textbf{Énoncé.} Montrer que l'équation $\cos x = x$ admet une unique solution dans l'intervalle $\left[0\,;\,\dfrac{\pi}{2}\right]$.
\tcblower
\textbf{Solution détaillée.} On pose $f(x) = \cos x - x$. Résoudre $\cos x = x$ revient à résoudre $f(x) = 0$.

\textbf{Étape 1 — Continuité.} La fonction $\cos$ est continue sur $\mathbb{R}$ et $x \mapsto x$ est un polynôme, donc $f$, différence de deux fonctions continues, est continue sur $\mathbb{R}$, en particulier sur $\left[0\,;\,\dfrac{\pi}{2}\right]$.

\textbf{Étape 2 — Images des bornes.}
\[ f(0) = \cos 0 - 0 = 1 > 0 \qquad ; \qquad f\left(\tfrac{\pi}{2}\right) = \cos\tfrac{\pi}{2} - \tfrac{\pi}{2} = -\tfrac{\pi}{2} < 0. \]

\textbf{Étape 3 — Comparaison.} $k = 0$ est compris entre $-\dfrac{\pi}{2}$ et $1$.

\textbf{Étape 4 — Existence (TVI).} D'après le théorème des valeurs intermédiaires, il existe au moins un réel $c \in \left[0\,;\,\dfrac{\pi}{2}\right]$ tel que $f(c) = 0$, c'est-à-dire $\cos c = c$.

\textbf{Étape 5 — Unicité (monotonie).} $f'(x) = -\sin x - 1$. Pour $x \in \left]0\,;\,\dfrac{\pi}{2}\right[$, $\sin x > 0$, donc $f'(x) = -(\sin x + 1) < 0$ : $f$ est \textbf{strictement décroissante} sur $\left[0\,;\,\dfrac{\pi}{2}\right]$. Elle ne peut donc pas s'annuler deux fois.

\textbf{Conclusion.} L'équation $\cos x = x$ admet une \textbf{unique} solution dans $\left[0\,;\,\dfrac{\pi}{2}\right]$. (Une valeur approchée : $c \approx 0{,}739$.)
\end{exoresolu}

\courssub{Corollaire et théorème des bornes}

\begin{propriete}[Corollaire du TVI : équation $f(x) = 0$]
Si $f$ est continue sur $[a\,;\,b]$ et si $f(a)$ et $f(b)$ sont de \textbf{signes contraires} (c'est-à-dire $f(a) \times f(b) < 0$), alors l'équation $f(x) = 0$ admet au moins une solution dans $]a\,;\,b[$.
\end{propriete}

C'est le cas particulier $k = 0$ du TVI, et c'est la forme la plus utilisée en pratique : un simple \textbf{changement de signe} entre les bornes garantit une racine. C'est ce critère que l'on utilise pour localiser les solutions d'équations que l'on ne sait pas résoudre algébriquement.

\begin{propriete}[Théorème des bornes (admis)]
Si $f$ est continue sur un \textbf{segment} $[a\,;\,b]$, alors $f$ est \cle{bornée} sur $[a\,;\,b]$ et \textbf{atteint ses bornes} : il existe deux réels $x_1$ et $x_2$ dans $[a\,;\,b]$ tels que, pour tout $x \in [a\,;\,b]$ :
\[ f(x_1) = m \leq f(x) \leq M = f(x_2), \]
où $m$ est le minimum et $M$ le maximum de $f$ sur $[a\,;\,b]$.
\end{propriete}

\textbf{Interprétation pour Monsieur Kamga.} Si la température $T(t)$ est une fonction continue du temps sur la plage horaire $[t_1\,;\,t_2]$, alors elle admet un minimum $m$ et un maximum $M$ \emph{effectivement atteints}. Pour garantir que la température reste dans $[45\,;\,55]$, il suffit donc de vérifier que $m \geq 45$ et $M \leq 55$ : pas besoin de surveiller chaque instant !

\begin{attention}[Le mot « segment » est indispensable]
Sur un intervalle non fermé ou non borné, le théorème des bornes peut être faux : $f(x) = \dfrac{1}{x}$ est continue sur $]0\,;\,1]$ mais n'y est pas majorée (elle tend vers $+\infty$ en $0^+$). De même $g(x) = x$ sur $]0\,;\,1[$ est bornée mais n'atteint ni $0$ ni $1$. Retenir : \textbf{continuité + segment} $\Rightarrow$ bornes atteintes.
\end{attention}

\courssub{Arbre de décision : vérifier la continuité d'une fonction}

\begin{popfigure}
\begin{center}
\begin{tikzpicture}[
  font=\small,
  decision/.style={diamond, draw=popblue, fill=popblueL, text width=3.4cm, align=center, inner sep=1pt, aspect=2},
  action/.style={rectangle, rounded corners, draw=popgreen, fill=popgreenL, text width=3.6cm, align=center, inner sep=4pt},
  stop/.style={rectangle, rounded corners, draw=poporange, fill=poporangeL, text width=3.6cm, align=center, inner sep=4pt},
  arr/.style={-{Stealth}, thick, popdark}
]
\node[action] (start) {Fonction $f$ à étudier sur $I$};
\node[decision, below=0.9cm of start] (dom) {$f$ est-elle définie en tout point de $I$ ?};
\node[decision, below=1.1cm of dom] (comp) {$f$ est-elle usuelle ou somme, produit, quotient, composée de fonctions continues ?};
\node[decision, below=1.1cm of comp] (denom) {Le dénominateur (ou l'intérieur des racines) s'annule-t-il sur $I$ ?};
\node[action, below left=1.1cm and -0.5cm of denom] (ok) {$f$ est continue sur $I$ : TVI et théorème des bornes applicables (si $I$ segment)};
\node[stop, below right=1.1cm and -0.5cm of denom] (ko) {Points de rupture : étudier les limites en ces points};
\draw[arr] (start) -- (dom);
\draw[arr] (dom) -- node[right]{oui} (comp);
\draw[arr] (dom.east) -- ++(2.2,0) node[above,pos=0.5]{non} |- (ko.east);
\draw[arr] (comp) -- node[right]{oui} (denom);
\draw[arr] (comp.east) -- ++(2.2,0) node[above,pos=0.5]{non} |- (ko.east);
\draw[arr] (denom) -- node[above left]{non} (ok);
\draw[arr] (denom) -- node[above right]{oui} (ko);
\end{tikzpicture}
\end{center}
\legende{Arbre de décision : vérifier la continuité d'une fonction sur un intervalle $I$.}
\end{popfigure}

\begin{exercice}[À faire]
Soit $g$ la fonction définie par $g(x) = x^3 + 2x - 3$.
\begin{enumerate}
\item Justifier que $g$ est continue sur $\mathbb{R}$.
\item Montrer que l'équation $g(x) = 0$ admet au moins une solution dans $[0\,;\,2]$.
\item Calculer $g(1)$. En déduire la solution de l'équation dans $[0\,;\,2]$.
\end{enumerate}
\end{exercice}

\begin{corrige}[Exercice]
\begin{enumerate}
\item $g$ est une fonction polynôme, donc continue sur $\mathbb{R}$.
\item $g(0) = -3 < 0$ et $g(2) = 8 + 4 - 3 = 9 > 0$. Comme $g$ est continue sur le segment $[0\,;\,2]$ et que $0$ est compris entre $g(0)$ et $g(2)$, le TVI garantit l'existence d'au moins une solution de $g(x) = 0$ dans $[0\,;\,2]$.
\item $g(1) = 1 + 2 - 3 = 0$. Donc $x = 1$ est une solution de l'équation dans $[0\,;\,2]$. (On peut même montrer qu'elle est unique : $g'(x) = 3x^2 + 2 > 0$ pour tout $x$, donc $g$ est strictement croissante sur $\mathbb{R}$.)
\end{enumerate}
\end{corrige}

\begin{savaistu}[Le TVI, un théorème aux applications concrètes]
Le théorème des valeurs intermédiaires est dû au mathématicien tchèque Bernard Bolzano (1817). Au-delà des mathématiques, il justifie des faits de la vie courante : si tu montes de Douala (altitude $13$ m) au sommet du mont Cameroun ($4\,040$ m) par un sentier continu, tu passes forcément par l'altitude $1\,000$ m. En informatique, la dichotomie issue de sa démonstration est à la base des algorithmes de recherche de racines utilisés dans les logiciels de CAO et les simulateurs industriels... comme celui qui pourrait aider Monsieur Kamga à piloter sa cuve de fermentation.
\end{savaistu}

\begin{aretenir}[L'essentiel de la section]
\begin{itemize}
\item $f$ continue en $a$ signifie $\lim_{x \to a} f(x) = f(a)$ ; $f$ continue sur $I$ signifie continue en tout point de $I$ (continuité à droite/gauche aux bornes fermées).
\item Les fonctions usuelles (polynômes, $\sin$, $\cos$, $\exp$, $\ln$, $\sqrt{\ }$) sont continues sur leur ensemble de définition ; sommes, produits, quotients (dénominateur non nul) et composées de fonctions continues sont continues.
\item \textbf{TVI} : $f$ continue sur le \textbf{segment} $[a\,;\,b]$ et $k$ entre $f(a)$ et $f(b)$ $\Rightarrow$ il existe $c \in [a\,;\,b]$ avec $f(c) = k$ (existence seulement ; unicité via stricte monotonie).
\item \textbf{Théorème des bornes} : continue sur un segment $\Rightarrow$ bornée et bornes atteintes.
\item Pièges : intervalle ouvert, fonction discontinue (partie entière, $\frac{1}{x}$), conclure à l'unicité sans monotonie.
\end{itemize}
\end{aretenir}

\coursec{Fonction continue et strictement monotone : Bijection et fonction réciproque}

Dans la section précédente, nous avons découvert la notion de \cle{continuité} sur un intervalle : une fonction continue ne fait pas de « saut », sa courbe se trace sans lever le crayon, et le théorème des valeurs intermédiaires garantit l'existence de solutions à l'équation $f(x)=k$. Mais ce théorème ne dit rien sur l'\emph{unicité} de la solution. Dans un atelier, lorsqu'on règle une machine pour obtenir une pièce de dimension exacte, on veut être certain qu'il n'existe qu'\textbf{un seul} réglage possible. C'est exactement ce que garantit la \cle{monotonie stricte} combinée à la continuité : c'est l'objet du \cle{théorème de la bijection}, pilier de cette section.

\courssub{Rappel : fonctions strictement monotones}

\begin{definition}[Monotonie stricte]
Soit $f$ une fonction définie sur un intervalle $I$.
\begin{itemize}
\item $f$ est \cle{strictement croissante} sur $I$ si pour tous réels $a$ et $b$ de $I$ : $a < b \Rightarrow f(a) < f(b)$.
\item $f$ est \cle{strictement décroissante} sur $I$ si pour tous réels $a$ et $b$ de $I$ : $a < b \Rightarrow f(a) > f(b)$.
\item $f$ est \cle{strictement monotone} sur $I$ si elle est strictement croissante ou strictement décroissante sur $I$.
\end{itemize}
\end{definition}

L'intuition est simple : une fonction strictement monotone ne « revient jamais en arrière ». Par conséquent, elle ne peut pas prendre deux fois la même valeur : si $a \neq b$, alors $f(a) \neq f(b)$. C'est cette propriété d'\cle{injectivité} qui va assurer l'unicité des solutions des équations.

\begin{propriete}[Lien entre dérivée et monotonie]
Soit $f$ une fonction dérivable sur un intervalle $I$.
\begin{itemize}
\item Si $f'(x) > 0$ pour tout $x \in I$ (sauf éventuellement en des points isolés où $f'$ s'annule), alors $f$ est strictement croissante sur $I$.
\item Si $f'(x) < 0$ pour tout $x \in I$ (sauf éventuellement en des points isolés), alors $f$ est strictement décroissante sur $I$.
\end{itemize}
\end{propriete}

\begin{attention}[Réciproque fausse en général]
Si $f$ est strictement croissante et dérivable, on n'a pas forcément $f'(x) > 0$ partout. Exemple classique : $f(x) = x^3$ est strictement croissante sur $\mathbb{R}$, pourtant $f'(0) = 0$. La dérivée peut s'annuler en un point isolé sans détruire la stricte monotonie. En revanche, si $f'$ s'annule sur tout un intervalle, la fonction est constante sur cet intervalle et n'est plus strictement monotone.
\end{attention}

\begin{experience}[Découverte graphique : le test de la ligne horizontale]
Dans un repère, tracez la courbe d'une fonction continue et strictement croissante (par exemple $y=x^3$). Placez une règle horizontale (parallèle à l'axe des abscisses) et faites-la glisser verticalement.
\begin{enumerate}
\item Combien de points d'intersection la règle a-t-elle avec la courbe lorsque la règle se trouve entre les valeurs extrêmes de la fonction ?
\item Que se passe-t-il si la fonction n'est pas strictement monotone (ex: $y=x^2$ sur $\mathbb{R}$) ?
\item Conclure sur le lien entre \emph{continuité}, \emph{monotonie stricte} et \emph{unicité} de l'intersection avec une horizontale $y=k$.
\end{enumerate}
Cette activité visuelle illustre le cœur du théorème de la bijection : une horizontale coupe la courbe en \textbf{au plus un point} (injectivité) et \textbf{au moins un point} si $k$ est dans l'image (surjectivité sur l'image), donc \textbf{exactement un point}.
\end{experience}

\courssub{Le théorème de la bijection}

Imaginons la courbe d'une fonction continue et strictement croissante sur un intervalle $I$. Puisqu'elle monte sans jamais redescendre et sans saut, toute horizontale d'équation $y = k$ située entre les valeurs extrêmes coupe la courbe en \textbf{exactement un point}. C'est toute la substance du théorème suivant.

\begin{propriete}[Théorème de la bijection (admis)]
Soit $f$ une fonction \cle{continue} et \cle{strictement monotone} sur un intervalle $I$. Alors :
\begin{enumerate}
\item $f$ réalise une \cle{bijection} de $I$ sur $f(I)$, c'est-à-dire que pour tout réel $k$ de $f(I)$, l'équation $f(x) = k$ admet une \textbf{unique} solution dans $I$ ;
\item l'ensemble image $f(I)$ est un intervalle de même nature que $I$ (ses bornes sont les limites ou les valeurs de $f$ aux bornes de $I$) ;
\item $f$ admet une \cle{fonction réciproque} $f^{-1}$, définie sur $f(I)$ et à valeurs dans $I$ ;
\item $f^{-1}$ est \cle{continue} sur $f(I)$ et possède le \cle{même sens de variation} que $f$.
\end{enumerate}
Ce théorème est admis au Probatoire et au Baccalauréat technique ; sa démonstration n'est pas exigible, mais son application doit être parfaitement maîtrisée.
\end{propriete}

\begin{definition}[Fonction réciproque]
Soit $f$ une bijection de $I$ sur $J = f(I)$. La \cle{fonction réciproque} de $f$, notée $f^{-1}$, est l'unique fonction de $J$ dans $I$ telle que :
\[ f^{-1}\bigl(f(x)\bigr) = x \quad \text{pour tout } x \in I \qquad \text{et} \qquad f\bigl(f^{-1}(y)\bigr) = y \quad \text{pour tout } y \in J. \]
Autrement dit : $f(x) = y \iff x = f^{-1}(y)$.
\end{definition}

\begin{aretenir}[Propriété fondamentale des courbes]
Dans un repère orthonormé, les courbes représentatives de $f$ et de $f^{-1}$ sont \cle{symétriques par rapport à la droite d'équation $y = x$} (la première bissectrice). En effet, le point $(a\,;\,b)$ appartient à la courbe de $f$ si et seulement si le point $(b\,;\,a)$ appartient à celle de $f^{-1}$.
\end{aretenir}

\begin{attention}[Réciproque $\neq$ inverse]
Erreur fréquente et grave : ne jamais confondre
\begin{itemize}
\item la \textbf{réciproque} $f^{-1}(x)$, qui « défait » ce que fait $f$ (exemple : la réciproque de $x \mapsto x^3$ est $x \mapsto \sqrt[3]{x}$) ;
\item l'\textbf{inverse} $\dfrac{1}{f(x)}$, qui est un quotient (exemple : l'inverse de $x^3$ est $\dfrac{1}{x^3}$).
\end{itemize}
Ce sont deux objets totalement différents. Autre piège : une fonction strictement monotone sur $\mathbb{R}$ n'est pas forcément bijective \emph{de $\mathbb{R}$ sur $\mathbb{R}$} : par exemple $f(x)=\arctan x$ est strictement croissante sur $\mathbb{R}$, mais $f(\mathbb{R}) = \left]-\dfrac{\pi}{2}\,;\,\dfrac{\pi}{2}\right[ \neq \mathbb{R}$. Elle réalise une bijection de $\mathbb{R}$ sur son image seulement.
\end{attention}

\begin{popfigure}
\begin{center}
\begin{tikzpicture}[scale=1.05]
\begin{axis}[
  axis lines=middle, xlabel={$x$}, ylabel={$y$},
  xmin=-2.6, xmax=2.6, ymin=-2.6, ymax=2.6,
  xtick={-2,-1,1,2}, ytick={-2,-1,1,2},
  width=0.72\linewidth, grid=none, axis equal image]
  \addplot[popcurve,domain=-1.45:1.45,samples=200]{x^3};
  % Inverse plotted parametrically: (t^3, t) for t in [-1.45, 1.45]
  \addplot[poporange,very thick,domain=-1.45:1.45,samples=200]({x^3},x);
  \addplot[popmuted,dashed,domain=-2.5:2.5,samples=2]{x};
  \node[popblue] at (axis cs:-1.3,-2.5) {$y=x^3$};
  \node[poporange] at (axis cs:2.45,1.15) {$y=\sqrt[3]{x}$};
  \node[popmuted] at (axis cs:-2.3,-2.05) {\small $y=x$};
\end{axis}
\end{tikzpicture}
\end{center}
\legende{Les courbes de $f(x)=x^3$ et de sa réciproque $f^{-1}(x)=\sqrt[3]{x}$ sont symétriques par rapport à la droite $y=x$.}
\end{popfigure}

\courssub{Méthode : montrer qu'une fonction est bijective et déterminer $f(I)$}

\begin{methode}[Étudier la bijectivité d'une fonction sur un intervalle]
Pour montrer qu'une fonction $f$ réalise une bijection d'un intervalle $I$ sur son image :
\begin{enumerate}
\item \textbf{Continuité} : justifier que $f$ est continue sur $I$ (somme, produit, composée de fonctions continues usuelles).
\item \textbf{Monotonie stricte} : calculer $f'(x)$, étudier son signe sur $I$ et conclure que $f$ est strictement monotone.
\item \textbf{Conclusion} : invoquer le théorème de la bijection : $f$ est une bijection de $I$ sur $f(I)$.
\item \textbf{Ensemble image} : déterminer $f(I)$ grâce aux valeurs ou limites de $f$ aux bornes de $I$ (tableau de variations).
\end{enumerate}
Pour résoudre ensuite $f(x)=k$ avec $k \in f(I)$ : conclure à l'existence et à l'\textbf{unicité} de la solution, puis éventuellement en donner un encadrement par balayage.
\end{methode}

\begin{exoresolu}[Une bijection de $\mathbb{R}$ sur $\mathbb{R}$]
\textbf{Énoncé.} Soit $f$ la fonction définie sur $\mathbb{R}$ par $f(x) = 2x + \sin x$.
\begin{enumerate}
\item Montrer que $f$ réalise une bijection de $\mathbb{R}$ sur $\mathbb{R}$.
\item Déterminer l'intervalle image de $f$ sur $[0\,;\,\pi]$.
\end{enumerate}
\tcblower
\textbf{Solution détaillée.}
\begin{enumerate}
\item \textbf{Continuité.} Les fonctions $x \mapsto 2x$ et $x \mapsto \sin x$ sont continues sur $\mathbb{R}$, donc leur somme $f$ est continue sur $\mathbb{R}$.

\textbf{Monotonie.} $f$ est dérivable sur $\mathbb{R}$ et $f'(x) = 2 + \cos x$. Or pour tout réel $x$, $\cos x \geq -1$, donc :
\[ f'(x) = 2 + \cos x \geq 2 - 1 = 1 > 0. \]
La dérivée est strictement positive sur $\mathbb{R}$ : $f$ est \textbf{strictement croissante} sur $\mathbb{R}$.

\textbf{Limites aux bornes.} Comme $-1 \leq \sin x \leq 1$, on a $2x - 1 \leq f(x) \leq 2x + 1$. Donc :
\[ \lim_{x \to +\infty} f(x) = +\infty \quad \text{et} \quad \lim_{x \to -\infty} f(x) = -\infty. \]

\textbf{Conclusion.} $f$ est continue et strictement croissante sur $\mathbb{R}$ : d'après le théorème de la bijection, $f$ réalise une bijection de $\mathbb{R}$ sur $f(\mathbb{R}) = \mathbb{R}$.
\item Sur $[0\,;\,\pi]$, $f$ est continue et strictement croissante, donc :
\[ f\bigl([0\,;\,\pi]\bigr) = \bigl[f(0)\,;\,f(\pi)\bigr] = \bigl[0\,;\,2\pi\bigr], \]
car $f(0) = 0 + \sin 0 = 0$ et $f(\pi) = 2\pi + \sin \pi = 2\pi$.
\end{enumerate}
\end{exoresolu}

\courssub{Restrictions des fonctions trigonométriques : arcsin, arccos, arctan}

Les fonctions $\sin$, $\cos$ et $\tan$ ne sont \textbf{pas} monotones sur $\mathbb{R}$ tout entier (elles oscillent) : elles n'y sont donc pas bijectives. L'idée, historique et fondamentale, consiste à les \cle{restreindre} à un intervalle où elles sont continues et strictement monotones, afin de pouvoir leur associer une fonction réciproque.

\begin{propriete}[Fonctions trigonométriques réciproques (admises)]
\begin{itemize}
\item La fonction $\sin$ restreinte à $\left[-\dfrac{\pi}{2}\,;\,\dfrac{\pi}{2}\right]$ est continue et strictement croissante, à valeurs dans $[-1\,;\,1]$. Sa réciproque est \cle{arcsin} : $[-1\,;\,1] \to \left[-\dfrac{\pi}{2}\,;\,\dfrac{\pi}{2}\right]$, strictement croissante.
\item La fonction $\cos$ restreinte à $[0\,;\,\pi]$ est continue et strictement décroissante, à valeurs dans $[-1\,;\,1]$. Sa réciproque est \cle{arccos} : $[-1\,;\,1] \to [0\,;\,\pi]$, strictement décroissante.
\item La fonction $\tan$ restreinte à $\left]-\dfrac{\pi}{2}\,;\,\dfrac{\pi}{2}\right[$ est continue et strictement croissante, à valeurs dans $\mathbb{R}$. Sa réciproque est \cle{arctan} : $\mathbb{R} \to \left]-\dfrac{\pi}{2}\,;\,\dfrac{\pi}{2}\right[$, strictement croissante.
\end{itemize}
Ainsi : $\arcsin\left(\dfrac{1}{2}\right) = \dfrac{\pi}{6}$, $\arccos(0) = \dfrac{\pi}{2}$, $\arctan(1) = \dfrac{\pi}{4}$.
\end{propriete}

\begin{popfigure}
\begin{center}
\begin{tikzpicture}[scale=1.0]
\begin{axis}[
  axis lines=middle, xlabel={$x$}, ylabel={$y$},
  xmin=-4.4, xmax=4.4, ymin=-2.1, ymax=2.1,
  xtick={-3.1416,-1.5708,1.5708,3.1416},
  xticklabels={$-\pi$,$-\frac{\pi}{2}$,$\frac{\pi}{2}$,$\pi$},
  ytick={-1,1}, width=0.85\linewidth]
  \addplot[popmuted,dashed,domain=-4.2:4.2,samples=200]{sin(deg(x))};
  \addplot[popblue,very thick,domain=-1.5708:1.5708,samples=100]{sin(deg(x))};
  \addplot[poporange,very thick,domain=-1:1,samples=100]{rad(asin(x))};
  \node[popmuted] at (axis cs:3.6,0.75) {\small $\sin x$};
  \node[popblue] at (axis cs:1.05,1.35) {\small $\sin$ restreint};
  \node[poporange] at (axis cs:-0.55,-1.5) {\small $\arcsin$};
\end{axis}
\end{tikzpicture}
\end{center}
\legende{En gras bleu : la restriction de $\sin$ à $\left[-\frac{\pi}{2}\,;\,\frac{\pi}{2}\right]$, strictement croissante. En orange : sa réciproque $\arcsin$, définie sur $[-1\,;\,1]$.}
\end{popfigure}

\begin{center}
\rowcolors{1}{popcream}{white}
\begin{tabularx}{\linewidth}{|l|X|X|l|}
\hline
\rowcolor{popblueL} \textbf{Fonction} & \textbf{Intervalle source $I$} & \textbf{Intervalle image $f(I)$} & \textbf{Réciproque} \\
\hline
$\sin$ restreint & $\left[-\dfrac{\pi}{2}\,;\,\dfrac{\pi}{2}\right]$ & $[-1\,;\,1]$ & $\arcsin$ \\
\hline
$\cos$ restreint & $[0\,;\,\pi]$ & $[-1\,;\,1]$ & $\arccos$ \\
\hline
$\tan$ restreint & $\left]-\dfrac{\pi}{2}\,;\,\dfrac{\pi}{2}\right[$ & $\mathbb{R}$ & $\arctan$ \\
\hline
$x \mapsto x^2$ restreint & $[0\,;\,+\infty[$ & $[0\,;\,+\infty[$ & $\sqrt{x}$ \\
\hline
$x \mapsto x^n$ ($n$ impair) & $\mathbb{R}$ & $\mathbb{R}$ & $\sqrt[n]{x}$ \\
\hline
$\exp$ & $\mathbb{R}$ & $]0\,;\,+\infty[$ & $\ln$ \\
\hline
\end{tabularx}
\end{center}

\courssub{Expression analytique de la réciproque}

Dans quelques cas privilégiés, on connaît une \emph{formule explicite} pour la réciproque :
\begin{itemize}
\item la réciproque de $x \mapsto x^2$ sur $[0\,;\,+\infty[$ est $x \mapsto \sqrt{x}$ ;
\item la réciproque de $x \mapsto x^3$ sur $\mathbb{R}$ est $x \mapsto \sqrt[3]{x}$ ;
\item la réciproque de $\exp$ sur $\mathbb{R}$ est $\ln$ sur $]0\,;\,+\infty[$ ;
\item les réciproques des restrictions trigonométriques sont $\arcsin$, $\arccos$, $\arctan$.
\end{itemize}

\begin{attention}[Le cas général : pas de formule explicite]
Pour la plupart des bijections, il est \textbf{impossible} d'exprimer $f^{-1}(x)$ par une formule. Par exemple, pour $f(x) = 2x + \sin x$, résoudre $y = 2x + \sin x$ en $x$ ne se fait pas algébriquement. On se contente alors :
\begin{itemize}
\item d'une \textbf{lecture graphique} (symétrie par rapport à $y=x$) ;
\item d'une \textbf{résolution numérique} : encadrement de la solution de $f(x)=k$ par balayage à la calculatrice, ou méthode de dichotomie.
\end{itemize}
L'existence de $f^{-1}$ est garantie par le théorème, même sans formule : c'est toute la puissance du raisonnement.
\end{attention}

\begin{exercice}[Application au réglage d'une machine]
La température (en degrés Celsius) d'un four après $t$ minutes est modélisée par $g(t) = 20 + 180\left(1 - e^{-0{,}1 t}\right)$ pour $t \geq 0$.
\begin{enumerate}
\item Étudier la continuité et la monotonie de $g$ sur $[0\,;\,+\infty[$.
\item Montrer que l'équation $g(t) = 150$ admet une unique solution $t_0$ et donner un encadrement de $t_0$ à la minute près.
\end{enumerate}
\end{exercice}

\begin{corrige}[Exercice : réglage d'une machine]
\begin{enumerate}
\item $g$ est continue sur $[0\,;\,+\infty[$ comme composée et somme de fonctions continues. $g'(t) = 18\,e^{-0{,}1 t} > 0$ pour tout $t \geq 0$ : $g$ est strictement croissante.
\item $g(0) = 20$ et $\displaystyle\lim_{t \to +\infty} g(t) = 200$. Comme $150 \in [20\,;\,200[$, le théorème de la bijection garantit une unique solution $t_0$. Par balayage : $g(12) \approx 145{,}8$ ; $g(13) \approx 151{,}0$ ; on obtient $12 < t_0 < 13$. Le four atteint $\SI{150}{\celsius}$ entre la 12\textsuperscript{ème} et la 13\textsuperscript{ème} minute.
  \par \textbf{Valeur exacte :} $1 - e^{-0{,}1t} = \frac{130}{180} = \frac{13}{18}$, d'où $e^{-0{,}1t} = \frac{5}{18}$ et $t_0 = 10\ln\frac{18}{5} \approx 12{,}8$ min, ce qui confirme l'encadrement.
\end{enumerate}
\end{corrige}

\begin{savaistu}[D'où viennent arcsin, arccos et arctan ?]
Le théorème de la bijection est la base théorique de la définition des fonctions $\arcsin$, $\arccos$ et $\arctan$ : sans lui, ces fonctions n'existeraient tout simplement pas ! En classe de Terminale technique, on utilise leurs propriétés (domaines, variations, valeurs remarquables) sans redémontrer leur existence. Ces fonctions sont omniprésentes dans les métiers techniques : calcul d'angles d'inclinaison en topographie et en charpente, déphasage en électricité (circuits en courant alternatif), trajectoire et balistique, réglage des axes d'une machine à commande numérique. Chaque fois qu'un technicien « remonte » d'une valeur de sinus, cosinus ou tangente à un angle, il utilise une fonction réciproque.
\end{savaistu}

\begin{aretenir}[L'essentiel de la section]
\begin{itemize}
\item $f$ continue et strictement monotone sur $I$ $\Rightarrow$ $f$ est une bijection de $I$ sur l'intervalle $f(I)$.
\item L'équation $f(x) = k$, avec $k \in f(I)$, admet alors une \textbf{unique} solution dans $I$.
\item $f^{-1}$ est définie sur $f(I)$, continue, et a le même sens de variation que $f$ ; sa courbe est symétrique de celle de $f$ par rapport à $y = x$.
\item $\sin$, $\cos$, $\tan$ deviennent bijectives par restriction à $\left[-\tfrac{\pi}{2}\,;\,\tfrac{\pi}{2}\right]$, $[0\,;\,\pi]$, $\left]-\tfrac{\pi}{2}\,;\,\tfrac{\pi}{2}\right[$ : leurs réciproques sont $\arcsin$, $\arccos$, $\arctan$.
\item Ne jamais confondre $f^{-1}(x)$ (réciproque) et $\dfrac{1}{f(x)}$ (inverse).
\end{itemize}
\end{aretenir}

\coursec{Limites des fonctions trigonométriques : limites remarquables et encadrement}

Après avoir étudié les propriétés des fonctions continues et strictement monotones, et vu comment la continuité permet d'inverser une bijection, nous abordons maintenant un chapitre technique incontournable : le calcul des limites des fonctions trigonométriques. En Terminale technique, que ce soit pour modéliser un phénomène vibratoire (mécanique), un signal alternatif (électricité) ou une oscillation amortie, la maîtrise des limites en $0$ et en $\pm\infty$ de $\sin$, $\cos$ et $\tan$ est un outil de base. La clé de voûte de tout cet édifice est la limite fondamentale $\lim_{x \to 0} \frac{\sin x}{x} = 1$, à condition de travailler en \cle{radians}.

\courssub{La limite fondamentale $\boldsymbol{\lim_{x \to 0} \frac{\sin x}{x} = 1}$ et sa démonstration géométrique}

C'est le résultat le plus important de ce chapitre. Il ne s'obtient pas par un calcul algébrique simple, mais par un raisonnement géométrique d'encadrement (théorème des gendarmes).

\begin{experience}[Démonstration géométrique de $\lim_{x \to 0} \frac{\sin x}{x} = 1$]
Considérons le cercle trigonométrique de centre $O$ et de rayon $1$. Soit $x \in ]0, \frac{\pi}{2}[$ un angle en \cle{radians}.
\begin{enumerate}
    \item On place le point $M$ sur le cercle tel que $\angle(\vec{u}, \vec{OM}) = x$. Ses coordonnées sont $(\cos x, \sin x)$.
    \item On trace la tangente au cercle en $A(1,0)$. Elle coupe la droite $(OM)$ en $B$.
    \item On note $H$ le projeté orthogonal de $M$ sur l'axe des abscisses $(OA)$.
\end{enumerate}
Comparons les aires de trois figures emboîtées :
\begin{itemize}
    \item Triangle rectangle $OAM$ (dans le secteur) : $\mathcal{A}_{\triangle OAM} = \frac{1}{2} \times 1 \times \sin x = \frac{\sin x}{2}$.
    \item Secteur circulaire $OAM$ d'angle $x$ : $\mathcal{A}_{\text{secteur}} = \frac{x}{2\pi} \times \pi \times 1^2 = \frac{x}{2}$.
    \item Triangle rectangle $OAB$ (contenant le secteur) : $\mathcal{A}_{\triangle OAB} = \frac{1}{2} \times 1 \times \tan x = \frac{\tan x}{2}$.
\end{itemize}
Par inclusion, on a l'inégalité d'aires : $\mathcal{A}_{\triangle OAM} \le \mathcal{A}_{\text{secteur}} \le \mathcal{A}_{\triangle OAB}$.
Soit : $\frac{\sin x}{2} \le \frac{x}{2} \le \frac{\tan x}{2}$.
En multipliant par $2$ et en divisant par $\sin x > 0$ : $1 \le \frac{x}{\sin x} \le \frac{1}{\cos x}$.
En inversant (les termes sont positifs) : $\cos x \le \frac{\sin x}{x} \le 1$.
Or $\lim_{x \to 0^+} \cos x = 1$. Par le théorème des gendarmes : $\lim_{x \to 0^+} \frac{\sin x}{x} = 1$.
La fonction $x \mapsto \frac{\sin x}{x}$ est paire, donc la limite à gauche est aussi $1$.
\emph{Conclusion :} $\lim_{x \to 0} \frac{\sin x}{x} = 1$ (avec $x$ en radians).
\end{experience}

\begin{popfigure}
\begin{tikzpicture}[scale=1.5, >=Stealth]
    % Circle
    \draw[popblue, thick] (0,0) circle (1);
    % Axes
    \draw[->] (-1.3,0) -- (1.5,0) node[right] {$x$};
    \draw[->] (0,-0.3) -- (0,1.5) node[above] {$y$};
    % Points
    \coordinate (O) at (0,0);
    \coordinate (A) at (1,0);
    \coordinate (M) at (45:1); % angle 45 deg = pi/4 rad
    \coordinate (H) at (0.707, 0);
    % Tangent at A: x=1. Line OM: y=x. Intersection B(1,1)
    \coordinate (B) at (1,1);
    % Lines
    \draw[popdark, thick] (O) -- (M) node[midway, above left] {$1$};
    \draw[popdark, thick] (O) -- (B);
    \draw[popdark, thick] (A) -- (B);
    \draw[popdark, dashed] (M) -- (H);
    % Angle arc
    \draw[poporange, thick] (0.2,0) arc (0:45:0.2) node[midway, right, poporange] {$x$};
    % Labels
    \node[below left] at (O) {$O$};
    \node[below] at (A) {$A(1,0)$};
    \node[above right] at (M) {$M(\cos x, \sin x)$};
    \node[above right] at (B) {$B(1, \tan x)$};
    \node[below] at (H) {$H$};
    % Right angle marks
    \draw (0.9,0) -- (0.9,0.1) -- (1,0.1);
    \draw (0.707, 0) -- (0.707, 0.05) -- (0.757, 0.05);
    % Area labels
    \node[popblue] at (0.3, 0.15) {$\frac{\sin x}{2}$};
    \node[poporange] at (0.5, 0.35) {$\frac{x}{2}$};
    \node[popgreen] at (0.8, 0.6) {$\frac{\tan x}{2}$};
\end{tikzpicture}
\legende{Encadrement géométrique des aires : $\triangle OAM \subset \text{Secteur } OAM \subset \triangle OAB$.}
\end{popfigure}

\begin{attention}[Piège mortel : Degrés vs Radians]
La limite $\lim_{x \to 0} \frac{\sin x}{x} = 1$ est \textbf{valable uniquement si $x$ est exprimé en radians}.
Si $x$ est en degrés, $\sin(x^\circ) = \sin\left(\frac{\pi x}{180}\right)$. Alors $\frac{\sin(x^\circ)}{x} = \frac{\pi}{180} \frac{\sin(\pi x/180)}{\pi x/180} \to \frac{\pi}{180} \approx 0,01745 \neq 1$.
\textbf{Au Probatoire / Baccalauréat technique :} Vérifiez systématiquement l'unité. Si l'énoncé ne précise pas, c'est \textbf{radians} par défaut. Une calculatrice en mode DEG donnera un résultat faux pour ce type de limite.
\end{attention}

\courssub{Conséquences immédiates : les trois limites remarquables}

À partir de la limite fondamentale, on déduit deux autres limites indispensables, souvent utilisées sans redémonstration.

\begin{propriete}[Limites remarquables trigonométriques en $0$ (à connaître par cœur)]
Toutes les limites sont prises pour $x \to 0$ avec $x$ en \textbf{radians}.
\begin{enumerate}
    \item $\lim_{x \to 0} \frac{\sin x}{x} = 1$.
    \item $\lim_{x \to 0} \frac{\tan x}{x} = 1$.
        \emph{Preuve :} $\frac{\tan x}{x} = \frac{\sin x}{x} \cdot \frac{1}{\cos x} \to 1 \times 1 = 1$.
    \item $\lim_{x \to 0} \frac{1 - \cos x}{x^2} = \frac{1}{2}$.
        \emph{Preuve :} $1 - \cos x = 2\sin^2\frac{x}{2}$. Donc $\frac{1-\cos x}{x^2} = 2 \frac{\sin^2(x/2)}{x^2} = \frac{1}{2} \left(\frac{\sin(x/2)}{x/2}\right)^2 \to \frac{1}{2} \times 1^2 = \frac{1}{2}$.
\end{enumerate}
\end{propriete}

\begin{aretenir}[Boîte à outils : Les 3 limites de base]
\begin{center}
\begin{tabular}{|c|c|}\hline
\textbf{Forme} & \textbf{Limite ($x \to 0$, rad)} \\ \hline
$\dfrac{\sin x}{x}$ & $1$ \\ \hline
$\dfrac{\tan x}{x}$ & $1$ \\ \hline
$\dfrac{1 - \cos x}{x^2}$ & $\dfrac{1}{2}$ \\ \hline
\end{tabular}
\end{center}
\textbf{Équivalents au voisinage de $0$ (Développements Limités d'ordre 1) :}
$\sin x \sim x$, \quad $\tan x \sim x$, \quad $1 - \cos x \sim \frac{x^2}{2}$.
Ces équivalents permettent de calculer très vite les limites de type $\frac{0}{0}$.
\end{aretenir}

\courssub{Changement de variable : la technique du « $\boldsymbol{u}$ »}

Dans les exercices, l'argument de la fonction trigonométrique n'est jamais simplement $x$, mais une fonction $u(x)$ qui tend vers $0$. Il faut \textbf{forcer l'apparition de $\frac{\sin u}{u}$ ou $\frac{1-\cos u}{u^2}$}.

\begin{methode}[Calculer une limite trigonométrique en $0$ par changement de variable]
\textbf{Objectif :} Calculer $\lim_{x \to a} f(x)$ où $f$ contient $\sin(u(x))$, $\tan(u(x))$ ou $1-\cos(u(x))$ avec $u(x) \to 0$.
\begin{enumerate}
    \item \textbf{Identifier $u(x)$} tel que $u(x) \to 0$ quand $x \to a$.
    \item \textbf{Multiplier et diviser} par $u(x)$ (ou $u(x)^2$ pour $1-\cos$) pour faire apparaître la forme remarquable.
    \item \textbf{Appliquer la limite remarquable} : $\frac{\sin u}{u} \to 1$, $\frac{\tan u}{u} \to 1$, $\frac{1-\cos u}{u^2} \to \frac{1}{2}$.
    \item \textbf{Calculer la limite du reste} (partie algébrique).
\end{enumerate}
\end{methode}

\begin{exemplebox}[Application : $\lim_{x \to 0} \frac{\sin(3x)}{5x}$ et $\lim_{x \to 0} \frac{1-\cos(2x)}{x^2}$]
\begin{enumerate}
    \item $\lim_{x \to 0} \frac{\sin(3x)}{5x}$.
        Posons $u = 3x \to 0$. On veut $\frac{\sin u}{u}$.
        $\frac{\sin(3x)}{5x} = \frac{\sin(3x)}{3x} \times \frac{3}{5} \to 1 \times \frac{3}{5} = \boxed{\frac{3}{5}}$.
        \textbf{Règle rapide :} $\lim_{x \to 0} \frac{\sin(ax)}{bx} = \frac{a}{b}$.
    \item $\lim_{x \to 0} \frac{1-\cos(2x)}{x^2}$.
        Posons $u = 2x \to 0$. On veut $\frac{1-\cos u}{u^2}$.
        $\frac{1-\cos(2x)}{x^2} = \frac{1-\cos(2x)}{(2x)^2} \times \frac{(2x)^2}{x^2} = \frac{1-\cos u}{u^2} \times 4 \to \frac{1}{2} \times 4 = \boxed{2}$.
        \textbf{Règle rapide :} $\lim_{x \to 0} \frac{1-\cos(ax)}{bx^2} = \frac{a^2}{2b}$.
\end{enumerate}
\end{exemplebox}

\begin{attention}[Erreurs fréquentes sur le changement de variable]
\begin{itemize}
    \item \textbf{Oublier de multiplier ET diviser :} Écrire $\frac{\sin(3x)}{x} = \frac{\sin(3x)}{3x}$ (faux, il manque le facteur $3$).
    \item \textbf{Appliquer $\frac{\sin x}{x} \to 1$ quand $x \to \infty$ :} $\lim_{x \to +\infty} \frac{\sin x}{x} = 0$ (théorème des gendarmes car $|\sin x| \le 1$), mais $\frac{\sin x}{x}$ n'a pas de forme $\frac{\sin u}{u}$ avec $u\to 0$ ici.
    \item \textbf{Confondre $1-\cos x \sim x^2/2$ et $\cos x \sim 1$ :} $\cos x - 1 \sim -x^2/2$, mais on ne peut pas remplacer $\cos x$ par $1$ dans un numérateur si le dénominateur est $x$ ou $x^2$ (forme indéterminée).
\end{itemize}
\end{attention}

\courssub{Limites en $\boldsymbol{\pm\infty}$ : oscillation et encadrement}

Contrairement aux fonctions polynomiales ou rationnelles, $\sin x$ et $\cos x$ n'ont \textbf{pas de limite} en $+\infty$ ni en $-\infty$.

\begin{propriete}[Comportement en $\pm\infty$]
\begin{itemize}
    \item $\sin x$ et $\cos x$ : \textbf{Pas de limite} en $\pm\infty$ (oscillation bornée entre $-1$ et $1$).
    \item $\tan x$ : \textbf{Pas de limite} en $\pm\infty$ (périodicité $\pi$ + asymptotes verticales en $\frac{\pi}{2} + k\pi$).
    \item $\frac{\sin x}{x}$ et $\frac{\cos x}{x}$ : Limite \textbf{$0$} en $\pm\infty$ (numérateur borné, dénominateur $\to \infty$).
\end{itemize}
\end{propriete}

C'est ici que le \textbf{Théorème de l'encadrement (Gendarmes)} devient un outil de calcul puissant.

\begin{propriete}[Théorème de l'encadrement (Rappel)]
Soient $f, g, h$ définies sur un intervalle $I$ (sauf peut-être en $a$).
Si $g(x) \le f(x) \le h(x)$ sur $I$ et $\lim_{x \to a} g(x) = \lim_{x \to a} h(x) = \ell$ (fini ou infini), alors $\lim_{x \to a} f(x) = \ell$.
\end{propriete}

\begin{exoresolu}[Application des gendarmes : $x \sin(1/x)$ en $0$ et $x \cos x$ en $+\infty$]
\begin{enumerate}
    \item \textbf{Étudier $\lim_{x \to 0} x \sin\left(\frac{1}{x}\right)$.}
        On sait que $\forall x \neq 0, \quad -1 \le \sin\left(\frac{1}{x}\right) \le 1$.
        \textbf{Cas $x > 0$ :} En multipliant par $x > 0$ : $-x \le x \sin(1/x) \le x$.
        Comme $\lim_{x \to 0^+} (-x) = \lim_{x \to 0^+} x = 0$, par les gendarmes : $\lim_{x \to 0^+} x \sin(1/x) = 0$.
        \textbf{Cas $x < 0$ :} En multipliant par $x < 0$, l'inégalité s'inverse : $x \le x \sin(1/x) \le -x$.
        Les bornes tendent vers $0$. Donc $\lim_{x \to 0^-} x \sin(1/x) = 0$.
        \textbf{Conclusion :} $\lim_{x \to 0} x \sin(1/x) = \boxed{0}$.
    \item \textbf{Étudier $\lim_{x \to +\infty} \frac{\cos x}{x}$.}
        $-1 \le \cos x \le 1$. Pour $x > 0$ : $-\frac{1}{x} \le \frac{\cos x}{x} \le \frac{1}{x}$.
        Bornes $\to 0$. Donc $\lim_{x \to +\infty} \frac{\cos x}{x} = \boxed{0}$.
\end{enumerate}
\end{exoresolu}

\begin{popfigure}
\begin{tikzpicture}
\begin{axis}[
    width=\linewidth, height=6cm,
    axis lines=middle,
    xlabel={$x$}, ylabel={$y$},
    xmin=-1.5, xmax=1.5, ymin=-1.2, ymax=1.2,
    xtick={-1,-0.5,0.5,1}, ytick={-1,-0.5,0.5,1},
    tick label style={font=\small},
    domain=-1.5:1.5, samples=1000,
    clip=true
]
% Enveloppes y=x and y=-x
\addplot[poporange, dashed, thick, domain=0.01:1.5] {x};
\addplot[poporange, dashed, thick, domain=-1.5:-0.01] {x};
\addplot[poporange, dashed, thick, domain=0.01:1.5] {-x};
\addplot[poporange, dashed, thick, domain=-1.5:-0.01] {-x};
% Function x*sin(1/x)
\addplot[popblue, very thick, domain=0.01:1.5] {x*sin(1/x)};
\addplot[popblue, very thick, domain=-1.5:-0.01] {x*sin(1/x)};
% Point at 0 (hole)
\addplot[only marks, mark=*, mark size=2, popblue, fill=white] coordinates {(0,0)};
\node[popblue, anchor=south west] at (0.1, 0.1) {$y=x\sin(1/x)$};
\node[poporange, anchor=north west] at (0.1, 1.0) {$y=x$};
\node[poporange, anchor=south west] at (0.1, -1.0) {$y=-x$};
\end{axis}
\end{tikzpicture}
\legende{Graphique de $y = x \sin(1/x)$ encadré par $y=x$ et $y=-x$. La fonction oscille infiniment mais l'amplitude tend vers $0$.}
\end{popfigure}

\courssub{Limites composées et continuité : $\boldsymbol{\lim \sin(u(x)) = \sin(\lim u(x))}$}

Les fonctions $\sin$, $\cos$, $\tan$ sont \textbf{continues} sur leur ensemble de définition. Cela permet de « faire entrer la limite à l'intérieur » sous réserve que la limite interne appartienne au domaine de définition.

\begin{propriete}[Limite de fonction trigonométrique composée]
Soit $u$ une fonction telle que $\lim_{x \to a} u(x) = L$.
\begin{itemize}
    \item Si $L \in \mathbb{R}$, alors $\lim_{x \to a} \sin(u(x)) = \sin L$ et $\lim_{x \to a} \cos(u(x)) = \cos L$ (continuité sur $\mathbb{R}$).
    \item Si $L \in \mathbb{R} \setminus \{\frac{\pi}{2} + k\pi\}$, alors $\lim_{x \to a} \tan(u(x)) = \tan L$ (continuité sur son domaine).
    \item \textbf{Attention :} Si $L = \frac{\pi}{2} + k\pi$, $\tan(u(x))$ n'a pas de limite finie (forme $\frac{1}{0}$).
\end{itemize}
\end{propriete}

\begin{exemplebox}[Limites composées : cas simples et pièges]
\begin{enumerate}
    \item $\lim_{x \to 1} \sin(x^2 - 1)$. $u(x)=x^2-1 \to 0$. $\sin(0)=0$. \textbf{Résultat : $0$.}
    \item $\lim_{x \to 0} \cos\left(\frac{1}{x}\right)$. $u(x)=1/x$ n'a pas de limite en $0$ ($\pm\infty$). $\cos$ oscille. \textbf{Pas de limite.}
    \item $\lim_{x \to 0} \tan\left(\frac{\pi}{4} + x\right)$. $u(x) \to \pi/4$. $\tan(\pi/4)=1$. \textbf{Résultat : $1$.}
    \item $\lim_{x \to 0} \tan\left(\frac{\pi}{2} - x\right)$. $u(x) \to \pi/2$ (valeur interdite pour $\tan$).
        $\tan(\pi/2 - x) = \cot x = \frac{\cos x}{\sin x} \sim \frac{1}{x} \to +\infty$ (pour $x>0$). \textbf{Limite infinie.}
\end{enumerate}
\end{exemplebox}

\courssub{Calcul de limites par Développements Limités (DL) au premier ordre — Outil Premium}

Bien que le programme officiel parle de « limites remarquables », l'utilisation des \textbf{équivalents} (DL d'ordre 1) est la méthode standard en Terminale technique pour traiter rapidement les formes indéterminées $\frac{0}{0}$ ou $\frac{\infty}{\infty}$ impliquant des fonctions trigonométriques. C'est la méthode attendue dans les concours et le Baccalauréat pour aller vite et juste.

\begin{aretenir}[Règles des équivalents (au voisinage de $0$, radians)]
Si $u(x) \to 0$ quand $x \to a$ :
\begin{align*}
    \sin(u(x)) &\sim u(x) \\
    \tan(u(x)) &\sim u(x) \\
    1 - \cos(u(x)) &\sim \frac{u(x)^2}{2}
\end{align*}
\textbf{Principe :} Dans une limite de quotient, on peut remplacer une fonction par son équivalent \textbf{uniquement si c'est une factorisation (produit/quotient)}. Interdit dans une somme/soustraction (sauf factorisation préalable).
\end{aretenir}

\begin{methode}[Calculer une limite trigonométrique par équivalents (DL1)]
\begin{enumerate}
    \item Vérifier la forme indéterminée ($0/0$ ou $\infty/\infty$).
    \item Identifier les « petits » arguments $u(x) \to 0$.
    \item Remplacer $\sin(u)$ par $u$, $\tan(u)$ par $u$, $1-\cos(u)$ par $u^2/2$.
    \item Simplifier algébriquement et calculer la limite.
\end{enumerate}
\end{methode}

\begin{exoresolu}[Exemple synthèse : $\lim_{x \to 0} \frac{\tan(2x) - \sin(3x)}{x^3}$]
\textbf{Énoncé :} Calculer $\lim_{x \to 0} \frac{\tan(2x) - \sin(3x)}{x^3}$.

\textbf{Solution détaillée :}
\begin{enumerate}
    \item \textbf{Forme :} Numérateur $\tan 0 - \sin 0 = 0$, Dénominateur $0$. Forme $\frac{0}{0}$.
    \item \textbf{Équivalents (DL1) :} $x \to 0 \implies 2x \to 0, 3x \to 0$.
        $\tan(2x) \sim 2x$ \quad et \quad $\sin(3x) \sim 3x$.
    \item \textbf{Piège :} Si on remplace directement : $\frac{2x - 3x}{x^3} = \frac{-x}{x^3} = -\frac{1}{x^2} \to -\infty$.
        \textbf{Attention !} Le numérateur est une \textbf{différence}. Les équivalents d'ordre 1 ne suffisent pas car les termes dominants ($2x$ et $3x$) s'annulent partiellement. Il faut aller à l'ordre supérieur (DL2 ou DL3) ou utiliser les limites remarquables précisément.
    \item \textbf{Méthode rigoureuse (DL3) :}
        Rappels : $\sin u = u - \frac{u^3}{6} + o(u^3)$, $\tan u = u + \frac{u^3}{3} + o(u^3)$.
        $\tan(2x) = 2x + \frac{(2x)^3}{3} + o(x^3) = 2x + \frac{8}{3}x^3 + o(x^3)$.
        $\sin(3x) = 3x - \frac{(3x)^3}{6} + o(x^3) = 3x - \frac{27}{6}x^3 + o(x^3) = 3x - \frac{9}{2}x^3 + o(x^3)$.
        Numérateur $= (2x - 3x) + \left(\frac{8}{3} + \frac{9}{2}\right)x^3 + o(x^3) = -x + \left(\frac{16+27}{6}\right)x^3 + o(x^3) = -x + \frac{43}{6}x^3 + o(x^3)$.
        Quotient $= \frac{-x + \frac{43}{6}x^3}{x^3} = -\frac{1}{x^2} + \frac{43}{6} \to -\infty$.
    \item \textbf{Conclusion :} La limite est $\boxed{-\infty}$ (signe $-$ car $x^2>0$).
\end{enumerate}
\textbf{Note pédagogique :} Cet exemple montre la limite des équivalents d'ordre 1 sur une différence. Au Probatoire, si la limite est infinie, l'ordre 1 suffit souvent à voir le signe de l'infini, mais ici l'annulation des termes d'ordre 1 impose de monter d'un cran. En Terminale Tech, on accepte souvent l'argument : « Au numérateur, $\tan(2x)-\sin(3x) \sim 2x-3x = -x$, donc la fraction $\sim -x/x^3 = -1/x^2 \to -\infty$ », à condition de justifier que le reste est négligeable devant $-x$.
\end{exoresolu}

\courssub{Représentation graphique et lecture de limites}

La représentation graphique permet de visualiser les comportements asymptotiques.

\begin{popfigure}
\begin{tikzpicture}
\begin{axis}[
    width=\linewidth, height=7cm,
    axis lines=middle,
    xlabel={$x$ (rad)}, ylabel={$y$},
    xmin=-2*pi, xmax=2*pi,
    ymin=-1.5, ymax=1.5,
    xtick={-6.283,-3.1415,0,3.1415,6.283},
    xticklabels={$-2\pi$, $-\pi$, $0$, $\pi$, $2\pi$},
    ytick={-1,0,1},
    tick label style={font=\small},
    domain=-2*pi:2*pi, samples=400,
    clip=false
]
% sin(x)/x
\addplot[popblue, very thick, domain=-2*pi:-0.1] {sin(deg(x))/x};
\addplot[popblue, very thick, domain=0.1:2*pi] {sin(deg(x))/x};
% Hole at 0
\addplot[only marks, mark=*, mark size=2.5, popblue, fill=white] coordinates {(0,1)};
\node[popblue, anchor=south west] at (0.3, 1.1) {$y=\frac{\sin x}{x}$};
\node[popblue, anchor=north west] at (0.3, 0.9) {Trou en $(0,1)$};
% Asymptotes horizontal y=0
\addplot[popmuted, dashed, domain=-2*pi:2*pi] {0};
\node[popmuted, anchor=south east] at (6, 0.1) {$y=0$ (asymptote)};
\end{axis}
\end{tikzpicture}
\legende{Courbe de $y = \frac{\sin x}{x}$. Limite $1$ en $0$ (trou), limite $0$ en $\pm\infty$ (amortissement). Zéros en $k\pi$ ($k \neq 0$).}
\end{popfigure}

\courssub{Application technique : Approximation des petits angles}

En technique, la limite $\frac{\sin x}{x} \to 0$ justifie l'approximation $\sin \theta \approx \theta$ (radians) pour les petits angles. C'est la base de la linéarisation de nombreux systèmes physiques.

\begin{exemplebox}[Pendule simple et signal électrique]
\begin{enumerate}
    \item \textbf{Pendule simple (Mécanique) :} L'équation différentielle est $\ddot{\theta} + \frac{g}{l}\sin\theta = 0$. Pour de petites oscillations ($\theta \approx 0$), $\sin\theta \sim \theta$. L'équation devient $\ddot{\theta} + \frac{g}{l}\theta = 0$ (oscillateur harmonique). La période est $T = 2\pi\sqrt{\frac{l}{g}}$, indépendante de l'amplitude (isochronisme).
    \item \textbf{Circuit RL série (Électricité) :} Sous tension alternative $u(t) = U\sqrt{2}\sin(\omega t)$, le courant est $i(t) = I\sqrt{2}\sin(\omega t - \varphi)$. Pour $\omega \to 0$ (régime continu), $\varphi = \arctan(\frac{L\omega}{R}) \sim \frac{L\omega}{R} \to 0$. Le circuit se comporte comme une résistance pure.
\end{enumerate}
\end{exemplebox}

\begin{savaistu}[La clé du calcul différentiel]
La limite $\lim_{x \to 0} \frac{\sin x}{x} = 1$ n'est pas qu'un exercice de limite. C'est \textbf{la} limite qui permet de démontrer que la dérivée de $\sin x$ est $\cos x$ :
$(\sin x)' = \lim_{h \to 0} \frac{\sin(x+h)-\sin h}{h} = \lim_{h \to 0} \frac{\sin x \cos h + \cos x \sin h - \sin x}{h}$
$= \sin x \underbrace{\lim_{h \to 0} \frac{\cos h - 1}{h}}_{=0} + \cos x \underbrace{\lim_{h \to 0} \frac{\sin h}{h}}_{=1} = \cos x$.
Sans cette limite (et le passage aux radians), tout le calcul différentiel des fonctions trigonométriques s'effondre. C'est pour cela que les radians sont l'unité \textbf{naturelle} des mathématiques et de la physique (pulsation $\omega$ en rad/s).
\end{savaistu}

\courssub{Entraînement : Limites trigonométriques}

\begin{exercice}[Exercices d'application]
Calculer les limites suivantes (préciser $+\infty$, $-\infty$ ou « n'existe pas »).
\begin{enumerate}
    \item $\lim_{x \to 0} \frac{\sin(5x)}{\tan(2x)}$
    \item $\lim_{x \to 0} \frac{1 - \cos(4x)}{x^2}$
    \item $\lim_{x \to +\infty} \frac{\sin(x^2)}{x}$
    \item $\lim_{x \to 0} \frac{\tan(x) - \sin(x)}{x^3}$ \textit{(Indice : $\tan x - \sin x = \sin x (\frac{1}{\cos x} - 1)$)}
    \item $\lim_{x \to \pi/2} \frac{\cos x}{x - \pi/2}$ \textit{(Indice : poser $h = x - \pi/2$)}
\end{enumerate}
\end{exercice}

\begin{corrige}[Exercices d'application]
\begin{enumerate}
    \item $\frac{\sin(5x)}{\tan(2x)} = \frac{\sin(5x)}{5x} \frac{5x}{2x} \frac{2x}{\tan(2x)} \to 1 \times \frac{5}{2} \times 1 = \boxed{2,5}$.
    \item $\frac{1-\cos(4x)}{x^2} = \frac{1-\cos(4x)}{(4x)^2} \times 16 \to \frac{1}{2} \times 16 = \boxed{8}$.
    \item $|\sin(x^2)| \le 1$, $x \to +\infty$. Par gendarmes : $-\frac{1}{x} \le \frac{\sin(x^2)}{x} \le \frac{1}{x} \to 0$. \textbf{Limite $\boxed{0}$}.
    \item $\tan x - \sin x = \sin x \frac{1-\cos x}{\cos x} \sim x \cdot \frac{x^2/2}{1} = \frac{x^3}{2}$.
        $\frac{\tan x - \sin x}{x^3} \sim \frac{x^3/2}{x^3} = \boxed{\frac{1}{2}}$.
    \item Posons $h = x - \pi/2 \to 0$. $\cos x = \cos(\pi/2 + h) = -\sin h \sim -h$.
        $\frac{\cos x}{x - \pi/2} = \frac{-\sin h}{h} \to \boxed{-1}$.
\end{enumerate}
\end{corrige}

\coursec{Limites des fonctions irrationnelles : rationalisation et composition}

Dans la section précédente, nous avons appris à lever les indéterminations faisant intervenir les fonctions trigonométriques, grâce à la limite remarquable $\lim_{x\to 0}\dfrac{\sin x}{x}=1$ et aux théorèmes d'encadrement. Nous abordons maintenant une autre famille de fonctions très présente dans les métiers techniques : les \cle{fonctions irrationnelles}, c'est-à-dire celles qui contiennent la variable sous un radical, comme $f(x)=\sqrt{x^2+x}-x$. Ces fonctions apparaissent dès qu'on calcule une diagonale de tôle, une longueur de câble, une distance dans un plan d'atelier ou une impédance en électricité. Leur étude demande une vigilance particulière sur le \cle{domaine de définition} et une technique de calcul spécifique : la \cle{méthode de l'expression conjuguée}.

\courssub{Domaine de définition et continuité des fonctions irrationnelles}

\textbf{Intuition.} On ne peut pas prendre la racine carrée d'un nombre strictement négatif (dans $\mathbb{R}$). Avant tout calcul de limite, il faut donc se poser la question : \emph{pour quelles valeurs de $x$ l'expression sous la racine est-elle positive ou nulle ?}

\begin{definition}[Fonction irrationnelle et domaine de définition]
Une \cle{fonction irrationnelle} est une fonction dont l'expression contient la variable sous un radical. Pour une racine d'indice pair (racine carrée, racine quatrième...), le domaine de définition est l'ensemble des réels $x$ tels que le \cle{radicande} (l'expression sous la racine) soit positif ou nul :
\[ D_f=\{x\in\mathbb{R}\ ;\ P(x)\geq 0\} \quad \text{pour } f(x)=\sqrt{P(x)}. \]
\end{definition}

\begin{exemplebox}[Détermination d'un domaine]
Soit $f(x)=\sqrt{x^2+x}$. On résout $x^2+x\geq 0$, c'est-à-dire $x(x+1)\geq 0$. Le trinôme est positif à l'extérieur des racines $-1$ et $0$ :
\[ D_f=]-\infty;-1]\cup[0;+\infty[. \]
Toute question de limite pour cette fonction ne peut donc se poser qu'en $-\infty$, en $-1$, en $0$ ou en $+\infty$.
\end{exemplebox}

\begin{propriete}[Continuité des fonctions irrationnelles]
La fonction $x\mapsto\sqrt{x}$ est continue sur $[0;+\infty[$. Par composition avec une fonction polynôme (continue sur $\mathbb{R}$), toute fonction $x\mapsto\sqrt{P(x)}$ est \textbf{continue sur son domaine de définition}. En particulier, si $a\in D_f$, alors :
\[ \lim_{x\to a}\sqrt{P(x)}=\sqrt{P(a)}. \]
La limite se calcule alors par simple substitution.
\end{propriete}

\begin{exemplebox}[Limite par substitution]
$\displaystyle\lim_{x\to 3}\sqrt{x^2+x}=\sqrt{9+3}=\sqrt{12}=2\sqrt{3}$, car $3\in D_f$ et la fonction est continue en $3$.
\end{exemplebox}

\begin{attention}[Toujours vérifier le domaine d'abord]
Avant de calculer une limite, vérifiez que le point visé appartient bien au domaine (ou en est une borne). Chercher $\lim_{x\to 2}\sqrt{x^2+x}$ ne pose pas de problème, mais écrire $\lim_{x\to -0{,}5}\sqrt{x^2+x}$ n'a \textbf{aucun sens} : $-0{,}5\notin D_f$ car $(-0{,}5)^2+(-0{,}5)=-0{,}25<0$.
\end{attention}

\courssub{Limites aux bornes du domaine : comportement de $\sqrt{P(x)}$ à l'infini}

\textbf{Intuition.} Quand $x$ devient très grand, un polynôme se comporte comme son terme de plus haut degré. Sous une racine carrée, c'est encore vrai : $\sqrt{x^2+3x+1}$ « ressemble » à $\sqrt{x^2}=|x|$ quand $x$ est grand.

\begin{propriete}[Limite de $\sqrt{P(x)}$ en $+\infty$]
Soit $P(x)=a_nx^n+\dots+a_0$ un polynôme de degré $n$ avec $a_n>0$. Alors, en $+\infty$ :
\[ \sqrt{P(x)}\sim \sqrt{a_n}\;x^{n/2} \qquad\text{et}\qquad \lim_{x\to+\infty}\sqrt{P(x)}=+\infty. \]
Autrement dit, $\sqrt{P(x)}$ se comporte comme la racine du terme dominant. Attention au signe : en $-\infty$, $\sqrt{x^2}=|x|=-x$.
\end{propriete}

\begin{experience}[Démonstration guidée : factoriser le terme dominant]
On veut montrer que $\lim_{x\to+\infty}\sqrt{x^2+3x+1}=+\infty$ et préciser son comportement.
\begin{enumerate}
\item Pour $x>0$, factorisons $x^2$ sous la racine : $x^2+3x+1=x^2\left(1+\dfrac{3}{x}+\dfrac{1}{x^2}\right)$.
\item Comme $x>0$, $\sqrt{x^2}=x$, donc $\sqrt{x^2+3x+1}=x\sqrt{1+\dfrac{3}{x}+\dfrac{1}{x^2}}$.
\item Quand $x\to+\infty$, $\dfrac{3}{x}\to 0$ et $\dfrac{1}{x^2}\to 0$, donc la racine tend vers $\sqrt{1}=1$.
\item Par produit : $x\times(\text{quantité qui tend vers }1)\to+\infty$. \hfill $\square$
\end{enumerate}
Cette technique de \textbf{factorisation du terme dominant} est la première arme contre les indéterminations irrationnelles.
\end{experience}

\courssub{La forme indéterminée $\infty-\infty$ : la méthode de la conjuguée}

\textbf{Intuition.} Considérons $f(x)=\sqrt{x^2+x}-x$ en $+\infty$. D'un côté $\sqrt{x^2+x}\to+\infty$, de l'autre $x\to+\infty$ : c'est une forme indéterminée $\infty-\infty$. L'idée consiste à multiplier par l'\cle{expression conjuguée} pour faire disparaître la racine grâce à l'identité remarquable $(a-b)(a+b)=a^2-b^2$.

\begin{methode}[Lever une indétermination $\infty-\infty$ ou $\dfrac{0}{0}$ avec une racine]
\begin{enumerate}
\item Repérer la forme indéterminée (soustraction de deux quantités qui tendent vers l'infini, ou quotient $0/0$ avec une racine).
\item Multiplier le numérateur \textbf{et} le dénominateur par l'\cle{expression conjuguée} : on change le signe devant la racine. Le conjugué de $\sqrt{A}-B$ est $\sqrt{A}+B$.
\item Utiliser $(a-b)(a+b)=a^2-b^2$ au numérateur : la racine disparaît.
\item Simplifier, factoriser le terme dominant, puis conclure.
\end{enumerate}
\end{methode}

\begin{exoresolu}[Limite de $\sqrt{x^2+3x+1}-x$ en $+\infty$]
\textbf{Énoncé.} Calculer $\displaystyle\lim_{x\to+\infty}\left(\sqrt{x^2+3x+1}-x\right)$.
\tcblower
\textbf{Solution détaillée.} On a une forme $\infty-\infty$. On multiplie par le conjugué $\sqrt{x^2+3x+1}+x$ :
\begin{align*}
\sqrt{x^2+3x+1}-x &= \frac{\left(\sqrt{x^2+3x+1}-x\right)\left(\sqrt{x^2+3x+1}+x\right)}{\sqrt{x^2+3x+1}+x} \\
&= \frac{(x^2+3x+1)-x^2}{\sqrt{x^2+3x+1}+x} = \frac{3x+1}{\sqrt{x^2+3x+1}+x}.
\end{align*}
Pour $x>0$, on factorise $x$ en haut et en bas (avec $\sqrt{x^2}=x$) :
\[ \frac{3x+1}{\sqrt{x^2+3x+1}+x}=\frac{x\left(3+\frac{1}{x}\right)}{x\left(\sqrt{1+\frac{3}{x}+\frac{1}{x^2}}+1\right)}=\frac{3+\frac{1}{x}}{\sqrt{1+\frac{3}{x}+\frac{1}{x^2}}+1}. \]
Quand $x\to+\infty$, le numérateur tend vers $3$ et le dénominateur vers $\sqrt{1}+1=2$. Donc :
\[ \lim_{x\to+\infty}\left(\sqrt{x^2+3x+1}-x\right)=\frac{3}{2}. \]
\textbf{Interprétation graphique :} la droite d'équation $y=\dfrac{3}{2}$ est asymptote horizontale à la courbe en $+\infty$.
\end{exoresolu}

\begin{popfigure}
\begin{center}
\begin{tikzpicture}
\begin{axis}[width=0.85\linewidth,height=6cm,axis lines=middle,xlabel={$x$},ylabel={$y$},xmin=-0.3,xmax=10,ymin=-0.2,ymax=1.2,samples=200,grid=both,grid style={popline!40},legend style={at={(0.97,0.35)},anchor=east,font=\small}]
\addplot[popcurve,domain=0:10]{sqrt(max(0,x^2+x))-x};
\addlegendentry{$y=\sqrt{x^2+x}-x$}
\addplot[poporange,thick,dashed,domain=-0.3:10]{0.5};
\addlegendentry{asymptote $y=\dfrac{1}{2}$}
\end{axis}
\end{tikzpicture}
\end{center}
\legende{La courbe de $f(x)=\sqrt{x^2+x}-x$ se rapproche de son asymptote horizontale $y=\dfrac{1}{2}$ quand $x\to+\infty$ (même calcul que l'exercice résolu avec $3x$ remplacé par $x$ : la limite vaut $\frac{1}{2}$).}
\end{popfigure}

\courssub{La forme indéterminée $\dfrac{0}{0}$ : conjuguée au numérateur}

La même technique fonctionne quand numérateur et dénominateur tendent tous les deux vers $0$. Exemple classique : $\lim_{x\to 0}\dfrac{\sqrt{1+x}-1}{x}$.

\begin{align*}
\frac{\sqrt{1+x}-1}{x} &= \frac{(\sqrt{1+x}-1)(\sqrt{1+x}+1)}{x(\sqrt{1+x}+1)} = \frac{(1+x)-1}{x(\sqrt{1+x}+1)} = \frac{x}{x(\sqrt{1+x}+1)} = \frac{1}{\sqrt{1+x}+1}.
\end{align*}
En simplifiant par $x$ (licite car $x\neq 0$ au voisinage de $0$), on obtient une expression continue en $0$ :
\[ \lim_{x\to 0}\frac{\sqrt{1+x}-1}{x}=\frac{1}{\sqrt{1}+1}=\frac{1}{2}. \]

\begin{exoresolu}[Une forme $0/0$ d'ordre supérieur]
\textbf{Énoncé.} Calculer $\displaystyle\lim_{x\to 0}\frac{\sqrt{1+2x}-1-x}{x^2}$.
\tcblower
\textbf{Solution détaillée.} En $0$ : $\sqrt{1+0}-1-0=0$ et $x^2\to 0$ : forme $\dfrac{0}{0}$. Multiplions par le conjugué de $\sqrt{1+2x}-(1+x)$, c'est-à-dire $\sqrt{1+2x}+(1+x)$ :
\begin{align*}
\frac{\sqrt{1+2x}-(1+x)}{x^2} &= \frac{(1+2x)-(1+x)^2}{x^2\left(\sqrt{1+2x}+1+x\right)} = \frac{1+2x-1-2x-x^2}{x^2\left(\sqrt{1+2x}+1+x\right)} \\
&= \frac{-x^2}{x^2\left(\sqrt{1+2x}+1+x\right)} = \frac{-1}{\sqrt{1+2x}+1+x}.
\end{align*}
Le dénominateur tend vers $\sqrt{1}+1+0=2$, donc :
\[ \lim_{x\to 0}\frac{\sqrt{1+2x}-1-x}{x^2}=-\frac{1}{2}. \]
\end{exoresolu}

\begin{popfigure}
\begin{center}
\begin{tikzpicture}
\begin{axis}[width=0.85\linewidth,height=6cm,axis lines=middle,xlabel={$x$},ylabel={$y$},xmin=-1,xmax=3,ymin=0.2,ymax=0.9,samples=200,grid=both,grid style={popline!40}]
\addplot[popcurve,domain=-0.99:-0.02]{(sqrt(max(0,1+x))-1)/x};
\addplot[popcurve,domain=0.02:3]{(sqrt(max(0,1+x))-1)/x};
\addplot[poporange,thick,dashed,domain=-1:3]{0.5};
\addplot[popdark,only marks,mark=o,mark size=2pt,thick] coordinates {(0,0.5)};
\end{axis}
\end{tikzpicture}
\end{center}
\legende{Courbe de $g(x)=\dfrac{\sqrt{1+x}-1}{x}$ : la fonction n'est pas définie en $0$ (trou, symbolisé par le rond vide), mais la limite en $0$ existe et vaut $\dfrac{1}{2}$ (droite en pointillés).}
\end{popfigure}

\begin{attention}[Erreurs fréquentes avec les racines carrées]
\begin{itemize}
\item Écrire $\sqrt{x^2}=x$ : c'est \textbf{faux} pour $x<0$ ! On a toujours $\sqrt{x^2}=|x|$. En $-\infty$, $\sqrt{x^2}=-x$.
\item Oublier le signe moins en développant le conjugué : $(\sqrt{A}-B)(\sqrt{A}+B)=A-B^2$, pas $A+B^2$.
\item Simplifier par $x$ sans préciser que $x\neq 0$ (légitime car on cherche une limite, mais il faut le dire).
\item Chercher la limite en un point hors du domaine de définition.
\end{itemize}
\end{attention}

\courssub{Limites par changement de variable}

\begin{methode}[Changement de variable]
Pour calculer $\lim_{x\to 0^+}f(\sqrt{x})$, on pose $t=\sqrt{x}$ (c'est-à-dire $x=t^2$). Quand $x\to 0^+$, alors $t\to 0^+$, et on est ramené à $\lim_{t\to 0^+}f(t)$, souvent plus simple. De même, en $+\infty$, poser $t=\sqrt{x}$ transforme $\sqrt{x}$ en $t$ et $x$ en $t^2$.
\end{methode}

\begin{exemplebox}[Exemple chiffré]
Calculer $\displaystyle\lim_{x\to 0^+}\frac{\sqrt{x}}{x+2\sqrt{x}}$. Posons $t=\sqrt{x}$, donc $x=t^2$ et $t\to 0^+$ :
\[ \frac{t}{t^2+2t}=\frac{t}{t(t+2)}=\frac{1}{t+2}\xrightarrow[t\to 0^+]{}\frac{1}{2}. \]
\end{exemplebox}

\courssub{Comparaison de croissance et asymptotes}

\begin{propriete}[Croissances comparées en $+\infty$]
Pour $x\to+\infty$, les puissances de $x$ se classent ainsi : si $a<\dfrac{1}{2}$, alors
\[ x^a \ll \sqrt{x} \ll x \ll x^2, \]
c'est-à-dire $\lim_{x\to+\infty}\dfrac{\sqrt{x}}{x^a}=+\infty$, $\lim_{x\to+\infty}\dfrac{x}{\sqrt{x}}=+\infty$, $\lim_{x\to+\infty}\dfrac{x^2}{x}=+\infty$. Une puissance plus grande « écrase » toujours une puissance plus petite.
\end{propriete}

\textbf{Utilité.} Ces comparaisons justifient le comportement asymptotique : puisque $\sqrt{x^2+x}\sim x$ en $+\infty$, la différence $\sqrt{x^2+x}-x$ est un « reste » petit devant $x$, qui tend vers une constante (ici $\frac{1}{2}$) : d'où l'asymptote horizontale. De même, $\dfrac{\sqrt{x}}{x}=\dfrac{1}{\sqrt{x}}\to 0$ : toute courbe de $\dfrac{\sqrt{x}}{x}$ admet l'axe des abscisses comme asymptote en $+\infty$.

\begin{popfigure}
\begin{center}
\begin{tikzpicture}[font=\small,node distance=6mm,
decision/.style={diamond,draw=popdark,fill=poppurpleL,aspect=2.2,inner sep=1pt,align=center},
action/.style={rectangle,rounded corners,draw=popdark,fill=popblueL,align=center,inner sep=4pt},
result/.style={rectangle,rounded corners,draw=popdark,fill=popgreenL,align=center,inner sep=4pt},
fleche/.style={-{Stealth},thick,popdark}]
\node[decision, align=center] (fi){Forme\\indéterminée ?};
\node[action,below left=8mm and 24mm of fi, align=center] (subst){Non : substitution\\(continuité)};
\node[decision,below=8mm of fi, align=center] (type){Type $\infty-\infty$\\ou $\frac{0}{0}$ ?};
\node[action,below left=8mm and 20mm of type, align=center] (conj){Multiplier par\\l'expression conjuguée};
\node[action,below right=8mm and 20mm of type, align=center] (fact){Factoriser le\\terme dominant};
\node[action,below=8mm of conj, align=center] (cv){Si besoin : changement\\de variable $t=\sqrt{x}$};
\node[result,below=8mm of cv, align=center] (fin){Conclure et vérifier\\le domaine};
\draw[fleche] (fi) -- node[left,font=\scriptsize]{non} (subst);
\draw[fleche] (fi) -- node[right,font=\scriptsize]{oui} (type);
\draw[fleche] (type) -- node[left,font=\scriptsize]{racine} (conj);
\draw[fleche] (type) -- node[right,font=\scriptsize]{polynômes} (fact);
\draw[fleche] (conj) -- (cv);
\draw[fleche] (cv) -- (fin);
\draw[fleche] (fact) |- (fin.east);
\end{tikzpicture}
\end{center}
\legende{Arbre de décision pour le calcul d'une limite de fonction irrationnelle.}
\end{popfigure}

\begin{savaistu}[Et pour les racines cubiques ?]
La méthode de la conjuguée se généralise : pour une racine cubique, on utilise l'identité $a^3-b^3=(a-b)(a^2+ab+b^2)$. Par exemple, pour $\sqrt[3]{x+1}-1$, on multiplierait par $\sqrt[3]{(x+1)^2}+\sqrt[3]{x+1}+1$. Cette technique n'est pas au programme de Terminale technique, mais elle montre la puissance des identités remarquables : un même outil algébrique résout toute une famille de problèmes de limites !
\end{savaistu}

\begin{exercice}[À faire]
\begin{enumerate}
\item Déterminer le domaine de définition de $f(x)=\sqrt{x^2-4}$, puis calculer $\lim_{x\to+\infty}\left(\sqrt{x^2-4}-x\right)$.
\item Calculer $\displaystyle\lim_{x\to 0}\frac{\sqrt{4+x}-2}{x}$.
\item En posant $t=\sqrt{x}$, calculer $\displaystyle\lim_{x\to 0^+}\frac{x-\sqrt{x}}{x+\sqrt{x}}$.
\end{enumerate}
\end{exercice}

\begin{corrige}[Exercice]
\begin{enumerate}
\item $x^2-4\geq 0\iff x\leq -2$ ou $x\geq 2$, donc $D_f=]-\infty;-2]\cup[2;+\infty[$. Par la conjuguée : $\sqrt{x^2-4}-x=\dfrac{(x^2-4)-x^2}{\sqrt{x^2-4}+x}=\dfrac{-4}{\sqrt{x^2-4}+x}\to 0$ car le dénominateur tend vers $+\infty$. La droite $y=0$ est asymptote en $+\infty$.
\item Conjuguée : $\dfrac{(\sqrt{4+x}-2)(\sqrt{4+x}+2)}{x(\sqrt{4+x}+2)}=\dfrac{(4+x)-4}{x(\sqrt{4+x}+2)}=\dfrac{x}{x(\sqrt{4+x}+2)}=\dfrac{1}{\sqrt{4+x}+2}\to\dfrac{1}{4}$.
\item Avec $x=t^2$, $t\to 0^+$ : $\dfrac{t^2-t}{t^2+t}=\dfrac{t(t-1)}{t(t+1)}=\dfrac{t-1}{t+1}\to\dfrac{-1}{1}=-1$.
\end{enumerate}
\end{corrige}

\begin{aretenir}[L'essentiel de la section]
\begin{itemize}
\item Domaine de $\sqrt{P(x)}$ : $P(x)\geq 0$. Toute fonction irrationnelle est continue sur son domaine : la limite en un point intérieur se fait par substitution.
\item En $+\infty$, $\sqrt{P(x)}\sim\sqrt{a_n}\,x^{n/2}$ (racine du terme dominant).
\item Forme $\infty-\infty$ ou $\dfrac{0}{0}$ avec une racine : multiplier numérateur et dénominateur par l'\cle{expression conjuguée}, puis utiliser $(a-b)(a+b)=a^2-b^2$.
\item Changement de variable $t=\sqrt{x}$ pour se ramener à une limite connue.
\item Vigilance permanente : $\sqrt{x^2}=|x|$, et toujours vérifier le domaine avant de conclure.
\end{itemize}
\end{aretenir}

\coursec{Représentation graphique : méthodologie complète et lecture graphique}

Après avoir maîtrisé le calcul des limites, notamment pour les fonctions irrationnelles où la technique de rationalisation s'avère indispensable, nous disposons désormais de tous les outils analytiques pour dresser le portrait complet d'une fonction. La représentation graphique n'est pas une simple formalité de dessin : c'est la synthèse visuelle de l'étude algébrique, un langage universel pour communiquer le comportement d'un phénomène technique (charge d'une poutre, réponse d'un filtre, évolution d'une population). Au Probatoire, un tracé soigné, respectant l'échelle, les asymptotes en pointillés et les points caractéristiques, est un gage de rigueur qui valorise votre copie.

\courssub{Algorithme complet d'étude d'une fonction}

Toute étude de fonction suit une logique immuable, du global vers le local. Respecter cet ordre évite les calculs inutiles et structure la rédaction.

\begin{methode}[Étude complète d'une fonction $f$]
\begin{enumerate}
    \item \textbf{Domaine de définition $D_f$ :} Identifier les restrictions (dénominateur non nul, radicande $\ge 0$ pour racine paire, argument $>0$ pour $\ln$). L'écrire sous forme d'intervalle(s) ou d'ensemble.
    \item \textbf{Parité et Périodicité (Symétries) :}
    \begin{itemize}
        \item Paire : $f(-x)=f(x)$ $\Rightarrow$ symétrie par rapport à l'axe $(Oy)$ (étudier sur $[0,+\infty[$).
        \item Impaire : $f(-x)=-f(x)$ $\Rightarrow$ symétrie par rapport à l'origine $O$ (étudier sur $[0,+\infty[$).
        \item Périodique de période $T$ : $f(x+T)=f(x)$ $\Rightarrow$ étudier sur une période (ex: $[0, 2\pi[$ pour $\sin, \cos$).
    \end{itemize}
    \item \textbf{Limites aux bornes de $D_f$ et Asymptotes :}
    \begin{itemize}
        \item Calculer $\lim_{x \to \text{bornes}} f(x)$ (finies ou infinies).
        \item Rechercher les asymptotes (verticales, horizontales, obliques) -- voir \S 5.2.
    \end{itemize}
    \item \textbf{Dérivée $f'$, signe et Tableau de variations :}
    \begin{itemize}
        \item Calculer $f'(x)$ sur $D_f$ (règles de dérivation, dérivation des fonctions composées $u(v(x))$).
        \item Résoudre $f'(x)=0$ et étudier le signe de $f'(x)$ (tableau de signes).
        \item Construire le \textbf{tableau de variations} (flèches $\nearrow$ si $f'>0$, $\searrow$ si $f'<0$).
    \end{itemize}
    \item \textbf{Valeurs remarquables :} Intersections avec les axes ($f(0)$, solutions de $f(x)=0$), extrema (valeurs aux bornes des variations), points d'inflexion éventuels.
    \item \textbf{Tracé de la courbe $\mathcal{C}_f$ :} Voir la méthode encadrée ci-dessous.
\end{enumerate}
\end{methode}

\begin{methode}[Tracer la courbe représentative $\mathcal{C}_f$]
\begin{enumerate}
    \item \textbf{Repère orthonormé :} Choisir des unités adaptées (souvent différentes sur $x$ et $y$). Graduations régulières. Nommer les axes.
    \item \textbf{Asymptotes :} Tracer les asymptotes verticales (pointillés verticaux) et obliques/horizontales (pointillés obliques/horizontaux). Les étiqueter ($y=ax+b$).
    \item \textbf{Points caractéristiques :} Placer les points d'intersection avec les axes, les extrema, les points d'inflexion. Marquer les trous (cercles vides $\circ$) si discontinuité évitable (ex: $x \mapsto \frac{x^2-1}{x-1}$ en $x=1$).
    \item \textbf{Branches :} Tracer les branches de courbe en respectant :
    \begin{itemize}
        \item Le sens de variation (croissant/décroissant).
        \item La convexité apparente (courbure vers le haut/bas).
        \item L'approche des asymptotes (la courbe s'en rapproche sans la couper pour une verticale ; elle peut couper une horizontale/oblique).
    \end{itemize}
    \item \textbf{Finition :} Flèches aux extrémités des branches (sauf si domaine borné). Légende $\mathcal{C}_f$.
\end{enumerate}
\end{methode}

\begin{attention}[Pièges classiques au tracé]
\begin{itemize}
    \item \textbf{Couper une asymptote verticale :} \textbf{Interdit.} Si $\lim_{x \to x_0} f(x) = \pm\infty$, la courbe s'éloigne indéfiniment de la droite $x=x_0$. Elle ne la traverse jamais.
    \item \textbf{Oublier le domaine $D_f$ :} Ne tracer la courbe \textbf{que} sur $D_f$. Pour $f(x)=\sqrt{x}$, rien à gauche de $0$.
    \item \textbf{Confondre tangente verticale et asymptote verticale :}
    \begin{itemize}
        \item Asymptote verticale en $x_0$ : $\lim_{x \to x_0} f(x) = \pm\infty$ (la fonction « explose »).
        \item Tangente verticale en $x_0 \in D_f$ : $f$ continue en $x_0$, mais $\lim_{x \to x_0} f'(x) = \pm\infty$ (la pente devient infinie, ex: $f(x)=\sqrt[3]{x}$ en $0$ ou $f(x)=\sqrt{x}$ en $0^+$). La courbe a une tangente verticale mais reste finie.
    \end{itemize}
    \item \textbf{Échelle non respectée :} Si $1$ cm $= 1$ sur $Ox$ et $1$ cm $= 10$ sur $Oy$, l'indiquer clairement. Une courbe « aplatie » ou « étirée » fausse la lecture graphique.
\end{itemize}
\end{attention}

\courssub{Recherche des asymptotes : méthode systématique}

Une asymptote est une droite que la courbe approche indéfiniment. Trois cas existent.

\begin{propriete}[Asymptotes -- Définition et recherche]
Soit $f$ définie sur un voisinage de $+\infty$ (ou $-\infty$) ou de $x_0 \notin D_f$.
\begin{enumerate}
    \item \textbf{Asymptote verticale (AV) :} Droite $x = x_0$ ($x_0$ fini).
    $\mathcal{C}_f$ admet $x=x_0$ comme AV $\iff \lim_{x \to x_0^{\pm}} f(x) = \pm\infty$.
    \item \textbf{Asymptote horizontale (AH) :} Droite $y = \ell$ ($\ell$ fini).
    $\mathcal{C}_f$ admet $y=\ell$ comme AH en $+\infty$ $\iff \lim_{x \to +\infty} f(x) = \ell$.
    \item \textbf{Asymptote oblique (AO) :} Droite $y = ax + b$ ($a \neq 0$ fini).
    $\mathcal{C}_f$ admet $y=ax+b$ comme AO en $+\infty$ $\iff$
    \begin{align*}
        a &= \lim_{x \to +\infty} \frac{f(x)}{x} \quad (a \in \mathbb{R}^*) \\
        b &= \lim_{x \to +\infty} \bigl( f(x) - ax \bigr) \quad (b \in \mathbb{R})
    \end{align*}
    \textbf{Remarques :} Si $a=0$, on retombe sur l'horizontale. Si $a$ est infini, ce n'est pas une oblique (branche parabolique, hors programme). La même procédure s'applique en $-\infty$ (les coefficients $a, b$ peuvent différer).
\end{enumerate}
\end{propriete}

\begin{exoresolu}[Étude et tracé de $f(x) = \dfrac{x^2 - 1}{\sqrt{x^2 + 1}}$ sur $\mathbb{R}$]
\textbf{1. Domaine :} $x^2+1 > 0$ pour tout $x \in \mathbb{R}$. $D_f = \mathbb{R}$.
\textbf{2. Parité :} $f(-x) = \frac{(-x)^2-1}{\sqrt{(-x)^2+1}} = \frac{x^2-1}{\sqrt{x^2+1}} = f(x)$. \textbf{Fonction paire.} Symétrie par rapport à $(Oy)$. On étudie sur $[0, +\infty[$.
\textbf{3. Limites et Asymptotes :}
\begin{itemize}
    \item $x \to +\infty$ : $f(x) \sim \frac{x^2}{\sqrt{x^2}} = \frac{x^2}{|x|} = x$ (car $x>0$).
    \item Asymptote oblique en $+\infty$ :
    $a = \lim_{+\infty} \frac{f(x)}{x} = \lim \frac{x^2-1}{x\sqrt{x^2+1}} = \lim \frac{x^2}{x\cdot x} = 1$.
    $b = \lim_{+\infty} \bigl( f(x) - x \bigr) = \lim \frac{x^2-1 - x\sqrt{x^2+1}}{\sqrt{x^2+1}}$.
    Rationalisons le numérateur :
    $(x^2-1 - x\sqrt{x^2+1})(x^2-1 + x\sqrt{x^2+1}) = (x^2-1)^2 - x^2(x^2+1) = x^4 -2x^2+1 - x^4 - x^2 = -3x^2+1$.
    Donc $b = \lim_{+\infty} \frac{-3x^2+1}{\sqrt{x^2+1}(x^2-1 + x\sqrt{x^2+1})} \sim \frac{-3x^2}{x(x^2+x^2)} = \frac{-3x^2}{2x^3} = 0$.
    \textbf{AO en $+\infty$ : $y = x$.}
    \item Par parité, en $-\infty$ : $f(x) \sim -x$ (car $\sqrt{x^2}=|x|=-x$). \textbf{AO en $-\infty$ : $y = -x$.}
    \item Pas d'asymptote verticale (domaine $\mathbb{R}$).
\end{itemize}
\textbf{4. Dérivée et Variations (sur $[0,+\infty[$) :}
$f(x) = \frac{x^2-1}{(x^2+1)^{1/2}}$.
$f'(x) = \frac{2x(x^2+1)^{1/2} - (x^2-1)\cdot\frac{1}{2}(x^2+1)^{-1/2}\cdot 2x}{x^2+1}$
$= \frac{2x(x^2+1) - x(x^2-1)}{(x^2+1)^{3/2}} = \frac{2x^3+2x - x^3+x}{(x^2+1)^{3/2}} = \frac{x^3+3x}{(x^2+1)^{3/2}} = \frac{x(x^2+3)}{(x^2+1)^{3/2}}$.
Sur $[0,+\infty[$, dénominateur $>0$, $x^2+3>0$. Signe de $f'$ = signe de $x$.
$f'(x) = 0 \iff x=0$. $f'(x) > 0$ sur $]0,+\infty[$.
\textbf{Tableau de variations sur $[0,+\infty[$ :}
\begin{center}
\begin{tabular}{|c|ccc|}\hline
$x$ & $0$ & & $+\infty$ \\ \hline
$f'(x)$ & $0$ & $+$ & \\ \hline
$f$ & $-1$ & $\nearrow$ & $+\infty$ \\ \hline
\end{tabular}
\end{center}
Par parité, sur $]-\infty,0]$ : $f$ décroît de $+\infty$ à $-1$.
\textbf{5. Points remarquables :}
$f(0) = -1$ (minimum global, intersection $Oy$).
$f(x)=0 \iff x^2-1=0 \iff x=\pm 1$ (intersections $Ox$).
\textbf{6. Tracé :} Voir figure ci-dessous.
\tcblower
\textbf{Solution détaillée :} La courbe passe par $A(-1,0)$, $B(1,0)$, $C(0,-1)$ (minimum). Elle est tangente aux asymptotes $y=x$ (en $+\infty$) et $y=-x$ (en $-\infty$) « par le dessous » car $f(x) - x \to 0^-$ (numérateur $-3x^2+1 < 0$ pour $x$ grand). La convexité change (point d'inflexion vers l'origine), mais non demandé ici.
\end{exoresolu}

\begin{popfigure}
\begin{center}
\begin{tikzpicture}[scale=0.85]
\begin{axis}[
    axis lines=middle,
    xlabel=$x$, ylabel=$y$,
    xmin=-5, xmax=5,
    ymin=-5, ymax=5,
    xtick={-4,-3,-2,-1,0,1,2,3,4},
    ytick={-4,-3,-2,-1,0,1,2,3,4},
    tick label style={font=\scriptsize},
    axis line style={-Stealth},
    clip=false,
    domain=-5:5,
    samples=400,
    restrict y to domain=-5:5,
    width=0.9\linewidth,
    height=0.6\linewidth,
]
% Asymptotes
\addplot[dashed, popmuted, thick, domain=-5:5] {x};
\addplot[dashed, popmuted, thick, domain=-5:5] {-x};
\node[popmuted, anchor=south west, font=\scriptsize] at (axis cs:4,4) {$y=x$};
\node[popmuted, anchor=north west, font=\scriptsize] at (axis cs:-4,4) {$y=-x$};

% Function f(x) = (x^2-1)/sqrt(x^2+1)
\addplot[popcurve, very thick, domain=-5:5] {(x^2-1)/sqrt(max(0,x^2+1))};

% Remarkable points
\addplot[only marks, mark=*, mark size=2.5, popblue] coordinates {(0,-1) (1,0) (-1,0)};
\node[popblue, anchor=north, font=\scriptsize] at (axis cs:0,-1.3) {$(0,-1)$ min};
\node[popblue, anchor=south, font=\scriptsize] at (axis cs:1,0.3) {$(1,0)$};
\node[popblue, anchor=south, font=\scriptsize] at (axis cs:-1,0.3) {$(-1,0)$};

% Variation table sketch (manual placement)
\node[anchor=north west, font=\tiny, fill=popcream, draw=popline, rounded corners, inner sep=2pt, align=center] at (rel axis cs:0.02,0.98){
\begin{tabular}{|c|ccccc|}\hline
$x$ & $-\infty$ & & $0$ & & $+\infty$ \\ \hline
$f'$ & & $-$ & $0$ & $+$ & \\ \hline
$f$ & $+\infty$ & $\searrow$ & $-1$ & $\nearrow$ & $+\infty$ \\ \hline
\end{tabular}
};
\end{axis}
\end{tikzpicture}
\end{center}
\legende{Courbe de $f(x)=\frac{x^2-1}{\sqrt{x^2+1}}$ : asymptotes obliques $y=\pm x$ (pointillés), minimum global en $(0,-1)$, zéros en $\pm 1$, tableau de variations intégré.}
\end{popfigure}

\courssub{Cas particuliers : fonctions trigonométriques et irrationnelles}

\begin{experience}[Fonctions trigonométriques : réduire l'étude à une période]
Pour $f(x) = 2\sin x + \cos 2x$ sur $\mathbb{R}$.
\begin{enumerate}
    \item $f$ est $2\pi$-périodique (PPCM des périodes $2\pi$ et $\pi$). Étudier sur $[0, 2\pi[$.
    \item Dérivée : $f'(x) = 2\cos x - 2\sin 2x = 2\cos x - 4\sin x \cos x = 2\cos x(1 - 2\sin x)$.
    \item Signe sur $[0, 2\pi[$ : Zéros de $\cos x$ : $\pi/2, 3\pi/2$. Zéros de $1-2\sin x$ : $\sin x = 1/2 \Rightarrow x=\pi/6, 5\pi/6$.
    \item Tableau de variations sur $[0, 2\pi[$, puis reporter par translations de $2\pi$.
\end{enumerate}
\textbf{Astuce :} Au Probatoire, précisez « $f$ est $2\pi$-périodique, on l'étudie sur $[0, 2\pi[$ » et tracez une période, puis indiquez la répétition par flèches.
\end{experience}

\begin{experience}[Fonctions irrationnelles : domaine borné et tangente verticale]
Pour $f(x) = \sqrt{x(4-x)}$ sur $[0, 4]$.
\begin{enumerate}
    \item Domaine : $x(4-x) \ge 0 \Rightarrow x \in [0, 4]$.
    \item Limites aux bornes : $\lim_{x \to 0^+} f(x) = 0$, $\lim_{x \to 4^-} f(x) = 0$.
    \item Dérivée : $f'(x) = \frac{4-2x}{2\sqrt{x(4-x)}} = \frac{2-x}{\sqrt{x(4-x)}}$.
    \item Variations : $f' > 0$ sur $[0, 2[$, $f'(2)=0$, $f' < 0$ sur $]2, 4]$. Max en $x=2$, $f(2)=2$.
    \item \textbf{Tangentes verticales :} $\lim_{x \to 0^+} f'(x) = \frac{2}{0^+} = +\infty$. $\lim_{x \to 4^-} f'(x) = \frac{-2}{0^+} = -\infty$.
    La courbe part de $(0,0)$ avec une tangente verticale (axe $Oy$), monte jusqu'au sommet $(2,2)$, puis redescend vers $(4,0)$ avec tangente verticale (droite $x=4$).
    \textbf{Attention :} Pas d'asymptote verticale car $f(0)$ et $f(4)$ existent et sont finis ($=0$).
\end{enumerate}
\end{experience}

\courssub{Lecture graphique : résolution d'équations et d'inéquations}

La courbe $\mathcal{C}_f$ est l'outil visuel pour résoudre $f(x) = k$ ou $f(x) > k$ sans calcul algébrique lourd.

\begin{propriete}[Lecture graphique]
Soit $f$ une fonction définie sur $D$, $\mathcal{C}_f$ sa courbe, $k \in \mathbb{R}$.
\begin{itemize}
    \item \textbf{Équation $f(x) = k$ :} Les solutions sont les \textbf{abscisses} des points d'intersection de $\mathcal{C}_f$ avec la droite horizontale $y=k$.
    \item \textbf{Inéquation $f(x) > k$ (ou $\ge, <, \le$) :} L'ensemble des solutions est l'ensemble des abscisses $x$ pour lesquels le point de $\mathcal{C}_f$ se situe \textbf{strictement au-dessus} (ou sur, ou en-dessous) de la droite $y=k$.
    \item \textbf{Nombre de solutions :} Il est égal au nombre de points d'intersection (en comptant les tangences comme une solution double si nécessaire, mais visuellement un point de contact).
\end{itemize}
\end{propriete}

\begin{popfigure}
\begin{center}
\begin{tikzpicture}[scale=0.9]
\begin{axis}[
    axis lines=middle,
    xlabel=$x$, ylabel=$y$,
    xmin=-3.5, xmax=3.5,
    ymin=-1.5, ymax=3.5,
    xtick={-3,-2,-1,0,1,2,3},
    ytick={-1,0,1,2,3},
    tick label style={font=\scriptsize},
    axis line style={-Stealth},
    clip=false,
    domain=-3:3,
    samples=300,
    width=0.9\linewidth,
    height=0.6\linewidth,
]
% Function: f(x) = x^3 - 3x + 1
\addplot[popcurve, very thick, domain=-2.5:2.5] {x^3 - 3*x + 1};
\node[popblue, anchor=south west, font=\footnotesize] at (axis cs:2, 2.5) {$\mathcal{C}_f : y = x^3 - 3x + 1$};

% Horizontal line y = k (k=1)
\addplot[dashed, poporange, thick, domain=-3:3] {1};
\node[poporange, anchor=south west, font=\scriptsize] at (axis cs:2.2, 1.1) {$y = k$};

% Intersection points (x^3-3x+1=1 -> x=0, +-sqrt(3))
\addplot[only marks, mark=*, mark size=3, poporange] coordinates {(-1.732,1) (0,1) (1.732,1)};
\node[poporange, anchor=north, font=\scriptsize] at (axis cs:-1.732, 0.7) {$x_1$};
\node[poporange, anchor=north, font=\scriptsize] at (axis cs:0, 0.7) {$x_2$};
\node[poporange, anchor=north, font=\scriptsize] at (axis cs:1.732, 0.7) {$x_3$};

% Zone labels for f(x) > k and f(x) < k
\node[poporange!80!black, font=\scriptsize, fill=popcream, rounded corners, inner sep=2pt, opacity=0.9] at (axis cs:-2.2, 2.5) {$f(x) < k$};
\node[poporange!80!black, font=\scriptsize, fill=popcream, rounded corners, inner sep=2pt, opacity=0.9] at (axis cs:-0.8, -0.5) {$f(x) > k$};
\node[poporange!80!black, font=\scriptsize, fill=popcream, rounded corners, inner sep=2pt, opacity=0.9] at (axis cs:0.8, 2.5) {$f(x) < k$};
\node[poporange!80!black, font=\scriptsize, fill=popcream, rounded corners, inner sep=2pt, opacity=0.9] at (axis cs:2.2, -0.5) {$f(x) > k$};

% Legend
\node[poporange!80!black, anchor=north east, font=\scriptsize, fill=popcream, rounded corners, inner sep=2pt] at (rel axis cs:0.98, 0.98) {
    Zones : $\mathcal{C}_f$ au-dessus de $y=k \iff f(x)>k$
};
\end{axis}
\end{tikzpicture}
\end{center}
\legende{Résolution graphique de $f(x)=k$ (3 solutions $x_1, x_2, x_3$) et $f(x)>k$ (zones où la courbe est au-dessus de $y=k$).}
\end{popfigure}

\courssub{Convexité et utilisation des outils numériques}

\begin{aretenir}[Convexité (notion admise en Terminale Tech)]
On admet qu'une fonction deux fois dérivable est :
\begin{itemize}
    \item \textbf{Convexe} sur un intervalle $I$ si sa courbe se situe \textbf{au-dessus} de ses tangentes (ou $\mathcal{C}_f$ « tient l'eau »). $\iff f''(x) \ge 0$ sur $I$.
    \item \textbf{Concave} sur $I$ si sa courbe se situe \textbf{en-dessous} de ses tangentes. $\iff f''(x) \le 0$ sur $I$.
\end{itemize}
Un \textbf{point d'inflexion} correspond à un changement de convexité ($f''$ change de signe). Au Probatoire, on ne demande généralement pas le calcul formel de $f''$, mais on attend que le tracé respecte la convexité visible (courbure naturelle). Si vous calculez $f''$ pour préciser le tracé, c'est valorisé.
\end{aretenir}

\begin{methode}[Utilisation de la calculatrice / GeoGebra / Python (Avant l'épreuve)]
\begin{enumerate}
    \item \textbf{Conjecturer :} Tracer $f$ sur l'outil. Observer domaine, asymptotes, variations, extrema, convexité.
    \item \textbf{Vérifier :} Lire les coordonnées des points remarquables (valeurs exactes si possible, sinon approchées).
    \item \textbf{Valider :} Confronter vos calculs analytiques (limites, dérivée, variations) au graphique. Si divergence $\rightarrow$ chercher l'erreur de calcul.
    \item \textbf{Interdit en examen :} Recopier le graphique de la calculatrice sans l'étude analytique justifiée. L'outil est un \textbf{aide à la rédaction}, pas une preuve.
\end{enumerate}
\end{methode}

\begin{savaistu}[Le tracé, un art noté au Probatoire]
En épreuve du Probatoire / Baccalauréat Technique, la \textbf{qualité graphique} est explicitement notée (souvent 2 à 4 points sur 20\%). Les correcteurs vérifient :
\begin{itemize}
    \item Repère orthonormé, unités lisibles, axes nommés.
    \item Asymptotes en \textbf{pointillés} (pas traits pleins !), étiquetées.
    \item Points caractéristiques (intersections, extrema) \textbf{placés et nommés}.
    \item Courbe tracée \textbf{uniquement sur $D_f$} (arrêt net aux bornes, cercles vides si trou).
    \item Respect des variations et de la convexité apparente.
    \item Flèches aux bouts des branches non bornées.
\end{itemize}
Un candidat qui a une étude analytique imparfaite mais un tracé propre et cohérent avec \emph{ses} résultats récupère des points précieux. Inversement, une étude parfaite gâchée par un tracé bâclé (asymptote coupée, pas d'échelle, courbe sur tout $\mathbb{R}$ pour $\sqrt{x}$) perd la note de présentation.
\end{savaistu}

\courssub{Synthèse visuelle : Fiche méthode « Étude complète »}

\begin{popfigure}
\begin{center}
\begin{tikzpicture}[
    mindmap,
    grow cyclic,
    every node/.style={concept, execute at begin node=\small\bfseries},
    concept color=popblue,
    level 1/.append style={level distance=4.5cm, sibling angle=60, concept color=popblue!70},
    level 2/.append style={level distance=3.2cm, sibling angle=45, concept color=poporange!80, font=\tiny},
    level 3/.append style={level distance=2.8cm, sibling angle=40, concept color=popgreen!70, font=\tiny},
    text=popdark
]
\node[root concept, align=center]{ÉTUDE COMPLÈTE\\ DE $f$}
    child[concept color=poppurple] { node {1. DOMAINE $D_f$}
        child { node {Dénominateur $\neq 0$} }
        child { node {Radicande $\ge 0$} }
        child { node {Arg $\ln > 0$} }
    }
    child[concept color=poppurple] { node {2. SYMÉTRIES}
        child { node {Paire : $f(-x)=f(x)$} }
        child { node {Impaire : $f(-x)=-f(x)$} }
        child { node {Périodique $T$} }
    }
    child[concept color=poppurple] { node {3. LIMITES \& ASYMPTOTES}
        child { node {Bornes de $D_f$} }
        child { node {$\pm\infty$} }
        child { node {AV : $x=x_0$} }
        child { node {AH : $y=\ell$} }
        child { node {AO : $y=ax+b$} }
    }
    child[concept color=poppurple] { node {4. DÉRIVÉE \& VARIATIONS}
        child { node {Calculer $f'$} }
        child { node {Signe de $f'$} }
        child { node {Tableau de variations} }
    }
    child[concept color=poppurple] { node {5. POINTS REMARQUABLES}
        child { node {$f(0)$, $f(x)=0$} }
        child { node {Extrema} }
        child { node {Inflexion (si $f''$)} }
    }
    child[concept color=poppurple] { node {6. TRACÉ $\mathcal{C}_f$}
        child { node {Repère + unités} }
        child { node {Asymptotes (---)} }
        child { node {Points + trous ($\circ$)} }
        child { node {Branches + flèches} }
    };
\end{tikzpicture}
\end{center}
\legende{Checklist visuelle de l'étude complète d'une fonction. Suivez les branches dans l'ordre numérique.}
\end{popfigure}

\begin{exercice}[Entraînement au tracé]
Étudier la fonction $f(x) = x + \frac{1}{x}$ sur $D_f = \mathbb{R}^*$.
\begin{enumerate}
    \item Déterminer la parité, les limites aux bornes, les asymptotes.
    \item Calculer $f'(x)$, dresser le tableau de variations.
    \item Tracer $\mathcal{C}_f$ dans un repère orthonormé (unité : 2 cm).
    \item Résoudre graphiquement $f(x) = 3$ et $f(x) \le -2$.
\end{enumerate}
\end{exercice}

\begin{corrige}[Exercice 1]
\begin{enumerate}
    \item $f$ impaire ($f(-x)=-f(x)$). $\lim_{0^+} f = +\infty$, $\lim_{0^-} f = -\infty$ $\Rightarrow$ AV $x=0$ (axe $Oy$). $\lim_{\pm\infty} f(x)/x = 1$, $\lim_{\pm\infty} (f(x)-x) = 0$ $\Rightarrow$ AO $y=x$ en $\pm\infty$.
    \item $f'(x) = 1 - 1/x^2 = (x^2-1)/x^2$. Zéros : $x=\pm 1$. Signe : $+$ sur $]-\infty, -1[$, $-$ sur $]-1, 0[$, $-$ sur $]0, 1[$, $+$ sur $]1, +\infty[$. Max en $-1$ : $f(-1)=-2$. Min en $1$ : $f(1)=2$.
    \item Tracé : Branches dans quadrants 1 et 3 (impair). Asymptotes $x=0$ et $y=x$. Max $(-1,-2)$, Min $(1,2)$. \textbf{Concave} sur $]-\infty,0[$ ($f''(x)=2/x^3 < 0$), \textbf{convexe} sur $]0,+\infty[$ ($f''(x)>0$).
    \item $f(x)=3 \iff x+1/x=3 \iff x^2-3x+1=0 \Rightarrow x = \frac{3\pm\sqrt{5}}{2} \approx 0,38 ; 2,62$. Deux solutions. $f(x) \le -2 \iff x \in ]-\infty, 0[$ (lecture graphique : les deux branches du 3ème quadrant sont sous $y=-2$, avec l'égalité en $x=-1$).
\end{enumerate}
\end{corrige}

\coursec{Modélisation et résolution de problèmes concrets (Synthèse)}

Après avoir maîtrisé la représentation graphique et l'étude complète des fonctions (continuité, monotonie, limites, variations), nous sommes désormais outillés pour passer de l'abstraction mathématique à la résolution de problèmes réels. C'est l'objectif final de cette unité : utiliser les fonctions comme des \cle{modèles} pour comprendre, prévoir et optimiser des phénomènes techniques, physiques ou économiques. Cette section synthétise l'ensemble des savoirs de la leçon.

\courssub{Du réel au modèle mathématique : la démarche de modélisation}

La modélisation n'est pas un simple exercice de traduction ; c'est un cycle itératif. Avant de manipuler des formules, il faut comprendre la situation physique.

\begin{methode}[Résoudre un problème de modélisation]
\begin{enumerate}
    \item \textbf{Lire et Comprendre :} Identifier les grandeurs physiques en jeu, leurs unités (SI de préférence), la variable indépendante (souvent le temps $t$ ou une quantité $x$) et la grandeur cherchée. Repérer les contraintes physiques (domaine de définition réel).
    \item \textbf{Modéliser :} Écrire la fonction $f$ qui relie la variable $x$ à la grandeur $y$. Déterminer l'intervalle $I$ (souvent un segment $[a,b]$) correspondant au domaine physique.
    \item \textbf{Étudier la fonction sur $I$ :} Domaine, continuité, dérivée, variations, limites aux bornes. Appliquer le TVI (Théorème des Valeurs Intermédiaires) pour l'existence, la monotonie stricte pour l'unicité. Chercher les extrema (min/max) pour l'optimisation.
    \item \textbf{Répondre à la question :} Calculer une valeur, encadrer une solution, trouver un optimum, donner une interprétation graphique.
    \item \textbf{Conclure et Valider :} Rédiger une phrase complète dans le contexte de l'énoncé, avec les unités. Discuter la validité du modèle (hypothèses simplificatrices, domaine de validité).
\end{enumerate}
\end{methode}

\begin{popfigure}
\begin{tikzpicture}[node distance=1.8cm, >=Stealth, font=\small]
    \node[draw, rounded corners, fill=popblueL, thick] (real) {Réalité / Situation concrète};
    \node[draw, rounded corners, fill=popgreenL, thick, right=of real, align=center] (model){Modèle Mathématique\\ (Fonction $f$, Domaine $I$)};
    \node[draw, rounded corners, fill=poporangeL, thick, right=of model, align=center] (math){Résolution Mathématique\\ (TVI, Monotonie, Extrema, Limites)};
    \node[draw, rounded corners, fill=poppinkL, thick, below=of model, align=center] (interp){Interprétation\\ (Réponse chiffrée + Unités + Sens)};
    \node[draw, rounded corners, fill=popgoldL, thick, left=of interp, align=center] (valid){Validation / Critique\\ (Hypothèses, Précision, Amélioration)};
    
    \draw[thick, ->] (real) -- node[above] {1. Modéliser} (model);
    \draw[thick, ->] (model) -- node[above] {2. Résoudre} (math);
    \draw[thick, ->] (math) -- node[right] {3. Interpréter} (interp);
    \draw[thick, ->] (interp) -- node[below] {4. Valider} (valid);
    \draw[thick, ->] (valid) -- node[left] {5. Retour / Affiner} (real);
    
    \draw[thick, ->, popmuted, dashed] (math) to[bend left=30] node[above, popmuted] {Vérification cohérence} (real);
\end{tikzpicture}
\legende{Le cycle de la modélisation mathématique : un aller-retour constant entre le réel et le modèle.}
\end{popfigure}

\courssub{Exemples techniques et physiques : température, marée, projectile, pharmacocinétique}

Les fonctions trigonométriques et irrationnelles modèlent naturellement des phénomènes périodiques ou à variation de racine carrée.

\begin{exemplebox}[Température de fermentation (Fonction sinusoïdale)]
Dans une cuve de fermentation, la température $T$ (en $^\circ$C) varie au cours du temps $t$ (en heures) selon le modèle :
\[ T(t) = 50 + 5\sin\left(\frac{\pi t}{12}\right), \quad t \in [0, 24]. \]
Le processus impose que la température reste dans l'intervalle $[48 ; 55]$. Étudier la conformité du processus sur 24h.
\end{exemplebox}

\begin{exoresolu}[Analyse de la température de fermentation]
\textbf{Énoncé :} Déterminer les intervalles de temps où le processus n'est pas conforme (température $< 48^\circ$C ou $> 55^\circ$C).

\textbf{Solution détaillée :}
\begin{enumerate}
    \item \textbf{Domaine et Continuité :} $T$ est définie et continue sur $\mathbb{R}$, donc sur le segment $[0, 24]$.
    \item \textbf{Variations :} $T'(t) = 5 \cdot \frac{\pi}{12} \cos\left(\frac{\pi t}{12}\right) = \frac{5\pi}{12}\cos\left(\frac{\pi t}{12}\right)$.
        \begin{itemize}
            \item $T'(t) > 0 \iff \cos(\frac{\pi t}{12}) > 0 \iff \frac{\pi t}{12} \in [0, \frac{\pi}{2}[ \cup ]\frac{3\pi}{2}, 2\pi] \iff t \in [0, 6[ \cup ]18, 24]$.
            \item $T'(t) < 0 \iff t \in ]6, 18[$.
        \end{itemize}
        $T$ est croissante sur $[0,6]$, décroissante sur $[6,18]$, croissante sur $[18,24]$.
    \item \textbf{Extrema sur $[0,24]$ :}
        \begin{align*}
            T(0) &= 50 + 5\sin(0) = 50. \\
            T(6) &= 50 + 5\sin(\frac{\pi}{2}) = 55 \quad \text{(Maximum global)}. \\
            T(18) &= 50 + 5\sin(\frac{3\pi}{2}) = 45 \quad \text{(Minimum global)}. \\
            T(24) &= 50 + 5\sin(2\pi) = 50.
        \end{align*}
    \item \textbf{Recherche des seuils (TVI + Monotonie) :}
        \begin{itemize}
            \item \textbf{Seuil $55^\circ$C :} Atteint au maximum $t=6$. $T(t) \le 55$ partout. \textbf{Jamais dépassé}.
            \item \textbf{Seuil $48^\circ$C :} $T(0)=50 > 48$, $T(18)=45 < 48$. Par TVI sur $[0,18]$ (continu), $\exists! t_1 \in ]0,18[$ tel que $T(t_1)=48$ (monotonie par morceaux).
                Résolution : $50 + 5\sin(\frac{\pi t}{12}) = 48 \iff \sin(\frac{\pi t}{12}) = -0,4$.
                $\frac{\pi t}{12} = \arcsin(-0,4) \approx -0,4115$ (rejeté, négatif) ou $\pi - (-0,4115) \approx 3,553$.
                $t_1 \approx \frac{12}{\pi} \times 3,553 \approx 13,57$ h.
                Sur $[18,24]$, $T(18)=45 < 48$, $T(24)=50 > 48$. $\exists! t_2 \in ]18,24[$ tel que $T(t_2)=48$.
                $\frac{\pi t}{12} = 2\pi + \arcsin(-0,4) \approx 6,283 - 0,4115 = 5,8715$.
                $t_2 \approx \frac{12}{\pi} \times 5,8715 \approx 22,43$ h.
        \end{itemize}
    \item \textbf{Conclusion interprétée :} La température descend sous $48^\circ$C entre \textbf{13h34 et 22h26} (environ). Le processus n'est \textbf{pas conforme} sur cet intervalle d'environ 8h50. La température ne dépasse jamais $55^\circ$C.
\end{enumerate}
\end{exoresolu}

\begin{popfigure}
\begin{tikzpicture}
\begin{axis}[
    width=\linewidth, height=8cm,
    axis lines=middle,
    xlabel={$t$ (heures)}, ylabel={$T(t)$ ($^\circ$C)},
    xmin=0, xmax=24.5, ymin=43, ymax=57,
    xtick={0,6,12,13.57,18,22.43,24},
    xticklabels={$0$,$6$,$12$,$t_1$,$18$,$t_2$,$24$},
    ytick={45,48,50,55},
    yticklabels={$45$,$48$,$50$,$55$},
    domain=0:24, samples=300,
    clip=false,
    legend style={at={(0.98,0.98)}, anchor=north east, fill=white, draw=popline},
]
\addplot[popcurve, thick, popblue] {50 + 5*sin(deg(pi*x/12))};
\addlegendentry{$T(t) = 50 + 5\sin(\pi t/12)$}

% Zones hors normes hachurées
\addplot[fill=popred!20, draw=popred, pattern=north east lines, pattern color=popred, opacity=0.5, domain=13.57:22.43, samples=100] {50 + 5*sin(deg(pi*x/12))} \closedcycle;
\addlegendentry{\small Zone $T < 48^\circ$C (Non conforme)}

% Lignes de seuils
\draw[dashed, popgreen, thick] (axis cs:0,48) -- (axis cs:24,48) node[right, pos=0.95, black] {$48^\circ$C};
\draw[dashed, poporange, thick] (axis cs:0,55) -- (axis cs:24,55) node[right, pos=0.95, black] {$55^\circ$C};

% Points caractéristiques
\addplot[only marks, mark=*, mark size=2.5, popred] coordinates {(6,55) (18,45) (13.57,48) (22.43,48)};
\end{axis}
\end{tikzpicture}
\legende{Modélisation de la température de fermentation. Les zones hachurées indiquent la non-conformité ($T < 48^\circ$C).}
\end{popfigure}

\begin{exemplebox}[Hauteur d'une marée (Modèle sinusoïdal)]
La hauteur de l'eau $h(t)$ (en mètres) dans un port est modélisée par $h(t) = 3 + 2\cos\left(\frac{\pi}{6}t - \frac{\pi}{3}\right)$ pour $t \in [0, 24]$ ($t=0$ minuit). Un bateau a besoin de $4,5$ m d'eau pour entrer. Déterminer les créneaux horaires d'accès.
\end{exemplebox}

\begin{exemplebox}[Portée d'un projectile (Trigonométrie)]
Un projectile est lancé avec une vitesse initiale $v_0 = 50$ m/s sous un angle $\alpha$. Sa portée (distance horizontale max) est $D(\alpha) = \frac{v_0^2}{g}\sin(2\alpha)$ avec $g \approx 10$ m/s$^2$, $\alpha \in [0, \frac{\pi}{2}]$. Chercher l'angle pour atteindre une cible à $200$ m. (Équation $\sin(2\alpha) = 0,8$, deux solutions symétriques par rapport à $45^\circ$).
\end{exemplebox}

\begin{exemplebox}[Concentration médicamenteuse (Fonction irrationnelle / exponentielle)]
La concentration $C(t)$ (mg/L) d'un médicament dans le sang suit $C(t) = \frac{D}{V} \frac{k_a}{k_a - k_e}(e^{-k_e t} - e^{-k_a t})$ (modèle à 2 compartiments). Pour un modèle simplifié $C(t) = 2\sqrt{t}e^{-0,5t}$ sur $[0, +\infty[$. Étudier le pic de concentration (max) et le temps pour redescendre sous le seuil thérapeutique $0,5$ mg/L (TVI + monotonie sur les branches).
\end{exemplebox}

\courssub{Exemples économiques et de gestion : Coûts, Revenus, Optimisation}

En gestion, on cherche souvent à minimiser un coût ou maximiser un profit. Les fonctions sont souvent rationnelles ($C_m(x) = \frac{C(x)}{x}$) ou irrationnelles.

\begin{propriete}[Théorème des Bornes (Extrema sur un segment)]
Soit $f$ une fonction \textbf{continue} sur un \textbf{segment} $[a,b]$ (fermé et borné). Alors $f$ est \textbf{bornée} et \textbf{atteint ses bornes} : il existe $m, M \in [a,b]$ tels que $\forall x \in [a,b],\ f(m) \le f(x) \le f(M)$.
\par \textbf{Recherche pratique des candidats au min/max sur $[a,b]$ :}
\begin{enumerate}
    \item Les bornes $a$ et $b$.
    \item Les points $c \in ]a,b[$ où $f'(c) = 0$ (extrema locaux).
    \item Les points $c \in ]a,b[$ où $f'$ n'existe pas (ex: racine carrée en 0, valeur absolue).
\end{enumerate}
\emph{Preuve exigible au Probatoire : Non (admis), mais l'application est systématique.}
\end{propriete}

\begin{exemplebox}[Coût moyen et optimisation de production]
Une entreprise a un coût total de production $C(x) = 1000 + 50x + 0,1x^2$ (en milliers de FCFA) pour $x$ unités produites ($x \in [10, 500]$). Le coût moyen est $C_m(x) = \frac{C(x)}{x} = \frac{1000}{x} + 50 + 0,1x$.
\begin{enumerate}
    \item Étudier les variations de $C_m$ sur $[10, 500]$.
    \item Déterminer la production $x$ qui minimise le coût moyen et la valeur de ce minimum.
\end{enumerate}
\end{exemplebox}

\begin{exoresolu}[Optimisation du coût moyen]
\textbf{Solution :}
\begin{enumerate}
    \item \textbf{Domaine / Continuité / Dérivabilité :} $C_m$ est continue et dérivable sur $]0, +\infty[$, donc sur $[10, 500]$.
    \item \textbf{Dérivée :} $C_m'(x) = -\frac{1000}{x^2} + 0,1 = \frac{0,1x^2 - 1000}{x^2} = \frac{0,1(x^2 - 10000)}{x^2} = \frac{0,1(x-100)(x+100)}{x^2}$.
    \item \textbf{Signe de $C_m'$ sur $[10, 500]$ :} Dénominateur $>0$. Numérateur : $(x-100)(x+100)$.
        \begin{itemize}
            \item $x \in [10, 100[$ : $x-100 < 0 \Rightarrow C_m'(x) < 0$ (décroît).
            \item $x = 100$ : $C_m'(100) = 0$.
            \item $x \in ]100, 500]$ : $x-100 > 0 \Rightarrow C_m'(x) > 0$ (croît).
        \end{itemize}
    \item \textbf{Tableau de variations et candidats au minimum :}
    \begin{center}
    \begin{tabular}{|c|ccccc|}\hline
    $x$ & $10$ & & $100$ & & $500$ \\ \hline
    $C_m'(x)$ & & $-$ & $0$ & $+$ & \\ \hline
    $C_m(x)$ & $151$ & $\searrow$ & $\mathbf{70}$ & $\nearrow$ & $102$ \\ \hline
    \end{tabular}
    \end{center}
    \item \textbf{Conclusion :} Le coût moyen est minimal pour une production de \textbf{100 unités}. La valeur minimale est \textbf{70 kFCFA/unité} (soit 70 000 FCFA/unité). Produire moins de 100 unités (sous-capacité) ou plus de 100 unités (sur-coûts) augmente le coût unitaire.
\end{enumerate}
\end{exoresolu}

\begin{popfigure}
\begin{tikzpicture}
\begin{axis}[
    width=\linewidth, height=8cm,
    axis lines=middle,
    xlabel={$x$ (unités)}, ylabel={$C_m(x)$ (kFCFA)},
    xmin=0, xmax=520, ymin=60, ymax=170,
    xtick={10,100,500},
    ytick={70,102,151},
    yticklabels={$70$,$102$,$151$},
    domain=10:500, samples=300,
    clip=false,
    legend style={at={(0.98,0.98)}, anchor=north east, fill=white, draw=popline},
]
\addplot[popcurve, thick, popgreen] {1000/x + 50 + 0.1*x};
\addlegendentry{$C_m(x) = \frac{1000}{x} + 50 + 0,1x$}
% Minimum point
\addplot[only marks, mark=*, mark size=3, popred] coordinates {(100,70)} node[above right, black] {Min $(100; 70)$};
% Asymptote verticale suggestion (x=0 hors domaine)
\draw[dashed, popmuted] (axis cs:0,60) -- (axis cs:0,170);
\end{axis}
\end{tikzpicture}
\legende{Graphique du coût moyen $C_m(x)$. Le minimum global sur $[10,500]$ est atteint en $x=100$.}
\end{popfigure}

\courssub{Puissance du couple TVI + Monotonie : Existence et Unicité}

C'est le cœur de la résolution d'équations $f(x) = k$ sur un intervalle.

\begin{aretenir}[Stratégie « Existence + Unicité »]
Pour montrer que l'équation $f(x) = k$ admet \textbf{une unique solution} sur un intervalle $I$ :
\begin{enumerate}
    \item \textbf{Existence :} Vérifier que $f$ est \textbf{continue} sur $I$ (souvent un segment $[a,b]$) et que $k$ est \textbf{entre $f(a)$ et $f(b)$} (ou entre les bornes de l'image). Invoquer le \textbf{TVI}.
    \item \textbf{Unicité :} Montrer que $f$ est \textbf{strictement monotone} (croissante ou décroissante) sur $I$ (signe constant de $f'$ ou étude de $f(x_2)-f(x_1)$). Une fonction strictement monotone est injective, donc l'antécédent est unique.
\end{enumerate}
\emph{Astuce :} Si $f$ n'est pas monotone sur tout $I$, la découper en intervalles de monotonie et appliquer la méthode sur chaque morceau.
\end{aretenir}

\begin{exemplebox}[Bassin de vidange (Fonction irrationnelle : racine carrée)]
Un bassin se vide. La hauteur d'eau $h$ (en mètres) en fonction du temps $t$ (en heures) est modélisée par :
\[ h(t) = 2\sqrt{9 - t}, \quad t \in [0, 9]. \]
\begin{enumerate}
    \item À quel instant la hauteur est-elle de $1$ m ?
    \item Le bassin est-il vide à $t=9$ ?
    \item Calculer la vitesse de vidange (dérivée) à $t=5$ h et l'interpréter.
\end{enumerate}
\end{exemplebox}

\begin{exoresolu}[Vidange du bassin]
\textbf{Solution :}
\begin{enumerate}
    \item \textbf{Équation $h(t) = 1$ :} $2\sqrt{9-t} = 1 \iff \sqrt{9-t} = 0,5 \iff 9-t = 0,25 \iff t = 8,75$ h.
        \textbf{Vérification TVI/Unicité :} $h$ continue sur $[0,9]$, strictement décroissante ($h'(t) = -\frac{1}{\sqrt{9-t}} < 0$). $h(0)=6$, $h(9)=0$. $1 \in [0,6]$. $\Rightarrow$ \textbf{Unique solution $t=8,75$ h (soit 8h45)}.
    \item \textbf{Vidange à $t=9$ :} $h(9) = 2\sqrt{0} = 0$. \textbf{Oui, le bassin est vide à $t=9$ h}.
    \item \textbf{Vitesse de vidange :} $v(t) = h'(t) = 2 \cdot \frac{-1}{2\sqrt{9-t}} = -\frac{1}{\sqrt{9-t}}$.
        $v(5) = -\frac{1}{\sqrt{4}} = -0,5$ m/h.
        \textbf{Interprétation :} À $t=5$ h, le niveau de l'eau baisse de \textbf{0,5 mètre par heure}. Le signe négatif confirme la baisse. La vitesse augmente (en valeur absolue) au fur et à mesure que le bassin se vide ($\lim_{t\to 9^-} v(t) = -\infty$ : modèle physique limite, en réalité la loi de Torricelli donne une vitesse finie).
\end{enumerate}
\end{exoresolu}

\courssub{Rédaction de la conclusion : l'art de l'interprétation}

Au Probatoire, le calcul juste sans phrase de conclusion perd des points. La conclusion doit être une \textbf{phrase affirmative, contextualisée, avec unités}.

\begin{attention}[Erreurs fréquentes en modélisation (Sanctions au Probatoire)]
\begin{itemize}
    \item \textbf{Oublier les unités :} Écrire « $t = 8,75$ » au lieu de « $t = 8,75$ h » ou « 8 h 45 ».
    \item \textbf{Ignorer le domaine physique :} Donner $t = -2$ h ou $x = -50$ unités comme solution. \cle{Toujours vérifier que la solution appartient à l'intervalle de définition physique $I$}.
    \item \textbf{Ne pas justifier l'unicité :} La question demande « \emph{l'instant unique} » ou « \emph{la valeur de $x$} ». Répondre seulement par le calcul algébrique sans invoquer la monotonie stricte (ou TVI + monotonie) est incomplet.
    \item \textbf{Confondre « maximum » et « valeur maximale » :} Dire « le max est $x=100$ » (faux, c'est l'abscisse) au lieu de « le coût minimal est $70$ kFCFA pour $x=100$ ».
    \item \textbf{Notation dérivée :} Écrire $f'(x) = 0 \Rightarrow x = ...$ sans dire « $f'(x)$ s'annule pour $x=...$ ».
\end{itemize}
\end{attention}

\begin{exemplebox}[Modèle de rédaction type]
\begin{quote}
\textit{« La fonction $C_m$ modélisant le coût moyen est continue et dérivable sur $[10; 500]$. Elle admet un minimum global en $x = 100$ où $C_m(100) = 70$. \textbf{Le coût moyen minimal est de 70 000 FCFA par unité, atteint pour une production de 100 unités.} Pour toute autre production dans l'intervalle, le coût unitaire est supérieur. »}
\end{quote}
\end{exemplebox}

\courssub{Validation du modèle : esprit critique et limites}

Un modèle mathématique est une \textbf{simplification} de la réalité. La dernière étape (étape 5 de la méthode) consiste à en discuter la fiabilité.

\begin{savaistu}[Modélisation et Probatoire Technique]
Les sujets du Probatoire / Baccalauréat Technique comportent \textbf{quasiment systématiquement} un exercice de « Situation concrète » (souvent le 3ème ou 4ème exercice), noté sur \textbf{6 à 8 points}. La répartition type est :
\begin{itemize}
    \item 2-3 pts : Modélisation (écriture fonction, domaine, unités).
    \item 2-3 pts : Étude mathématique (dérivée, TVI, variations, calculs).
    \item 2 pts : \textbf{Interprétation, conclusion rédigée, unités, critique du modèle}.
\end{itemize}
Négliger la rédaction ou les unités, c'est perdre \textbf{le tiers des points} de l'exercice alors que les calculs sont justes.
\end{savaistu}

\begin{aretenir}[Points de vigilance pour la validation]
\begin{enumerate}
    \item \textbf{Hypothèses simplificatrices :} Linéarisation (loi de Hooke, loi d'Ohm), négligence des frottements, température/pression constantes, milieu homogène.
    \item \textbf{Domaine de validité :} Le modèle $h(t)=2\sqrt{9-t}$ prédit une vitesse infinie à la fin ($t\to 9$). Physiquement impossible. Le modèle n'est valide que tant que $h$ n'est pas trop petit (loi de Torricelli : $v \propto \sqrt{h}$).
    \item \textbf{Précision des paramètres :} $g=9,81$ ou $10$ ? Coefficient de frottement connu à $\pm 10\%$ ? Une erreur de $5\%$ sur un paramètre peut déplacer l'optimum.
    \item \textbf{Extrapolation interdite :} Ne jamais utiliser le modèle en dehors de l'intervalle $[a,b]$ ayant servi à sa construction (ex: prévoir la température à $t=30$h avec un modèle calibré sur $[0,24]$).
\end{enumerate}
\end{aretenir}

\begin{exemplebox}[Discussion critique : Exemple du bassin]
\textit{« Le modèle $h(t)=2\sqrt{9-t}$ suppose un débit de vidange proportionnel à $\sqrt{h}$ (loi de Torricelli) et une section constante du bassin. Il néglige la viscosité, la turbulence à l'orifice, et l'évaporation. La dérivée tendant vers $-\infty$ quand $t\to 9$ signale la rupture de validité du modèle pour les toutes dernières minutes. En réalité, le bassin se vide en un temps fini avec une vitesse finie. Les prévisions sont fiables pour $h > 0,5$ m (soit $t < 8,94$ h). »}
\end{exemplebox}

\begin{exercice}[Synthèse : Production d'énergie éolienne]
La puissance $P(v)$ (en kW) fournie par une éolienne en fonction de la vitesse du vent $v$ (en m/s) est modélisée par :
\[ P(v) = \begin{cases}
0 & \text{si } 0 \le v < 3 \text{ (seuil de démarrage)} \\
0,5 v^3 & \text{si } 3 \le v \le 12 \text{ (zone nominale)} \\
864 & \text{si } 12 < v \le 25 \text{ (puissance max, bridage)} \\
0 & \text{si } v > 25 \text{ (arrêt sécurité)}
\end{cases} \]
\begin{enumerate}
    \item Tracer la courbe représentative de $P$ sur $[0, 30]$.
    \item Étudier la continuité de $P$ en $v=3$, $v=12$, $v=25$.
    \item Le vent varie selon $v(t) = 10 + 5\sin(\frac{\pi t}{12})$ pour $t \in [0, 24]$. Exprimer la puissance en fonction du temps $P(t)$. Y a-t-il des instants où l'éolienne s'arrête ? Produit-elle sa puissance max ?
\end{enumerate}
\end{exercice}

\begin{corrige}[Exercice : Production d'énergie éolienne]
\begin{enumerate}
    \item \textbf{Graphique :} Par morceaux. Nulle sur $[0,3[$, cubique $0,5v^3$ sur $[3,12]$ (passant par $(3, 13,5)$ et $(12, 864)$), constante $864$ sur $]12, 25]$, nulle sur $]25, 30]$.
    \item \textbf{Continuité :}
        \begin{itemize}
            \item $v=3$ : $P(3^-) = 0$, $P(3) = 0,5 \times 27 = 13,5$. \textbf{Discontinue (saut)}. Physiquement : inertie du rotor.
            \item $v=12$ : $P(12^-) = 0,5 \times 1728 = 864$, $P(12) = 864$. \textbf{Continue}. Transition douce vers le bridage.
            \item $v=25$ : $P(25^-) = 864$, $P(25) = 0$. \textbf{Discontinue (saut)}. Arrêt d'urgence.
        \end{itemize}
    \item \textbf{Puissance temporelle :} $v(t) \in [5, 15]$.
        \begin{itemize}
            \item $v(t) \ge 5 > 3$ : jamais à l'arrêt par manque de vent.
            \item $v(t) \le 15 < 25$ : jamais à l'arrêt par excès de vent.
            \item $v(t) > 12 \iff 10 + 5\sin(\frac{\pi t}{12}) > 12 \iff \sin(\frac{\pi t}{12}) > 0,4$.
                Solutions sur $[0, 24]$ : $\frac{\pi t}{12} \in ]\arcsin(0,4), \pi - \arcsin(0,4)[ \approx ]0,4115, 2,73[$.
                $t \in ]1,57, 10,43[$ (environ \textbf{1h34 à 10h26}).
                Sur cet intervalle, $P(t) = 864$ kW (puissance max).
            \item Sinon ($t \in [0, 1,57] \cup [10,43, 24]$), $3 \le v(t) \le 12$, $P(t) = 0,5(10 + 5\sin(\frac{\pi t}{12}))^3$.
        \end{itemize}
        \textbf{Conclusion :} L'éolienne fonctionne en continu sur 24h. Elle produit sa puissance nominale maximale (864 kW) pendant environ \textbf{8h50} (de 1h34 à 10h26).
\end{enumerate}
\end{corrige}

\courssub{Bilan de l'unité : Votre « Boîte à outils » pour le Probatoire}

\begin{center}
\begin{tabularx}{\linewidth}{|l|X|X|}\hline
\rowcolor{popblueL} \textbf{Notion Clé} & \textbf{Quand l'utiliser ? (Mots-clés de l'énoncé)} & \textbf{Écriture attendue (Rigueur)} \\ \hline
\textbf{Continuité sur $[a,b]$} & « Fonction continue », « Sans rupture », « TVI applicable » & Vérifier $\lim_{x\to x_0} f(x) = f(x_0)$ ou opérations sur fonctions continues. \\ \hline
\textbf{TVI (Théorème Valeurs Intermédiaires)} & « Montrer qu'il existe au moins une solution », « Prouver que la courbe coupe la droite $y=k$ » & 1. $f$ continue sur $[a,b]$. 2. $k \in [f(a), f(b)]$ (ou $[f(b), f(a)]$). 3. $\Rightarrow \exists c \in [a,b], f(c)=k$. \\ \hline
\textbf{Monotonie stricte / Bijection} & « Montrer qu'il existe \textbf{une unique} solution », « L'instant précis », « La valeur exacte de $x$ » & TVI (existence) + $f$ strictement mono sur $I$ (unicité par injectivité). \\ \hline
\textbf{Limites remarquables} & $\frac{\sin X}{X} \to 1$, $\frac{\tan X}{X} \to 1$, $\frac{\ln(1+X)}{X} \to 1$, $\frac{e^X-1}{X} \to 1$ & Poser $X = \dots$, faire apparaître la forme standard. Ne pas écrire $\frac{0}{0}$. \\ \hline
\textbf{Fonctions irrationnelles (racines)} & Domaine, limites en $+\infty$ (raisonnement sur terme dominant), rationalisation & $\sqrt{P(x)} \sim \sqrt{a_n}x^{n/2}$. Rationaliser pour lever indétermination $\infty-\infty$ ou $0/0$. \\ \hline
\textbf{Optimisation sur $[a,b]$} & « Minimiser le coût », « Maximiser le profit/portée », « Production optimale » & 1. $f$ continue sur $[a,b]$. 2. Candidats : $a, b, \{f'=0\} \cap ]a,b[$. 3. Comparer valeurs. 4. Conclure avec unité. \\ \hline
\textbf{Représentation graphique} & « Tracer la courbe », « Lecture graphique », « Encadrer graphiquement » & Repère orthonormé, unités, échelle, points caractéristiques, asymptotes, variations respectées. \\ \hline
\end{tabularx}
\end{center}

\begin{methode}[Check-list finale avant de rendre la copie (Modélisation)]
\begin{enumerate}
    \item \textbf{Variables :} $x$ (ou $t$) défini ? Unités notées ? Domaine physique $[a,b]$ explicitement écrit ?
    \item \textbf{Fonction :} Expression $f(x)$ correcte ? Continuité / Dérivabilité sur $]a,b[$ justifiée ?
    \item \textbf{Calculs :} Dérivée simplifiée ? Signe étudié par produit/quotient de facteurs ? Tableau de variations complet (3 lignes, même nb colonnes) ?
    \item \textbf{Théorèmes :} TVI invoqué explicitement (continuité + encadrement) ? Monotonie stricte invoquée pour l'unicité ?
    \item \textbf{Réponses :} Valeurs numériques \textbf{avec unités} ? Appartenance au domaine $[a,b]$ vérifiée ?
    \item \textbf{Conclusion :} Phrase complète, langage courant, réponse à la question posée ?
    \item \textbf{Validation :} Une phrase sur les limites du modèle (si demandé) ?
\end{enumerate}
\end{methode}

\begin{savaistu}[L'héritage de Cauchy et l'ingénierie moderne]
Augustin-Louis Cauchy (1789-1857) a formalisé la notion de limite et de continuité pour mettre le calcul infinitésimal sur des bases rigoureuses, justement parce que les ingénieurs de son temps (ponts, routes, machines à vapeur) avaient besoin de \textbf{garanties mathématiques} pour leurs calculs de résistance des matériaux et de stabilité. Aujourd'hui, en Terminale Technique, vous utilisez ses théorèmes (TVI, Bornes, Monotonie) pour les mêmes raisons : \textbf{certifier qu'une pièce ne cassera pas, qu'un coût est minimal, qu'un processus reste dans les normes}. Les mathématiques ne sont pas une abstraction : ce sont les \textbf{certificats de conformité} du monde technique.
\end{savaistu}

Cette leçon sur les fonctions numériques d'une variable réelle s'achève. Vous possédez désormais la panoplie complète : \textbf{analyser} une fonction (domaine, limites, continuité, dérivée, variations, graphique), \textbf{résoudre} des équations et inéquations (TVI, monotonie), et \textbf{modéliser} des situations réelles pour prendre des décisions techniques ou économiques fondées. C'est le cœur du métier d'ingénieur ou de technicien supérieur. Bon courage pour le Probatoire !

\newpage

\coursec{Activité d'intégration}

\coursec{Fonctions numériques d'une variable réelle : Application à l'irrigation goutte-à-goutte}

\courssub{Objectifs de l'activité}
\begin{itemize}
    \item Mobiliser les notions de \cle{continuité}, de \cle{dérivée}, de \cle{variations}, de \cle{limites} et de \cle{théorème des valeurs intermédiaires (TVI)} dans une situation technique authentique.
    \item Modéliser un problème d'\cle{optimisation} (débit maximal) et de \cle{seuils} (débit minimal requis) à l'aide d'une fonction numérique définie sur un intervalle fermé.
    \item Interpréter mathématiquement des contraintes physiques (dimensions, unités, domaine de définition) et formuler une recommandation technique argumentée.
    \item Renforcer la rigueur de la rédaction mathématique : justifications des hypothèses des théorèmes, tableaux de variations complets, encadrements numériques.
\end{itemize}

\courssub{Situation}
Dans le cadre d'un projet d'agriculture résiliente au changement climatique dans la région de l'Extrême-Nord (Maroua), un jeune ingénieur agronome conçoit un système d'irrigation goutte-à-goutte à basse pression. Il dispose de tuyaux en PVC rigide, de section circulaire de rayon \textbf{R = 10 cm} et de longueur \textbf{L = 1 m}. Le principe est simple : l'eau arrive par gravité depuis un réservoir surélevé, remplit le tuyau horizontal bouché à ses deux extrémités, et s'écoule par une fente longitudinale unique, découpée sur toute la longueur du tuyau, vers les rangées de cultures.

La largeur de cette fente, notée \textbf{x} (en centimètres), est le paramètre de conception que l'ingénieur doit choisir. L'eau dans le tuyau maintient une hauteur constante \textbf{H = 20 cm} (mesurée du fond du tuyau à la surface libre) grâce à un trop-plein. D'après la loi de Torricelli corrigée pour une fente, le débit volumique \textbf{Q} (en litres par heure, L/h) s'exprime par la formule :
\[
Q(x) = k \cdot x \cdot \sqrt{2g\left(H - \frac{x}{2}\right)} \qquad \text{pour } x \in [0,\, 2H]
\]
où :
\begin{itemize}
    \item \textbf{g = 981 cm/s²} (accélération de la pesanteur, unités cohérentes avec x et H en cm) ;
    \item \textbf{k = 0,05} (constante de calibration sans dimension, intégrant la longueur du tuyau L = 100 cm et la conversion cm³/s $\to$ L/h : $k = \frac{L \times 3600}{1000} = 360$, mais ici ajustée à 0,05 pour tenir compte du coefficient de contraction du jet et des pertes de charge singulières, ce qui donne des débits réalistes pour du goutte-à-goutte).
\end{itemize}
Le domaine de définition physique est \textbf{x $\in$ [0 ; 40]} cm (la fente ne peut pas être plus large que le diamètre du tuyau, et l'argument de la racine doit rester positif).

L'agriculteur a besoin d'un débit minimal de \textbf{50 L/h} pour alimenter correctement son champ. L'ingénieur doit déterminer la largeur optimale \textbf{x$_{\text{opt}}$} qui maximise le débit, calculer ce débit maximal, et préciser l'intervalle de largeurs admissibles pour garantir au moins 50 L/h.

\courssub{Consignes}
\begin{enumerate}
    \item \textbf{Continuité et bornes :} Vérifier que la fonction $Q$ est continue sur l'intervalle $[0\,;\, 40]$. Calculer $Q(0)$ et $Q(40)$. Interpréter physiquement ces valeurs.
    \item \textbf{Étude des variations :} Calculer la dérivée $Q'(x)$ sur $]0\,;\, 40[$. Étudier son signe et dresser le tableau de variations complet de $Q$ sur $[0\,;\, 40]$.
    \item \textbf{Optimisation :} Déterminer la largeur exacte $x_{\text{opt}}$ (en cm) qui maximise le débit. Calculer la valeur approchée de $x_{\text{opt}}$ au dixième de cm. En déduire le débit maximal $Q_{\text{max}}$ (en L/h, arrondi à l'unité).
    \item \textbf{Seuil de 50 L/h (TVI) :} Montrer qu'il existe exactement deux largeurs $x_1$ et $x_2$ dans $[0\,;\, 40]$ telles que $Q(x) = 50$. Encadrer $x_1$ et $x_2$ avec une précision de 1 cm (par entiers consécutifs), puis affiner à 0,1 cm par dichotomie ou calculatrice.
    \item \textbf{Recommandation technique :} Rédiger une note technique courte (5-6 lignes) à destination de l'artisan soudeur, indiquant la largeur de coupe visée, la tolérance admissible pour rester au-dessus de 50 L/h, et les risques si la coupe est trop large ou trop étroite.
\end{enumerate}

\courssub{Ressources à mobiliser}
\begin{methode}[Rappels théoriques utiles]
\begin{itemize}
    \item \textbf{Continuité :} Produit, composition de fonctions continues (racine carrée, polynômes) sur un intervalle où elles sont définies.
    \item \textbf{Dérivée :} $(u \cdot v)' = u'v + uv'$ ; $\big(\sqrt{u}\big)' = \frac{u'}{2\sqrt{u}}$ ; $\big(u^\alpha\big)' = \alpha u' u^{\alpha-1}$.
    \item \textbf{Tableau de variations :} Une colonne par valeur remarquable (bornes, zéros de $f'$) + colonnes intermédiaires pour signes/flèches. Trois lignes : $x$, $f'(x)$, $f(x)$.
    \item \textbf{TVI (Théorème des Valeurs Intermédiaires) :} Si $f$ continue sur $[a,b]$, pour tout $k$ entre $f(a)$ et $f(b)$, $\exists c \in [a,b]$ tel que $f(c)=k$. Si $f$ est \emph{strictement} monotone, $c$ est unique.
    \item \textbf{Encadrement :} Méthode des valeurs successives (tableau de valeurs) ou dichotomie (milieu d'intervalle).
\end{itemize}
\end{methode}

\begin{attention}[Pièges fréquents]
\begin{itemize}
    \item Oublier de vérifier les hypothèses du TVI (continuité sur intervalle \emph{fermé}, valeur intermédiaire \emph{strictement} entre les bornes).
    \item Confondre $x$ (largeur en cm) et $Q(x)$ (débit en L/h) dans les encadrements.
    \item Négliger le domaine de définition : $\sqrt{H - x/2}$ impose $x \le 2H = 40$.
    \item Erreur de dérivation : oublier la dérivée de l'intérieur pour la racine (règle de la chaîne).
    \item Unités : $g$ en cm/s² car $x, H$ en cm. La constante $k$ absorbe les conversions.
\end{itemize}
\end{attention}

\courssub{Démarche et corrigé commenté}
\begin{exoresolu}[Démarche et corrigé commenté]
\textbf{1. Continuité et valeurs aux bornes}
La fonction $Q$ est définie sur $[0\,;\, 40]$ par $Q(x) = k \cdot x \cdot \sqrt{2g(H - x/2)}$.
\begin{itemize}
    \item $x \mapsto x$ est continue (polynôme) sur $\mathbb{R}$.
    \item $x \mapsto H - x/2 = 20 - 0,5x$ est continue (affine) sur $\mathbb{R}$. Sur $[0\,;\, 40]$, elle prend des valeurs dans $[0\,;\, 20]$, donc positives ou nulles.
    \item $t \mapsto \sqrt{t}$ est continue sur $[0\,;\, +\infty[$. La composition $x \mapsto \sqrt{20 - 0,5x}$ est donc continue sur $[0\,;\, 40]$.
    \item $x \mapsto \sqrt{2g(H - x/2)} = \sqrt{1962(20 - 0,5x)}$ est continue sur $[0\,;\, 40]$.
    \item Produit de fonctions continues : $Q$ est \textbf{continue sur $[0\,;\, 40]$}.
\end{itemize}
Calculs aux bornes :
\begin{align*}
Q(0) &= k \cdot 0 \cdot \sqrt{2g(20 - 0)} = \mathbf{0 \text{ L/h}}. \quad \text{(Pas de fente = pas d'écoulement)} \\
Q(40) &= k \cdot 40 \cdot \sqrt{2g(20 - 20)} = k \cdot 40 \cdot 0 = \mathbf{0 \text{ L/h}}. \quad \text{(Fente trop large = plus de charge hydraulique)}
\end{align*}

\textbf{2. Dérivée et variations}
Pour $x \in ]0\,;\, 40[$, posons $u(x) = x$ et $v(x) = \sqrt{2g(H - x/2)} = \sqrt{1962(20 - 0,5x)} = \sqrt{39240 - 981x}$.
$u'(x) = 1$.
$v(x) = (39240 - 981x)^{1/2} \implies v'(x) = \frac{1}{2}(39240 - 981x)^{-1/2} \cdot (-981) = -\frac{981}{2\sqrt{39240 - 981x}}$.
\begin{align*}
Q'(x) &= k \left[ u'(x)v(x) + u(x)v'(x) \right] \\
&= k \left[ \sqrt{39240 - 981x} + x \cdot \left( -\frac{981}{2\sqrt{39240 - 981x}} \right) \right] \\
&= k \cdot \frac{2(39240 - 981x) - 981x}{2\sqrt{39240 - 981x}} \\
&= k \cdot \frac{78480 - 1962x - 981x}{2\sqrt{39240 - 981x}} \\
&= k \cdot \frac{78480 - 2943x}{2\sqrt{39240 - 981x}} \\
&= k \cdot \frac{2943\left(\frac{80}{3} - x\right)}{2\sqrt{39240 - 981x}} \quad \left(\text{car } \frac{78480}{2943} = \frac{80}{3}\right) \\
&= \frac{2943k}{2} \cdot \frac{\frac{80}{3} - x}{\sqrt{39240 - 981x}}
\end{align*}
Sur $]0\,;\, 40[$, le dénominateur $2\sqrt{39240 - 981x} > 0$ et $\frac{2943k}{2} > 0$.
Le signe de $Q'(x)$ est donc le signe du numérateur $\frac{80}{3} - x$.
\begin{itemize}
    \item $Q'(x) > 0 \iff x < \frac{80}{3} \approx 26,67$ : $Q$ croît sur $\left[0\,;\, \frac{80}{3}\right]$.
    \item $Q'(x) = 0 \iff x = \frac{80}{3}$ : extremum.
    \item $Q'(x) < 0 \iff x > \frac{80}{3}$ : $Q$ décroît sur $\left[\frac{80}{3}\,;\, 40\right]$.
\end{itemize}

\noindent \textbf{Tableau de variations de $Q$ sur $[0\,;\, 40]$ :}
\begin{center}
\begin{tabular}{|c|ccccc|}\hline
$x$ & $0$ & & $\frac{80}{3}$ & & $40$ \\ \hline
$Q'(x)$ & & $+$ & $0$ & $-$ & \\ \hline
$Q(x)$ & $0$ & $\nearrow$ & $Q(\frac{80}{3})$ & $\searrow$ & $0$ \\ \hline
\end{tabular}
\end{center}

\textbf{3. Optimisation : $x_{\text{opt}}$ et $Q_{\text{max}}$}
Le maximum est atteint pour $x_{\text{opt}} = \frac{80}{3} \approx \mathbf{26,7 \text{ cm}}$.
Calcul de $Q_{\text{max}} = Q\left(\frac{80}{3}\right)$ :
$H - \frac{x_{\text{opt}}}{2} = 20 - \frac{40}{3} = \frac{20}{3} \approx 6,667 \text{ cm}$.
$\sqrt{2g(H - x_{\text{opt}}/2)} = \sqrt{1962 \times \frac{20}{3}} = \sqrt{13080} \approx 114,37 \text{ cm/s}$.
$Q_{\text{max}} = 0,05 \times \frac{80}{3} \times 114,37 \approx 0,05 \times 26,667 \times 114,37 \approx \mathbf{152,5 \text{ L/h}}$.
Arrondi : \textbf{$Q_{\text{max}} \approx 153 \text{ L/h}$}.

\textbf{4. Recherche des largeurs pour $Q(x) = 50$ L/h (TVI)}
$Q$ est continue sur $[0\,;\, 40]$, $Q(0)=0$, $Q(40)=0$, $Q_{\text{max}} \approx 153$.
Puisque $0 < 50 < 153$, par le TVI :
\begin{itemize}
    \item Sur $\left[0\,;\, \frac{80}{3}\right]$, $Q$ est continue et \textbf{strictement croissante} ($Q'$ $>$ 0), avec $Q(0)=0$ et $Q(80/3)\approx 153$. Il existe \textbf{un unique} $x_1 \in \left[0\,;\, \frac{80}{3}\right]$ tel que $Q(x_1)=50$.
    \item Sur $\left[\frac{80}{3}\,;\, 40\right]$, $Q$ est continue et \textbf{strictement décroissante} ($Q'$ $<$ 0), avec $Q(80/3)\approx 153$ et $Q(40)=0$. Il existe \textbf{un unique} $x_2 \in \left[\frac{80}{3}\,;\, 40\right]$ tel que $Q(x_2)=50$.
\end{itemize}

\noindent \textbf{Encadrement de $x_1$ (branche croissante) :}
On cherche $x \in [0\,;\, 26,7]$ tel que $Q(x)=50$.
\begin{center}
\begin{tabular}{|c|c|c|c|c|c|c|}\hline
$x$ (cm) & 5 & 5,4 & 5,43 & 5,44 & 5,5 & 6 \\ \hline
$Q(x)$ (L/h) & $\approx 54,8$ & $\approx 50,5$ & $\approx 50,0$ & $\approx 49,8$ & $\approx 48,1$ & $\approx 41,2$ \\ \hline
\end{tabular}
\end{center}
$Q(5,43) \approx 50,0 > 50 > 49,8 \approx Q(5,44)$. Donc $x_1 \in [5,43\,;\, 5,44]$.
\textbf{$x_1 \approx 5,4$ cm} (au dixième de cm).

\noindent \textbf{Encadrement de $x_2$ (branche décroissante) :}
On cherche $x \in [26,7\,;\, 40]$ tel que $Q(x)=50$.
\begin{center}
\begin{tabular}{|c|c|c|c|c|c|}\hline
$x$ (cm) & 39,0 & 39,3 & 39,35 & 39,4 & 39,5 \\ \hline
$Q(x)$ (L/h) & $\approx 61,1$ & $\approx 51,5$ & $\approx 50,0$ & $\approx 47,8$ & $\approx 43,7$ \\ \hline
\end{tabular}
\end{center}
$Q(39,35) \approx 50,0$. $Q(39,3) > 50 > Q(39,4)$. Donc $x_2 \in [39,3\,;\, 39,4]$.
\textbf{$x_2 \approx 39,4$ cm} (au dixième de cm).

\textbf{5. Recommandation technique}
\begin{aretenir}[Note technique pour l'atelier de découpe]
\textbf{Objet :} Paramètres de coupe fente PVC irrigation goutte-à-goutte (Maroua).
\begin{itemize}
    \item \textbf{Cible optimale :} Largeur de fente $x = \mathbf{26,7 \text{ cm}}$ (débit max $\approx 153 \text{ L/h}$).
    \item \textbf{Tolérance fonctionnelle (débit $\ge 50 \text{ L/h}$) :} Largeur comprise entre \textbf{5,4 cm et 39,4 cm}.
    \item \textbf{Consigne pratique :} Viser $26,7 \text{ cm} \pm 1 \text{ cm}$ (soit $[25,7\,;\, 27,7]$) pour rester très au-dessus du seuil et absorber les imprécisions de coupe.
    \item \textbf{Risques :} Coupe $< 5,4 \text{ cm}$ $\to$ débit insuffisant (sous-irrigation). Coupe $> 39,4 \text{ cm}$ $\to$ perte de charge, débit insuffisant et fragilisation mécanique du tuyau (risque de fente longitudinale par pression).
\end{itemize}
\end{aretenir}
\end{exoresolu}

\courssub{Bilan de l'activité}
Cette activité d'intégration a permis de mettre en œuvre l'ensemble des savoirs de la leçon dans un contexte professionnel réel (agriculture, hydraulique, fabrication) :
\begin{itemize}
    \item \textbf{Continuité :} Vérifiée rigoureusement par composition/produit sur un intervalle fermé, condition \emph{sine qua non} pour appliquer le TVI.
    \item \textbf{Dérivée et variations :} Le calcul de $Q'(x)$ a mobilisé la dérivation de produits et de fonctions composées (racine carrée), l'analyse du signe a nécessité une factorisation soignée. Le tableau de variations a synthétisé le comportement global (montée puis descente), typique des problèmes d'optimum avec contrainte géométrique.
    \item \textbf{Optimisation :} La recherche du maximum a utilisé le passage par la dérivée nulle et le changement de signe (croissant $\to$ décroissant), confirmant un maximum global sur l'intervalle fermé.
    \item \textbf{TVI et encadrement :} L'application du TVI a été distinguée sur chaque intervalle de monotonie stricte pour garantir l'unicité des solutions $x_1$ et $x_2$. La méthode de lecture de tableau de valeurs (ou dichotomie) a donné des encadrements numériques exploitables techniquement.
    \item \textbf{Modélisation :} L'interprétation physique des bornes ($Q=0$ pour $x=0$ et $x=40$) et de l'optimum a donné du sens aux nombres. La rédaction finale a traduit le langage mathématique en consignes d'atelier (tolérance, risques), compétence attendue en Terminale technique.
\end{itemize}

\begin{savaistu}[Le saviez-vous ?]
À Maroua, comme dans tout le Sahel, l'irrigation goutte-à-goutte permet d'économiser jusqu'à \textbf{60\% d'eau} par rapport à l'irrigation par gravité traditionnelle. Le dimensionnement précis des émetteurs (ici une simple fente) est crucial : une fente trop large ($x \to 40$ cm) vide le tuyau trop vite, la pression chute et l'eau n'atteint pas le bout de la ligne (100 m). Le modèle $Q(x)$ capture ce compromis entre \emph{surface d'écoulement} (proportionnelle à $x$) et \emph{charge hydraulique motrice} (proportionnelle à $\sqrt{H - x/2}$). C'est un exemple classique d'\textbf{optimum technico-économique} enseigné en Génie Rural.
\end{savaistu}

\begin{popfigure}
\begin{tikzpicture}[scale=0.75, domain=0:40]
\begin{axis}[
    width=\linewidth, height=10cm,
    axis lines=middle,
    xlabel={$x$ (cm)}, ylabel={$Q(x)$ (L/h)},
    xmin=0, xmax=42, ymin=0, ymax=160,
    xtick={0,5.4,26.7,39.4,40},
    xticklabels={$0$, $x_1\approx5{,}4$, $x_{\text{opt}}\approx26{,}7$, $x_2\approx39{,}4$, $40$},
    ytick={0,50,153},
    yticklabels={$0$, $50$, $Q_{\text{max}}\approx153$},
    grid=major, grid style={dashed, gray!30},
    samples=200,
    clip=false
]
\addplot[popcurve, thick, domain=0:40] {0.05*x*sqrt(max(0,1962*(20 - x/2)))};
\addplot[dashed, poporange, thick] coordinates {(0,50) (40,50)} node[pos=0.5, above, sloped] {Seuil 50 L/h};
\addplot[only marks, mark=*, mark size=3, popblue] coordinates {(0,0) (5.43,50) (26.67,153) (39.35,50) (40,0)};
\draw[popgreen, thick, <->] (5.43, -5) -- (39.35, -5) node[midway, below] {Zone admissible $Q \ge 50$};
\end{axis}
\end{tikzpicture}
\legende{Représentation graphique du débit $Q(x)$. La zone verte sous la courbe et au-dessus du seuil 50 L/h correspond aux largeurs de fente admissibles.}
\end{popfigure}

\newpage

\coursec{Méthodes et exercices résolus}

\coursec{Fonctions numériques d'une variable réelle}

\courssub{Continuité sur un intervalle : définitions et propriétés fondamentales}

\begin{definition}[Continuité en un point et sur un intervalle]
Soit $f$ une fonction définie sur un intervalle $I$ et $a \in I$.
\begin{itemize}
\item $f$ est \cle{continue en $a$} si $\displaystyle\lim_{x \to a} f(x) = f(a)$.
\item $f$ est \cle{continue sur $I$} si elle est continue en tout point de $I$.
\end{itemize}
\end{definition}

\begin{propriete}[Fonctions usuelles continues sur leur ensemble de définition]
Les fonctions suivantes sont continues sur leur ensemble de définition :
\begin{itemize}
\item Les fonctions \cle{polynômes} (sur $\mathbb{R}$).
\item Les fonctions \cle{rationnelles} (sur chaque intervalle de $D_f$).
\item Les fonctions \cle{trigonométriques} $\sin$, $\cos$ (sur $\mathbb{R}$) ; $\tan$ (sur $D_{\tan}$).
\item La fonction \cle{racine carrée} $x \mapsto \sqrt{x}$ (sur $[0\,;+\infty[$).
\item La fonction \cle{valeur absolue} $x \mapsto |x|$ (sur $\mathbb{R}$).
\end{itemize}
\end{propriete}

\begin{propriete}[Opérations conservant la continuité]
Si $f$ et $g$ sont continues sur un intervalle $I$, alors :
\begin{itemize}
\item $f+g$, $f-g$, $f \times g$ sont continues sur $I$.
\item $\dfrac{f}{g}$ est continue sur $I$ si $g(x) \neq 0$ sur $I$.
\item Si $f(I) \subset J$ et $g$ continue sur $J$, la \cle{composée} $g \circ f$ est continue sur $I$.
\end{itemize}
\end{propriete}

\begin{methode}[Prouver qu'une fonction est continue sur un intervalle]
\begin{enumerate}
\item \cle{Identifier} la nature de la fonction : polynôme, fonction rationnelle, fonction trigonométrique, fonction irrationnelle (racine carrée), ou somme/produit/quotient/composée de telles fonctions.
\item \cle{Déterminer l'ensemble de définition} $D_f$ : dénominateur non nul pour un quotient, quantité sous la racine positive ou nulle pour $\sqrt{u(x)}$.
\item \cle{Appliquer les théorèmes généraux} : les fonctions polynômes, sinus, cosinus sont continues sur $\mathbb{R}$ ; une fonction rationnelle est continue sur chaque intervalle de son ensemble de définition ; $x \mapsto \sqrt{x}$ est continue sur $[0\,;+\infty[$.
\item \cle{Conclure} par une phrase rédigée : « $f$ est continue sur $I$ comme somme (ou produit, quotient, composée) de fonctions continues sur $I$ ».
\end{enumerate}
\textbf{Réflexe Probatoire :} ne jamais écrire « $f$ est continue » sans préciser \emph{sur quel intervalle} et \emph{pourquoi}.
\end{methode}

\begin{exoresolu}[Continuité d'une fonction irrationnelle]
On considère la fonction $f$ définie par $f(x)=\sqrt{x-2}+\dfrac{1}{x-5}$.
\begin{enumerate}
\item Déterminer l'ensemble de définition $D_f$ de $f$.
\item Justifier que $f$ est continue sur l'intervalle $]5\,;+\infty[$.
\end{enumerate}
\tcblower
\textbf{1. Ensemble de définition.}

La fonction $f$ est définie lorsque les deux conditions suivantes sont réunies :
\begin{itemize}
\item la quantité sous la racine est positive ou nulle : $x-2 \geq 0$, c'est-à-dire $x \geq 2$ ;
\item le dénominateur est non nul : $x-5 \neq 0$, c'est-à-dire $x \neq 5$.
\end{itemize}
Donc $D_f=[2\,;5[\,\cup\,]5\,;+\infty[$.

\textbf{2. Continuité sur $]5\,;+\infty[$.}

\begin{itemize}
\item La fonction $x \mapsto x-2$ est une fonction polynôme, donc continue sur $\mathbb{R}$, et elle est strictement positive sur $]5\,;+\infty[$. La fonction $x \mapsto \sqrt{x}$ est continue sur $[0\,;+\infty[$. Par composition, $x \mapsto \sqrt{x-2}$ est continue sur $]5\,;+\infty[$.
\item La fonction $x \mapsto \dfrac{1}{x-5}$ est une fonction rationnelle dont le dénominateur ne s'annule pas sur $]5\,;+\infty[$ : elle y est donc continue.
\end{itemize}
\textbf{Conclusion :} $f$ est continue sur $]5\,;+\infty[$ comme somme de deux fonctions continues sur cet intervalle.
\end{exoresolu}

\courssub{Fonction continue et strictement monotone : Bijection et fonction réciproque}

\begin{propriete}[Théorème de la bijection (Théorème des valeurs intermédiaires - TVI)]
Soit $f$ une fonction \cle{continue} et \cle{strictement monotone} sur un intervalle $I$.
Alors $f$ réalise une \cle{bijection} de $I$ sur l'intervalle image $J = f(I)$.
\begin{itemize}
\item Si $f$ est strictement croissante : $J = [\lim_{x \to \inf I} f(x) \,;\, \lim_{x \to \sup I} f(x)]$ (bornes éventuellement infinies, crochets adaptés).
\item Si $f$ est strictement décroissante : $J = [\lim_{x \to \sup I} f(x) \,;\, \lim_{x \to \inf I} f(x)]$.
\end{itemize}
\emph{Conséquence :} Pour tout $k \in J$, l'équation $f(x)=k$ admet une \textbf{unique} solution dans $I$.
\end{propriete}

\begin{propriete}[Fonction réciproque]
Si $f$ est une bijection de $I$ sur $J$, sa \cle{fonction réciproque} $f^{-1}$ est définie de $J$ vers $I$ par :
\[ y = f(x) \iff x = f^{-1}(y). \]
$f^{-1}$ est continue et strictement monotone sur $J$ (même sens de variation que $f$). Les courbes $\mathcal{C}_f$ et $\mathcal{C}_{f^{-1}}$ sont symétriques par rapport à la droite $y=x$.
\end{propriete}

\begin{methode}[Montrer qu'une équation admet une unique solution et l'encadrer]
Pour résoudre $f(x)=k$ sur un intervalle $I$ :
\begin{enumerate}
\item \cle{Justifier la continuité} de $f$ sur $I$.
\item \cle{Étudier la monotonie} : calculer $f'(x)$, dresser le tableau de signes de $f'$ et montrer que $f$ est \emph{strictement} monotone sur $I$.
\item \cle{Calculer les limites (ou valeurs) aux bornes} de $I$ et vérifier que $k$ appartient à l'intervalle image $f(I)$.
\item \cle{Conclure par le théorème de la bijection} : $f$ réalise une bijection de $I$ sur $f(I)$, donc l'équation $f(x)=k$ admet une \emph{unique} solution $\alpha$ dans $I$.
\item \cle{Encadrer $\alpha$} : calculer $f(a)$ et $f(b)$ pour des valeurs bien choisies, puis conclure par stricte monotonie : $f(a)<k<f(b) \Rightarrow a<\alpha<b$ (si $f$ croissante).
\end{enumerate}
\textbf{Réflexe :} pour un encadrement à $10^{-1}$, tester les dixièmes successifs ($f(1{,}2)$, $f(1{,}3)$, ...).
\end{methode}

\begin{exoresolu}[Existence et unicité d'une solution]
Soit $f$ la fonction définie sur $\mathbb{R}$ par $f(x)=x^3+3x-1$.
\begin{enumerate}
\item Étudier les variations de $f$ sur $\mathbb{R}$.
\item Montrer que l'équation $f(x)=0$ admet une unique solution $\alpha$ dans $\mathbb{R}$.
\item Donner un encadrement de $\alpha$ d'amplitude $10^{-1}$.
\end{enumerate}
\tcblower
\textbf{1. Variations de $f$.}

$f$ est une fonction polynôme, donc dérivable sur $\mathbb{R}$. Pour tout réel $x$ :
\[ f'(x)=3x^2+3=3(x^2+1). \]
Pour tout $x\in\mathbb{R}$, $x^2+1>0$, donc $f'(x)>0$. La fonction $f$ est donc \cle{strictement croissante} sur $\mathbb{R}$.

\textbf{2. Existence et unicité de la solution.}

\begin{itemize}
\item $f$ est un polynôme : elle est \cle{continue} sur $\mathbb{R}$.
\item $f$ est \cle{strictement croissante} sur $\mathbb{R}$ d'après la question 1.
\item Limites aux bornes : $\displaystyle\lim_{x\to -\infty}f(x)=-\infty$ (terme dominant $x^3$) et $\displaystyle\lim_{x\to +\infty}f(x)=+\infty$.
\end{itemize}
Donc $f$ réalise une bijection de $\mathbb{R}$ sur $]-\infty\,;+\infty[=\mathbb{R}$. Comme $0\in\mathbb{R}$, le théorème de la bijection assure que l'équation $f(x)=0$ admet une \textbf{unique solution} $\alpha$ dans $\mathbb{R}$.

\textbf{3. Encadrement de $\alpha$ à $10^{-1}$ près.}

On calcule : $f(0)=-1<0$ et $f(1)=1+3-1=3>0$, donc $0<\alpha<1$.

On affine par dixièmes :
\[ f(0{,}3)=0{,}027+0{,}9-1=-0{,}073<0 \qquad ; \qquad f(0{,}4)=0{,}064+1{,}2-1=0{,}264>0. \]
On a donc $f(0{,}3)<0<f(0{,}4)$. Comme $f$ est strictement croissante, on en déduit :
\[ 0{,}3<\alpha<0{,}4. \]
\textbf{Conclusion :} l'équation $x^3+3x-1=0$ admet une unique solution $\alpha$, avec $\alpha\in\,]0{,}3\,;0{,}4[$.
\end{exoresolu}

\courssub{Limites des fonctions trigonométriques : limites remarquables et encadrement}

\begin{propriete}[Limites remarquables trigonométriques (en radians)]
\begin{align*}
\lim_{x \to 0} \frac{\sin x}{x} &= 1, \\
\lim_{x \to 0} \frac{\tan x}{x} &= 1, \\
\lim_{x \to 0} \frac{1 - \cos x}{x^2} &= \frac{1}{2}.
\end{align*}
\emph{Preuve de la 3e :} $1-\cos x = 2\sin^2(\frac{x}{2})$, donc $\frac{1-\cos x}{x^2} = \frac{1}{2}\left(\frac{\sin(x/2)}{x/2}\right)^2 \to \frac{1}{2}$.
\end{propriete}

\begin{propriete}[Théorème des gendarmes (Encadrement)]
Si $g(x) \leq f(x) \leq h(x)$ au voisinage de $a$ (ou en $\pm\infty$) et $\lim g = \lim h = L$, alors $\lim f = L$.
\emph{Application classique :} $-1 \leq \sin x \leq 1 \Rightarrow -\frac{1}{x} \leq \frac{\sin x}{x} \leq \frac{1}{x}$ pour $x>0$, d'où $\lim_{+\infty} \frac{\sin x}{x} = 0$.
\end{propriete}

\begin{methode}[Lever une indétermination avec les fonctions trigonométriques]
\begin{enumerate}
\item \cle{Repérer la forme indéterminée} (souvent $\dfrac{0}{0}$ en $0$).
\item \cle{Se ramener aux limites de référence} :
\[ \lim_{x\to 0}\frac{\sin x}{x}=1 \qquad ; \qquad \lim_{x\to 0}\frac{\tan x}{x}=1 \qquad ; \qquad \lim_{x\to 0}\frac{1-\cos x}{x^2}=\frac{1}{2}. \]
\item \cle{Utiliser le changement de variable} : si $x\to 0$, alors $\dfrac{\sin(ax)}{bx}=\dfrac{a}{b}\times\dfrac{\sin(ax)}{ax}$ avec $u=ax\to 0$.
\item \cle{Encadrer si nécessaire} : pour tout réel $x$, $-1\leq\sin x\leq 1$ et $-1\leq\cos x\leq 1$ ; on applique ensuite le théorème des gendarmes (utile pour les limites en l'infini).
\end{enumerate}
\textbf{Attention :} dans ces formules, $x$ est en \emph{radians}.
\end{methode}

\begin{exoresolu}[Limites trigonométriques]
Calculer les limites suivantes :
\[ \text{a) } \lim_{x\to 0}\frac{\sin(3x)}{2x} \qquad ; \qquad \text{b) } \lim_{x\to 0}\frac{1-\cos x}{x^2} \qquad ; \qquad \text{c) } \lim_{x\to +\infty}\frac{\sin x}{x}. \]
\tcblower
\textbf{a)} On fait apparaître la limite de référence en posant $u=3x$. Quand $x\to 0$, $u\to 0$ :
\[ \frac{\sin(3x)}{2x}=\frac{3}{2}\times\frac{\sin(3x)}{3x}=\frac{3}{2}\times\frac{\sin u}{u}. \]
Or $\displaystyle\lim_{u\to 0}\frac{\sin u}{u}=1$, donc $\displaystyle\lim_{x\to 0}\frac{\sin(3x)}{2x}=\frac{3}{2}\times 1=\frac{3}{2}$.

\textbf{b)} C'est une limite de référence du cours : $\displaystyle\lim_{x\to 0}\frac{1-\cos x}{x^2}=\frac{1}{2}$.

\emph{Justification par les formules :} $1-\cos x=2\sin^2\!\left(\dfrac{x}{2}\right)$, donc
\[ \frac{1-\cos x}{x^2}=\frac{2\sin^2(x/2)}{x^2}=\frac{1}{2}\left(\frac{\sin(x/2)}{x/2}\right)^2 \xrightarrow[x\to 0]{}\frac{1}{2}\times 1^2=\frac{1}{2}. \]

\textbf{c)} Pour tout $x>0$, on a $-1\leq\sin x\leq 1$, donc en divisant par $x>0$ :
\[ -\frac{1}{x}\leq\frac{\sin x}{x}\leq\frac{1}{x}. \]
Or $\displaystyle\lim_{x\to +\infty}\left(-\frac{1}{x}\right)=\lim_{x\to +\infty}\frac{1}{x}=0$. D'après le \cle{théorème des gendarmes} :
\[ \lim_{x\to +\infty}\frac{\sin x}{x}=0. \]
\textbf{Conclusion :} a) $\dfrac{3}{2}$ ; b) $\dfrac{1}{2}$ ; c) $0$.
\end{exoresolu}

\courssub{Limites des fonctions irrationnelles : rationalisation et composition}

\begin{propriete}[Rationalisation d'une différence de racines]
Pour $a \geq 0, b \geq 0$ :
\[ \sqrt{a} - \sqrt{b} = \frac{a-b}{\sqrt{a}+\sqrt{b}}. \]
Cette identité permet de lever les indéterminations de type $\dfrac{0}{0}$ ou $\infty - \infty$.
\end{propriete}

\begin{propriete}[Limites en l'infini avec racines]
Pour $x > 0$, $\sqrt{x^2} = x$. Pour $x < 0$, $\sqrt{x^2} = -x = |x|$.
On factorise par le terme dominant : $\sqrt{ax^2+bx+c} \sim_{+\infty} \sqrt{a}\,x$ (si $a>0$).
\end{propriete}

\begin{methode}[Lever une indétermination avec des racines carrées]
\begin{enumerate}
\item \cle{Constater l'indétermination} : forme $\dfrac{0}{0}$ ou $\infty-\infty$.
\item \cle{Rationaliser} : multiplier numérateur et dénominateur par l'\emph{expression conjuguée} :
\[ \sqrt{a}-\sqrt{b}=\frac{a-b}{\sqrt{a}+\sqrt{b}}. \]
\item \cle{Simplifier} puis conclure avec les opérations sur les limites.
\item \cle{En l'infini}, on peut aussi factoriser par le terme dominant : pour $x>0$, $\sqrt{x^2}=x$ (attention : $\sqrt{x^2}=|x|$ en général).
\end{enumerate}
\textbf{Réflexe :} dès qu'une différence de racines carrées apparaît, penser \emph{immédiatement} à la quantité conjuguée.
\end{methode}

\begin{exoresolu}[Rationalisation et limite en l'infini]
\begin{enumerate}
\item Calculer $\displaystyle\lim_{x\to 0}\frac{\sqrt{x+1}-1}{x}$.
\item Calculer $\displaystyle\lim_{x\to +\infty}\left(\sqrt{x^2+3x}-x\right)$.
\end{enumerate}
\tcblower
\textbf{1.} Quand $x\to 0$, le numérateur tend vers $\sqrt{1}-1=0$ et le dénominateur vers $0$ : forme indéterminée $\dfrac{0}{0}$. On multiplie par la quantité conjuguée $\sqrt{x+1}+1$ :
\begin{align*}
\frac{\sqrt{x+1}-1}{x} &= \frac{(\sqrt{x+1}-1)(\sqrt{x+1}+1)}{x(\sqrt{x+1}+1)} \\
&= \frac{(x+1)-1}{x(\sqrt{x+1}+1)} = \frac{x}{x(\sqrt{x+1}+1)} = \frac{1}{\sqrt{x+1}+1}.
\end{align*}
Cette dernière expression est continue en $0$, donc :
\[ \lim_{x\to 0}\frac{\sqrt{x+1}-1}{x}=\frac{1}{\sqrt{1}+1}=\frac{1}{2}. \]

\textbf{2.} Quand $x\to+\infty$, $\sqrt{x^2+3x}\to+\infty$ et $x\to+\infty$ : forme indéterminée $\infty-\infty$. On rationalise :
\begin{align*}
\sqrt{x^2+3x}-x &= \frac{(\sqrt{x^2+3x}-x)(\sqrt{x^2+3x}+x)}{\sqrt{x^2+3x}+x} = \frac{x^2+3x-x^2}{\sqrt{x^2+3x}+x} = \frac{3x}{\sqrt{x^2+3x}+x}.
\end{align*}
Pour $x>0$, $\sqrt{x^2+3x}=x\sqrt{1+\dfrac{3}{x}}$ (car $\sqrt{x^2}=x$ pour $x>0$). Donc :
\[ \frac{3x}{\sqrt{x^2+3x}+x}=\frac{3x}{x\left(\sqrt{1+\frac{3}{x}}+1\right)}=\frac{3}{\sqrt{1+\frac{3}{x}}+1}. \]
Comme $\dfrac{3}{x}\to 0$ quand $x\to+\infty$, on obtient :
\[ \lim_{x\to +\infty}\left(\sqrt{x^2+3x}-x\right)=\frac{3}{\sqrt{1}+1}=\frac{3}{2}. \]
\textbf{Conclusion :} 1. $\dfrac{1}{2}$ ; 2. $\dfrac{3}{2}$.
\end{exoresolu}

\courssub{Représentation graphique : méthodologie complète et lecture graphique}

\begin{methode}[Étudier et représenter une fonction]
\begin{enumerate}
\item \cle{Ensemble de définition} $D_f$ et, le cas échéant, parité ou périodicité (pour réduire le domaine d'étude).
\item \cle{Limites aux bornes} de $D_f$ ; en déduire les \emph{asymptotes} (verticales si limite infinie en un point, horizontales si limite finie en l'infini, obliques si limite infinie avec pente finie).
\item \cle{Dérivée} $f'$, étude de son signe, tableau de variations complet (avec limites et valeurs remarquables).
\item \cle{Points particuliers} : intersections avec les axes, tangentes horizontales (là où $f'=0$), centre de symétrie (si impaire).
\item \cle{Tracé} : placer d'abord les asymptotes (trait pointillé) et les points remarquables, puis esquisser la courbe en respectant le tableau de variations.
\item \cle{Lecture graphique} : pour résoudre graphiquement $f(x)=k$, on lit les abscisses des points d'intersection de la courbe avec la droite $y=k$.
\end{enumerate}
\end{methode}

\begin{aretenir}[Asymptotes]
\begin{itemize}
\item \textbf{Verticale} $x=a$ : $\lim_{x\to a^{\pm}} f(x) = \pm\infty$.
\item \textbf{Horizontale} $y=\ell$ : $\lim_{x\to \pm\infty} f(x) = \ell$ (fini).
\item \textbf{Oblique} $y=ax+b$ : $\lim_{x\to \pm\infty} \frac{f(x)}{x}=a$ et $\lim_{x\to \pm\infty} (f(x)-ax)=b$.
\end{itemize}
\end{aretenir}

\begin{exoresolu}[Étude complète et tracé]
Soit $f$ la fonction définie par $f(x)=\dfrac{2x}{x^2+1}$ et $\mathcal{C}_f$ sa courbe dans un repère orthonormé.
\begin{enumerate}
\item Donner $D_f$ et montrer que $f$ est impaire.
\item Calculer les limites de $f$ en $-\infty$ et $+\infty$. Interpréter graphiquement.
\item Étudier les variations de $f$ et dresser son tableau de variations.
\item Tracer l'allure de $\mathcal{C}_f$.
\end{enumerate}
\tcblower
\textbf{1.} Pour tout réel $x$, $x^2+1\geq 1>0$ : le dénominateur ne s'annule jamais, donc $D_f=\mathbb{R}$.

Pour tout $x\in\mathbb{R}$ : $f(-x)=\dfrac{2(-x)}{(-x)^2+1}=-\dfrac{2x}{x^2+1}=-f(x)$. Donc $f$ est \cle{impaire} : sa courbe est symétrique par rapport à l'origine, et on peut limiter l'étude à $[0\,;+\infty[$.

\textbf{2.} Pour $x\neq 0$, $f(x)=\dfrac{2x}{x^2\left(1+\frac{1}{x^2}\right)}=\dfrac{2}{x\left(1+\frac{1}{x^2}\right)}$. Quand $x\to\pm\infty$, $\dfrac{1}{x^2}\to 0$ et $\dfrac{2}{x}\to 0$, donc :
\[ \lim_{x\to -\infty}f(x)=0 \qquad \text{et} \qquad \lim_{x\to +\infty}f(x)=0. \]
\emph{Interprétation :} la droite $y=0$ (l'axe des abscisses) est \cle{asymptote horizontale} à $\mathcal{C}_f$ en $-\infty$ et en $+\infty$.

\textbf{3.} $f$ est une fonction rationnelle définie sur $\mathbb{R}$, donc dérivable sur $\mathbb{R}$. Avec la formule $\left(\dfrac{u}{v}\right)'=\dfrac{u'v-uv'}{v^2}$ :
\[ f'(x)=\frac{2(x^2+1)-2x(2x)}{(x^2+1)^2}=\frac{2-2x^2}{(x^2+1)^2}=\frac{2(1-x)(1+x)}{(x^2+1)^2}. \]
Le dénominateur est strictement positif, donc $f'(x)$ a le signe de $1-x^2$ : $f'(x)>0$ sur $]-1\,;1[$ et $f'(x)<0$ sur $]-\infty\,;-1[$ et $]1\,;+\infty[$.

Valeurs remarquables : $f(-1)=\dfrac{-2}{2}=-1$, $f(0)=0$, $f(1)=\dfrac{2}{2}=1$.

\begin{center}
\begin{tabular}{|c|ccccc|}\hline
$x$ & $-\infty$ & & $-1$ & & $0$ \\ \hline
$f'(x)$ & & $-$ & $0$ & $+$ & \\ \hline
$f$ & $0$ & $\searrow$ & $-1$ & $\nearrow$ & $0$ \\ \hline
\end{tabular}
\quad
\begin{tabular}{|c|cccc|}\hline
$x$ & $0$ & & $1$ & $+\infty$ \\ \hline
$f'(x)$ & & $+$ & $0$ & $-$ \\ \hline
$f$ & $0$ & $\nearrow$ & $1$ & $\searrow$ \quad $0$ \\ \hline
\end{tabular}
\end{center}

\textbf{4.} Tracé : on place l'asymptote $y=0$, les extrema $(-1\,;-1)$ et $(1\,;1)$ (tangentes horizontales), le point $(0\,;0)$.

\begin{popfigure}
\begin{tikzpicture}
\begin{axis}[width=0.9\linewidth,height=6cm,axis lines=middle,xlabel=$x$,ylabel=$y$,xmin=-5,xmax=5,ymin=-1.6,ymax=1.6,xtick={-1,1},ytick={-1,1},grid=both,grid style={popline!40}]
\addplot[popcurve,domain=-5:5,samples=200]{2*x/(x^2+1)};
\end{axis}
\end{tikzpicture}
\legende{Courbe représentative de $f(x)=\dfrac{2x}{x^2+1}$ : impaire, asymptote $y=0$, extrema en $-1$ et $1$.}
\end{popfigure}
\textbf{Conclusion :} $f$ est impaire, admet un minimum $-1$ en $x=-1$, un maximum $1$ en $x=1$, et l'axe des abscisses pour asymptote horizontale aux deux infinis.
\end{exoresolu}

\courssub{Modélisation et résolution de problèmes concrets (Synthèse)}

\begin{savaistu}[Le seuil de rentabilité dans l'industrie camerounaise]
Dans la gestion d'un atelier (menuiserie, métallerie, agro-alimentaire), la fonction bénéfice $B(x)$ est souvent un polynôme de degré 2 ou 3. La recherche du seuil de rentabilité ($B(x)=0$) et de l'optimum de production ($B'(x)=0$) est un classique de l'analyse fonctionnelle appliquée. Au Probatoire, on attend une rédaction rigoureuse : hypothèses du TVI vérifiées explicitement, conclusion dans le contexte avec unité (FCFA, m, kg, h).
\end{savaistu}

\begin{methode}[Résoudre un problème concret par les fonctions]
\begin{enumerate}
\item \cle{Traduire} l'énoncé : identifier la grandeur étudiée, la variable, et poser la fonction modèle $f$ avec son domaine (souvent un intervalle borné lié au contexte : temps, longueur, quantité).
\item \cle{Exploiter les outils du cours} : continuité et théorème de la bijection pour prouver qu'un seuil est atteint une seule fois ; dérivée et variations pour optimiser (maximum de bénéfice, minimum de coût).
\item \cle{Rédiger la démarche} : citer le théorème utilisé, vérifier explicitement ses hypothèses (continuité, stricte monotonie, appartenance de la valeur à l'intervalle image).
\item \cle{Conclure dans le contexte} : répondre par une phrase en langage courant, avec l'unité adaptée (FCFA, mètres, heures...).
\end{enumerate}
\end{methode}

\begin{exoresolu}[Seuil de rentabilité d'un atelier à Douala]
Un atelier de menuiserie produit des chaises. Pour $x$ dizaines de chaises produites par semaine ($x\in[0\,;10]$), le bénéfice hebdomadaire en milliers de FCFA est modélisé par :
\[ B(x)=-x^3+6x^2+9x-20. \]
\begin{enumerate}
\item Étudier les variations de $B$ sur $[0\,;10]$.
\item Montrer que l'équation $B(x)=0$ admet une unique solution $\alpha$ dans $[0\,;4]$, puis encadrer $\alpha$ entre deux entiers consécutifs.
\item Interpréter le résultat pour l'atelier.
\end{enumerate}
\tcblower
\textbf{1. Variations de $B$.}

$B$ est une fonction polynôme, donc dérivable sur $[0\,;10]$ :
\[ B'(x)=-3x^2+12x+9=-3(x^2-4x-3). \]
Le trinôme $x^2-4x-3$ a pour discriminant $\Delta=16+12=28$, donc pour racines $x_1=\dfrac{4-\sqrt{28}}{2}=2-\sqrt{7}\approx -0{,}65$ et $x_2=2+\sqrt{7}\approx 4{,}65$.

Sur $[0\,;10]$, $B'(x)$ a le signe opposé à celui du trinôme (coefficient $-3<0$) : $B'(x)>0$ sur $[0\,;2+\sqrt{7}[$ et $B'(x)<0$ sur $]2+\sqrt{7}\,;10]$.

Donc $B$ est \cle{strictement croissante} sur $[0\,;2+\sqrt{7}]$ et strictement décroissante sur $[2+\sqrt{7}\,;10]$, avec un maximum en $x=2+\sqrt{7}$.

\textbf{2. Existence et unicité de $\alpha$ dans $[0\,;4]$.}

Sur $[0\,;4]$, la fonction $B$ est strictement croissante (car $4<2+\sqrt{7}$). De plus :
\begin{itemize}
\item $B$ est un polynôme : elle est \cle{continue} sur $[0\,;4]$ ;
\item $B(0)=-20<0$ et $B(4)=-64+96+36-20=48>0$.
\end{itemize}
$B$ est continue et strictement monotone sur $[0\,;4]$, et $0$ est compris entre $B(0)=-20$ et $B(4)=48$. D'après le \cle{théorème de la bijection}, l'équation $B(x)=0$ admet une \textbf{unique solution} $\alpha$ dans $[0\,;4]$.

Encadrement : $B(1)=-1+6+9-20=-6<0$ et $B(2)=-8+24+18-20=14>0$. Comme $B$ est strictement croissante :
\[ 1<\alpha<2. \]

\textbf{3. Interprétation.}

L'atelier devient rentable (bénéfice positif) à partir d'environ $\alpha$ dizaines de chaises, c'est-à-dire entre 10 et 20 chaises par semaine. En dessous de ce seuil, l'atelier fonctionne à perte ; au-delà, il dégage un bénéfice, maximal pour environ $2+\sqrt{7}\approx 4{,}6$ dizaines, soit 46 chaises par semaine.
\end{exoresolu}

\begin{exercice}[Synthèse : Étude d'une fonction rationnelle et application]
Soit la fonction $f$ définie sur $\mathbb{R}^*_+$ par $f(x) = x + \dfrac{4}{x}$.
\begin{enumerate}
\item Étudier les variations de $f$ sur $\mathbb{R}^*_+$.
\item Montrer que pour tout $x > 0$, $f(x) \geq 4$. Quand l'égalité a-t-elle lieu ?
\item On considère un rectangle de périmètre $P = 20$ m. Exprimer son aire $A$ en fonction de sa longueur $L$. Retrouver le résultat de la question 2 pour déterminer les dimensions du rectangle de plus grande aire.
\end{enumerate}
\end{exercice}

\begin{corrige}[Exercice n° Synthèse]
\textbf{1.} $f'(x) = 1 - \dfrac{4}{x^2} = \dfrac{x^2-4}{x^2} = \dfrac{(x-2)(x+2)}{x^2}$.
Sur $\mathbb{R}^*_+$, le dénominateur $x^2 > 0$. Le numérateur est nul pour $x=2$, négatif sur $]0;2[$, positif sur $]2;+\infty[$.
Tableau de variations :
\begin{center}
\begin{tabular}{|c|ccccc|}\hline
$x$ & $0$ & $2$ & $+\infty$ \\ \hline
$f'(x)$ & & $-$ & $0$ & $+$ \\ \hline
$f$ & $+\infty$ & $\searrow$ & $4$ & $\nearrow$ & $+\infty$ \\ \hline
\end{tabular}
\end{center}

\textbf{2.} D'après le tableau, $f$ admet un minimum global en $x=2$ avec $f(2)=4$. Donc $\forall x>0,\ f(x) \geq 4$. Égalité pour $x=2$ uniquement.

\textbf{3.} Périmètre $P = 2(L+l) = 20 \Rightarrow l = 10 - L$.
Aire $A(L) = L \times l = L(10-L) = 10L - L^2 = -(L^2 - 10L) = -(L-5)^2 + 25$.
On peut aussi écrire $A(L) = 25 - (L-5)^2 \leq 25$. Maximum $25$ m$^2$ pour $L=5$ m (carré).
\emph{Lien avec q.2 :} $A(L) = 25 - (L-5)^2$. Posons $x = L/2$ (ou autre liaison). Ici, on note que $f(x)=x+4/x$ n'est pas directement l'aire mais l'étude de $x+4/x \geq 4$ est l'inégalité AM-GM qui sous-tend l'optimalité du carré pour un périmètre fixé. Pour $P=20$, $L+l=10$, aire max quand $L=l=5$.
\end{corrige}

\begin{attention}[Les 5 erreurs qui coûtent des points au Probatoire]
\begin{enumerate}
\item \textbf{Appliquer le théorème de la bijection sans vérifier les hypothèses.} Il faut rédiger explicitement : « $f$ est \emph{continue} et \emph{strictement monotone} sur $I$, et $k$ appartient à l'intervalle image ». Oublier la stricte monotonie ou la continuité fait perdre les points de justification.
\item \textbf{Confondre $\sqrt{x^2}$ avec $x$.} Pour $x<0$, $\sqrt{x^2}=-x=|x|$. Dans les limites en $-\infty$ avec des racines, cette erreur change le signe du résultat.
\item \textbf{Oublier la quantité conjuguée.} Devant $\sqrt{a}-\sqrt{b}$, diviser ou « simplifier » les racines séparément est faux : il faut multiplier par $\sqrt{a}+\sqrt{b}$ en haut \emph{et} en bas.
\item \textbf{Mal recopier les limites de référence trigonométriques.} $\lim\limits_{x\to 0}\dfrac{\sin x}{x}=1$ mais $\lim\limits_{x\to 0}\dfrac{1-\cos x}{x^2}=\dfrac{1}{2}$ (et non $1$). De plus, ces formules exigent des angles en \emph{radians} et un argument qui tend vers $0$.
\item \textbf{Conclure hors du contexte.} Dans un problème concret, une réponse « $\alpha\approx 1{,}4$ » sans phrase d'interprétation (unité, signification pour l'atelier ou le chantier) est incomplète : le correcteur attend une conclusion rédigée en langage courant.
\end{enumerate}
\end{attention}

\newpage

\coursec{Exercices}

\courssub{Je vérifie mes connaissances}

\begin{exercice}[QCM : continuité et limites]
Pour chaque question, une seule réponse est exacte. Entoure-la et justifie brièvement.
\begin{enumerate}
\item Si $f$ est continue sur $[a\,;b]$ et strictement monotone, alors l'équation $f(x)=k$ avec $k$ compris entre $f(a)$ et $f(b)$ :
\begin{enumerate}
\item n'a pas de solution ;
\item a une unique solution dans $[a\,;b]$ ;
\item a au moins deux solutions.
\end{enumerate}
\item $\displaystyle\lim_{x\to 0}\dfrac{\sin x}{x}$ vaut :
\begin{enumerate}
\item $0$ ; \item $1$ ; \item $+\infty$.
\end{enumerate}
\item Si $f$ est une bijection de $I$ sur $J$ et $f^{-1}$ sa réciproque, alors pour $a\in I$ avec $f'(a)\neq 0$ :
\begin{enumerate}
\item $(f^{-1})'(f(a))=\dfrac{1}{f'(a)}$ ;
\item $(f^{-1})'(f(a))=f'(a)$ ;
\item $(f^{-1})'(f(a))=-\dfrac{1}{f'(a)}$.
\end{enumerate}
\item $\displaystyle\lim_{x\to +\infty}\left(\sqrt{x^2+x}-x\right)$ vaut :
\begin{enumerate}
\item $+\infty$ ; \item $0$ ; \item $\dfrac{1}{2}$.
\end{enumerate}
\end{enumerate}
\end{exercice}

\begin{exercice}[Vrai ou faux]
Réponds par Vrai ou Faux en justifiant chaque réponse.
\begin{enumerate}
\item Toute fonction dérivable sur un intervalle $I$ est continue sur $I$.
\item Toute fonction continue sur un intervalle $I$ est dérivable sur $I$.
\item La fonction $x\mapsto \sqrt{x}$ est continue en $0$.
\item Si $f$ est continue et strictement croissante sur $\mathbb{R}$, alors sa courbe coupe toute droite horizontale en exactement un point.
\end{enumerate}
\end{exercice}

\begin{exercice}[Définitions]
\begin{enumerate}
\item Donner la définition de la continuité d'une fonction $f$ en un point $x_0$.
\item Énoncer le théorème des valeurs intermédiaires (TVI).
\item Qu'appelle-t-on fonction réciproque d'une bijection $f$ ? Quelle relation lie les courbes de $f$ et de $f^{-1}$ dans un repère orthonormé ?
\end{enumerate}
\end{exercice}

\begin{exercice}[Calculs directs de limites]
Calculer les limites suivantes en justifiant chaque étape :
\begin{enumerate}
\item $\displaystyle\lim_{x\to 0}\dfrac{\sin(3x)}{x}$ ;
\item $\displaystyle\lim_{x\to 0}\dfrac{\tan x}{x}$ ;
\item $\displaystyle\lim_{x\to +\infty}\left(\sqrt{x^2+4}-x\right)$ ;
\item $\displaystyle\lim_{x\to 4}\dfrac{\sqrt{x}-2}{x-4}$.
\end{enumerate}
\end{exercice}

\courssub{Je m'entraîne}

\begin{exercice}[Application du TVI]
Soit $f$ la fonction définie sur $\mathbb{R}$ par $f(x)=x^3+x-3$.
\begin{enumerate}
\item Étudier les variations de $f$ sur $\mathbb{R}$.
\item Montrer que l'équation $f(x)=0$ admet une unique solution $\alpha$ dans $\mathbb{R}$.
\item Justifier que $\alpha\in[1\,;2]$, puis donner un encadrement de $\alpha$ à $10^{-1}$ près.
\end{enumerate}
\end{exercice}

\begin{exercice}[Limites trigonométriques]
Un technicien étudie un signal modélisé par $f(x)=\dfrac{1-\cos x}{x^2}$ pour $x\neq 0$.
\begin{enumerate}
\item En utilisant l'identité $1-\cos x=2\sin^2\!\left(\dfrac{x}{2}\right)$, calculer $\displaystyle\lim_{x\to 0}f(x)$.
\item La fonction $f$ est-elle prolongeable par continuité en $0$ ? Si oui, préciser le prolongement.
\item Calculer $\displaystyle\lim_{x\to 0}\dfrac{\sin(5x)}{\sin(2x)}$.
\end{enumerate}
\end{exercice}

\begin{exercice}[Fonction irrationnelle]
Soit $g$ la fonction définie par $g(x)=\sqrt{x^2-4x+3}$.
\begin{enumerate}
\item Déterminer l'ensemble de définition $D_g$ de $g$.
\item Calculer les limites de $g$ aux bornes de $D_g$.
\item Étudier la continuité de $g$ sur $D_g$.
\item Calculer $\displaystyle\lim_{x\to +\infty}\dfrac{g(x)}{x}$ et interpréter graphiquement le résultat.
\end{enumerate}
\end{exercice}

\begin{exercice}[Bijection et réciproque]
Soit $h$ la fonction définie sur $[0\,;+\infty[$ par $h(x)=x^2+2x$.
\begin{enumerate}
\item Montrer que $h$ réalise une bijection de $[0\,;+\infty[$ sur un intervalle $J$ à préciser.
\item Expliciter la fonction réciproque $h^{-1}$.
\item Calculer $h(1)$, puis $(h^{-1})'(3)$ de deux façons différentes.
\item Tracer dans un repère orthonormé les courbes de $h$ et de $h^{-1}$.
\end{enumerate}
\end{exercice}

\courssub{Je me prépare au Probatoire}

\begin{exercice}[Type Probatoire 1 : analyse pure]
Soit $f$ la fonction définie sur $\mathbb{R}$ par
\[f(x)=\dfrac{2x+1}{\sqrt{x^2+4}-x}.\]
On note $(\mathcal{C})$ sa courbe représentative dans un repère orthonormé $(O\,;\vec{i},\vec{j})$ d'unité \SI{1}{cm}.
\begin{enumerate}
\item \begin{enumerate}
\item Justifier que pour tout réel $x$, $\sqrt{x^2+4}-x>0$. En déduire que $f$ est définie et continue sur $\mathbb{R}$.
\item Montrer que $f(x)=(2x+1)\,\dfrac{\sqrt{x^2+4}+x}{4}$ pour tout réel $x$.
\end{enumerate}
\item \begin{enumerate}
\item Calculer $\displaystyle\lim_{x\to +\infty}f(x)$.
\item Calculer $\displaystyle\lim_{x\to -\infty}f(x)$, puis $\displaystyle\lim_{x\to -\infty}\big(f(x)-0\big)$. Interpréter graphiquement.
\end{enumerate}
\item \begin{enumerate}
\item Calculer $f'(x)$ et montrer que $f'(x)$ a le signe de $2\sqrt{x^2+4}+2x+1$ sur $\mathbb{R}$.
\item Dresser le tableau de variations complet de $f$.
\end{enumerate}
\item Montrer que l'équation $f(x)=3$ admet une unique solution $\alpha$ dans $\mathbb{R}$, et donner un encadrement de $\alpha$ à $10^{-1}$ près.
\item Tracer $(\mathcal{C})$ en plaçant les éléments remarquables (asymptote, points d'intersection avec les axes).
\end{enumerate}
\end{exercice}

\begin{exercice}[Type Probatoire 2 : trigonométrie et modélisation]
Lors de la charge d'un condensateur dans un atelier d'électronique, la tension (en volts) à ses bornes est modélisée, pour $t\in[0\,;\,0{,}1]$ (en secondes), par :
\[u(t)=10-8\,e^{-2t}+2\sin(100\pi t).\]
\begin{enumerate}
\item Calculer $u(0)$ et $u(0{,}1)$ (on donnera une valeur approchée à $10^{-2}$ près).
\item \begin{enumerate}
\item Justifier que $u$ est continue sur $[0\,;0{,}1]$.
\item Calculer $u'(t)$ pour $t\in[0\,;0{,}1]$.
\item On admet que $u$ est strictement croissante sur $[0\,;0{,}02]$. Dresser le tableau de variations de $u$ sur cet intervalle.
\end{enumerate}
\item \begin{enumerate}
\item Montrer qu'il existe un unique réel $t_0\in[0\,;0{,}02]$ tel que $u(t_0)=5$.
\item À l'aide de la calculatrice, donner un encadrement de $t_0$ à $10^{-2}$ près.
\item Interpréter ce résultat pour le technicien : à quel instant la tension atteint-elle \SI{5}{V} ?
\end{enumerate}
\item Calculer $\displaystyle\lim_{t\to 0}\dfrac{u(t)-u(0)}{t}$. Que représente ce nombre pour le circuit ?
\end{enumerate}
\end{exercice}

\begin{exercice}[Type Probatoire 3 : optimisation en atelier]
Un tôlier veut fabriquer un bac ouvert (sans couvercle) à partir d'une feuille carrée de côté \SI{40}{cm}. Il découpe à chaque coin un carré de côté $x$ (en cm), puis replie les bords.
\begin{enumerate}
\item \begin{enumerate}
\item Justifier que $x\in[0\,;20]$.
\item Montrer que le volume du bac est $V(x)=x(40-2x)^2$, puis que $V(x)=4x^3-160x^2+1600x$.
\end{enumerate}
\item \begin{enumerate}
\item Étudier la continuité et la dérivabilité de $V$ sur $[0\,;20]$.
\item Calculer $V'(x)$ et vérifier que $V'(x)=12\left(x-\dfrac{20}{3}\right)(x-20)$.
\item Dresser le tableau de variations de $V$ sur $[0\,;20]$.
\end{enumerate}
\item En déduire la valeur de $x$ pour laquelle le volume est maximal, et calculer ce volume maximal (arrondi au cm$^3$).
\item Le client souhaite un bac de volume \SI{1500}{cm^3}.
\begin{enumerate}
\item Montrer, à l'aide du théorème des valeurs intermédiaires, qu'il existe au moins une valeur de $x$ dans $\left]0\,;\dfrac{20}{3}\right[$ réalisant ce volume.
\item Donner un encadrement de cette valeur à $10^{-1}$ près.
\end{enumerate}
\item Rédiger une conclusion pour le tôlier : quelle découpe conseilles-tu, et pourquoi ?
\end{enumerate}
\end{exercice}

\newpage

\coursec{Corrigés des exercices}

\coursec{Fonctions numériques d'une variable réelle}

\courssub{Je vérifie mes connaissances}

\begin{corrige}[Exercice 1 — QCM : continuité et limites]
\begin{enumerate}
\item \textbf{Réponse b.} $f$ est continue et strictement monotone sur $[a\,;b]$ : c'est une bijection de $[a\,;b]$ sur l'intervalle image. Comme $k$ est compris entre $f(a)$ et $f(b)$, le théorème de la bijection (corollaire du TVI) garantit l'existence d'une \textbf{unique} solution.
\item \textbf{Réponse b.} C'est la limite remarquable du cours : $\displaystyle\lim_{x\to 0}\dfrac{\sin x}{x}=1$.
\item \textbf{Réponse a.} Formule de la dérivée de la réciproque : $(f^{-1})'(f(a))=\dfrac{1}{f'(a)}$, valable car $f'(a)\neq 0$.
\item \textbf{Réponse c.} On rationalise :
\[\sqrt{x^2+x}-x=\dfrac{x}{\sqrt{x^2+x}+x}=\dfrac{1}{\sqrt{1+\frac{1}{x}}+1}\xrightarrow[x\to+\infty]{}\dfrac{1}{2}.\]
\end{enumerate}
\end{corrige}

\begin{corrige}[Exercice 2 — Vrai ou faux]
\begin{enumerate}
\item \textbf{Vrai.} C'est un théorème du cours : toute fonction dérivable en un point (ou sur un intervalle) est continue en ce point (sur cet intervalle).
\item \textbf{Faux.} La réciproque est fausse. Contre-exemple : $x\mapsto |x|$ est continue en $0$ mais n'y est pas dérivable (demi-tangentes distinctes).
\item \textbf{Vrai.} $\displaystyle\lim_{x\to 0^+}\sqrt{x}=0=\sqrt{0}$ : la fonction racine carrée est continue en $0$ (à droite, ce qui est la continuité en la borne de son ensemble de définition). Attention : elle n'est pas \emph{dérivable} en $0$.
\item \textbf{Faux.} Il faut que les limites aux bornes encadrent la droite. Contre-exemple : $f(x)=e^{x}$ est continue et strictement croissante sur $\mathbb{R}$, mais sa courbe ne rencontre jamais la droite $y=-1$. La conclusion n'est vraie que si $\displaystyle\lim_{-\infty}f=-\infty$ et $\displaystyle\lim_{+\infty}f=+\infty$.
\end{enumerate}
\end{corrige}

\begin{corrige}[Exercice 3 — Définitions]
\begin{enumerate}
\item $f$ est \cle{continue en $x_0$} si $f$ est définie en $x_0$ et si $\displaystyle\lim_{x\to x_0}f(x)=f(x_0)$.
\item \textbf{Théorème des valeurs intermédiaires} : si $f$ est continue sur $[a\,;b]$, alors pour tout réel $k$ compris entre $f(a)$ et $f(b)$, l'équation $f(x)=k$ admet \textbf{au moins une} solution dans $[a\,;b]$. Si de plus $f$ est strictement monotone, cette solution est \textbf{unique}.
\item Si $f$ est une bijection de $I$ sur $J$, sa \cle{réciproque} $f^{-1}$ est l'unique fonction de $J$ dans $I$ telle que $f^{-1}(y)=x \iff y=f(x)$. On a $f^{-1}\circ f=\mathrm{Id}_I$ et $f\circ f^{-1}=\mathrm{Id}_J$. Dans un repère orthonormé, les courbes de $f$ et de $f^{-1}$ sont \textbf{symétriques par rapport à la droite d'équation $y=x$} (première bissectrice).
\end{enumerate}
\end{corrige}

\begin{corrige}[Exercice 4 — Calculs directs de limites]
\begin{enumerate}
\item On fait apparaître la limite remarquable :
\[\dfrac{\sin(3x)}{x}=3\,\dfrac{\sin(3x)}{3x}.\]
En posant $u=3x$, $u\to 0$ quand $x\to 0$, donc $\displaystyle\lim_{x\to 0}\dfrac{\sin(3x)}{x}=3\times 1=3$.
\item $\dfrac{\tan x}{x}=\dfrac{\sin x}{x}\times\dfrac{1}{\cos x}\xrightarrow[x\to 0]{}1\times\dfrac{1}{1}=1$.
\item Forme indéterminée « $\infty-\infty$ » : on rationalise.
\[\sqrt{x^2+4}-x=\dfrac{(x^2+4)-x^2}{\sqrt{x^2+4}+x}=\dfrac{4}{\sqrt{x^2+4}+x}\xrightarrow[x\to+\infty]{}0.\]
\item On reconnaît un taux d'accroissement, ou on rationalise :
\[\dfrac{\sqrt{x}-2}{x-4}=\dfrac{\sqrt{x}-2}{(\sqrt{x}-2)(\sqrt{x}+2)}=\dfrac{1}{\sqrt{x}+2}\xrightarrow[x\to 4]{}\dfrac{1}{4}.\]
(C'est aussi $f'(4)$ pour $f=\sqrt{\cdot}$, avec $f'(x)=\frac{1}{2\sqrt{x}}$.)
\end{enumerate}
\end{corrige}

\courssub{Je m'entraîne}

\begin{corrige}[Exercice 5 — Application du TVI]
\begin{enumerate}
\item $f$ est une fonction polynôme, dérivable sur $\mathbb{R}$, et $f'(x)=3x^2+1>0$ pour tout $x$ (somme de termes strictement positifs). Donc $f$ est \textbf{strictement croissante} sur $\mathbb{R}$. De plus $\displaystyle\lim_{x\to-\infty}f(x)=-\infty$ et $\displaystyle\lim_{x\to+\infty}f(x)=+\infty$.
\begin{center}\begin{tabular}{|c|ccc|}\hline $x$ & $-\infty$ & & $+\infty$ \\ \hline $f'(x)$ & & $+$ & \\ \hline $f$ & $-\infty$ & $\nearrow$ & $+\infty$ \\ \hline\end{tabular}\end{center}
\item $f$ est \textbf{continue} (polynôme) et \textbf{strictement croissante} sur $\mathbb{R}$, et $0\in\left]-\infty\,;+\infty\right[$. D'après le théorème de la bijection, l'équation $f(x)=0$ admet une \textbf{unique} solution $\alpha$ dans $\mathbb{R}$.
\item $f(1)=1+1-3=-1<0$ et $f(2)=8+2-3=7>0$. Comme $f$ est strictement croissante et change de signe, $\alpha\in[1\,;2]$.\\
À la calculatrice : $f(1{,}2)=1{,}728+1{,}2-3=-0{,}072<0$ et $f(1{,}3)=2{,}197+1{,}3-3=0{,}497>0$.\\
Donc $\boxed{1{,}2<\alpha<1{,}3}$.
\end{enumerate}
\end{corrige}

\begin{corrige}[Exercice 6 — Limites trigonométriques]
\begin{enumerate}
\item Avec $1-\cos x=2\sin^2\!\left(\dfrac{x}{2}\right)$ :
\[\dfrac{1-\cos x}{x^2}=\dfrac{2\sin^2\!\left(\frac{x}{2}\right)}{x^2}=\dfrac{1}{2}\left(\dfrac{\sin\left(\frac{x}{2}\right)}{\frac{x}{2}}\right)^2.\]
En posant $u=\dfrac{x}{2}\to 0$, on obtient $\displaystyle\lim_{x\to 0}f(x)=\dfrac{1}{2}\times 1^2=\dfrac{1}{2}$.
\item $f$ n'est pas définie en $0$ mais admet une limite finie $\dfrac{1}{2}$ en $0$ : elle est donc \textbf{prolongeable par continuité} en $0$, par $\tilde{f}(0)=\dfrac{1}{2}$ et $\tilde{f}(x)=f(x)$ sinon.
\item On fait apparaître deux limites remarquables :
\[\dfrac{\sin(5x)}{\sin(2x)}=\dfrac{5x\cdot\frac{\sin(5x)}{5x}}{2x\cdot\frac{\sin(2x)}{2x}}=\dfrac{5}{2}\times\dfrac{\frac{\sin(5x)}{5x}}{\frac{\sin(2x)}{2x}}\xrightarrow[x\to 0]{}\dfrac{5}{2}\times\dfrac{1}{1}=\dfrac{5}{2}.\]
\end{enumerate}
\end{corrige}

\begin{corrige}[Exercice 7 — Fonction irrationnelle]
\begin{enumerate}
\item $g$ est définie lorsque $x^2-4x+3\geq 0$. Or $x^2-4x+3=(x-1)(x-3)$, trinôme positif à l'extérieur des racines. Donc
\[D_g=\left]-\infty\,;1\right]\cup\left[3\,;+\infty\right[.\]
\item $\displaystyle\lim_{x\to-\infty}g(x)=+\infty$ (car $x^2-4x+3\to+\infty$) ; $\displaystyle\lim_{x\to 1^-}g(x)=\sqrt{0}=0$ ; $\displaystyle\lim_{x\to 3^+}g(x)=0$ ; $\displaystyle\lim_{x\to+\infty}g(x)=+\infty$.
\item Le polynôme $x\mapsto x^2-4x+3$ est continu sur $\mathbb{R}$, et il est \textbf{positif} sur $D_g$ ; la fonction racine carrée est continue sur $[0\,;+\infty[$. Par composition, $g$ est \textbf{continue sur $D_g$} (continuité à droite en $3$ et à gauche en $1$).
\item Pour $x>0$ :
\[\dfrac{g(x)}{x}=\dfrac{\sqrt{x^2-4x+3}}{x}=\sqrt{\dfrac{x^2-4x+3}{x^2}}=\sqrt{1-\dfrac{4}{x}+\dfrac{3}{x^2}}\xrightarrow[x\to+\infty]{}1.\]
\textbf{Interprétation} : la courbe de $g$ a en $+\infty$ une \textbf{direction asymptotique} de pente $1$ (elle se rapproche d'une droite parallèle à $y=x$).
\end{enumerate}
\end{corrige}

\begin{corrige}[Exercice 8 — Bijection et réciproque]
\begin{enumerate}
\item $h$ est dérivable sur $[0\,;+\infty[$ et $h'(x)=2x+2>0$ pour $x\geq 0$ : $h$ est \textbf{continue et strictement croissante}. Comme $h(0)=0$ et $\displaystyle\lim_{x\to+\infty}h(x)=+\infty$, $h$ réalise une bijection de $[0\,;+\infty[$ sur $J=[0\,;+\infty[$.
\item Pour $y\geq 0$, résolvons $y=x^2+2x$ avec $x\geq 0$ : $(x+1)^2=y+1$, donc $x+1=\sqrt{y+1}$ (on garde la racine positive car $x+1\geq 1>0$), d'où
\[h^{-1}(y)=\sqrt{y+1}-1,\qquad\text{soit } h^{-1}(x)=\sqrt{x+1}-1\ \text{sur } [0\,;+\infty[.\]
\item $h(1)=1+2=3$.
\textbf{Première méthode} (formule de la réciproque) : $(h^{-1})'(3)=\dfrac{1}{h'(1)}=\dfrac{1}{2\times 1+2}=\dfrac{1}{4}$.\\
\textbf{Deuxième méthode} (dérivation directe) : $(h^{-1})'(x)=\dfrac{1}{2\sqrt{x+1}}$, donc $(h^{-1})'(3)=\dfrac{1}{2\sqrt{4}}=\dfrac{1}{4}$. Les deux résultats concordent.
\item On trace $y=x^2+2x$ sur $[0\,;+\infty[$, la droite $y=x$, puis la courbe symétrique.
\begin{popfigure}
\begin{center}
\begin{tikzpicture}[scale=1.1]
\draw[->,popmuted] (-0.5,0) -- (4.5,0) node[below left]{$x$};
\draw[->,popmuted] (0,-0.5) -- (0,4.5) node[below left]{$y$};
\draw[popmuted,dashed] (-0.4,-0.4) -- (4.3,4.3) node[above left,font=\small]{$y=x$};
\draw[popblue,thick,domain=0:1.55,samples=100] plot (\x,{\x*\x+2*\x}) node[above right,font=\small]{$\mathcal{C}_h$};
\draw[poporange,thick,domain=0:4.3,samples=100] plot (\x,{sqrt(max(0,\x+1))-1}) node[right,font=\small]{$\mathcal{C}_{h^{-1}}$};
\fill[popink] (1,3) circle (1.5pt) node[above left,font=\small]{$(1\,;3)$};
\fill[popink] (3,1) circle (1.5pt) node[below right,font=\small]{$(3\,;1)$};
\draw (1,-0.08)--(1,0.08) node[below,font=\small]{$1$};
\draw (-0.08,1)--(0.08,1) node[left,font=\small]{$1$};
\end{tikzpicture}
\end{center}
\legende{Courbes de $h$ et de $h^{-1}$, symétriques par rapport à la première bissectrice.}
\end{popfigure}
\end{enumerate}
\end{corrige}

\courssub{Je me prépare au Probatoire}

\begin{corrige}[Exercice 9 — Type Probatoire 1 : analyse pure]
\textbf{1.a)} Pour tout réel $x$, $x^2+4>x^2$, donc $\sqrt{x^2+4}>\sqrt{x^2}=|x|\geq x$, d'où $\sqrt{x^2+4}-x>0$. Le dénominateur ne s'annule jamais : $f$ est \textbf{définie sur $\mathbb{R}$}. De plus $x\mapsto\sqrt{x^2+4}-x$ est continue (composée et somme de fonctions continues) et ne s'annule pas, donc $f$, quotient de fonctions continues, est \textbf{continue sur $\mathbb{R}$}.

\textbf{1.b)} On multiplie par la quantité conjuguée, en utilisant $(\sqrt{x^2+4}-x)(\sqrt{x^2+4}+x)=x^2+4-x^2=4$ :
\[f(x)=(2x+1)\times\dfrac{\sqrt{x^2+4}+x}{(\sqrt{x^2+4}-x)(\sqrt{x^2+4}+x)}=(2x+1)\,\dfrac{\sqrt{x^2+4}+x}{4}.\]

\textbf{2.a)} En $+\infty$ : $\sqrt{x^2+4}+x\sim 2x\to+\infty$ et $2x+1\to+\infty$, donc par produit $\displaystyle\lim_{x\to+\infty}f(x)=+\infty$.

\textbf{2.b)} En $-\infty$ : $\sqrt{x^2+4}+x=\dfrac{4}{\sqrt{x^2+4}-x}\to 0^+$ (car $-x\to+\infty$). Plus précisément, pour $x<0$, $\sqrt{x^2+4}+x=\dfrac{4}{\sqrt{x^2+4}-x}\sim\dfrac{4}{-2x}=-\dfrac{2}{x}$, donc
\[f(x)\sim (2x)\times\dfrac{-2/x}{4}=-1.\]
Ainsi $\displaystyle\lim_{x\to-\infty}f(x)=-1$, c'est-à-dire $\displaystyle\lim_{x\to-\infty}\big(f(x)-(-1)\big)=0$.\\
\textbf{Interprétation} : la droite d'équation $y=-1$ est \textbf{asymptote horizontale} à $(\mathcal{C})$ en $-\infty$.

\textbf{3.a)} $f=\dfrac{N}{D}$ avec $N=2x+1$, $D=\sqrt{x^2+4}-x$. On a $N'=2$ et
\[D'(x)=\dfrac{x}{\sqrt{x^2+4}}-1=\dfrac{x-\sqrt{x^2+4}}{\sqrt{x^2+4}}=-\dfrac{D(x)}{\sqrt{x^2+4}}.\]
Donc
\[f'(x)=\dfrac{N'D-ND'}{D^2}=\dfrac{2D+\dfrac{N\,D}{\sqrt{x^2+4}}}{D^2}=\dfrac{2\sqrt{x^2+4}+2x+1}{D(x)\sqrt{x^2+4}}.\]
Le dénominateur est strictement positif sur $\mathbb{R}$, donc $f'(x)$ a le \textbf{signe de $\varphi(x)=2\sqrt{x^2+4}+2x+1$}.

\textbf{3.b)} Étudions $\varphi$ : $\varphi'(x)=\dfrac{2x}{\sqrt{x^2+4}}+2>0$ car $|x|<\sqrt{x^2+4}$. De plus $\displaystyle\lim_{x\to-\infty}\varphi(x)=1$ (car $2\sqrt{x^2+4}\sim -2x$). $\varphi$ étant strictement croissante et de limite $1$ en $-\infty$, on a $\varphi(x)>1>0$ pour tout $x$. Donc $f'(x)>0$ : $f$ est \textbf{strictement croissante} sur $\mathbb{R}$.
\begin{center}\begin{tabular}{|c|ccc|}\hline $x$ & $-\infty$ & & $+\infty$ \\ \hline $f'(x)$ & & $+$ & \\ \hline $f$ & $-1$ & $\nearrow$ & $+\infty$ \\ \hline\end{tabular}\end{center}

\textbf{4.} $f$ est continue et strictement croissante sur $\mathbb{R}$, et $3\in\left]-1\,;+\infty\right[$ : d'après le théorème de la bijection, l'équation $f(x)=3$ admet une \textbf{unique} solution $\alpha$.\\
Avec la forme rationalisée : $f(1)=\dfrac{3(\sqrt{5}+1)}{4}\approx 2{,}43<3$ ; $f(1{,}1)\approx 2{,}71<3$ ; $f(1{,}2)=\dfrac{3{,}4(\sqrt{5{,}44}+1{,}2)}{4}\approx 3{,}00>3$. Donc $\boxed{1{,}1<\alpha<1{,}2}$.

\textbf{5.} Éléments remarquables : asymptote $y=-1$ en $-\infty$ ; $f(0)=\dfrac{1}{2}$ ; intersection avec l'axe des abscisses en $x=-\dfrac{1}{2}$ ; passage par $\left(1\,;\frac{3(\sqrt5+1)}{4}\right)$.
\begin{popfigure}
\begin{center}
\begin{tikzpicture}[scale=0.55]
\draw[->,popmuted] (-6.5,0) -- (4.5,0) node[below left]{$x$};
\draw[->,popmuted] (0,-2) -- (0,6) node[below left]{$y$};
\draw[popgreen,dashed] (-6.3,-1) -- (4.3,-1) node[right,font=\small]{$y=-1$};
\draw[popblue,thick,domain=-6.2:1.9,samples=200] plot (\x,{(2*\x+1)*(sqrt(max(0,\x*\x+4))+\x)/4});
\fill[popink] (0,0.5) circle (2pt);
\fill[popink] (-0.5,0) circle (2pt);
\draw (1,-0.12)--(1,0.12) node[below,font=\small]{$1$};
\draw (-0.12,1)--(0.12,1) node[left,font=\small]{$1$};
\node[popblue,font=\small] at (-4.5,3.5) {$(\mathcal{C})$};
\end{tikzpicture}
\end{center}
\legende{Courbe de $f$ : croissante, asymptote $y=-1$ en $-\infty$.}
\end{popfigure}

\textbf{Barème indicatif (sur 10)} : question 1 : 2 pts ; question 2 : 2 pts ; question 3 : 2,5 pts ; question 4 : 1,5 pt ; question 5 : 2 pts.

\textbf{Points de vigilance} :
\begin{itemize}
\item Pour la continuité, citer les théorèmes généraux (somme, composée, quotient dont le dénominateur ne s'annule pas).
\item En $-\infty$, il est interdit d'écrire « $\sqrt{x^2+4}+x$ est une somme » sans rationaliser : c'est une forme indéterminée. Penser à $\sqrt{x^2}=|x|=-x$ pour $x<0$.
\item Pour l'unicité de $\alpha$, la \textbf{stricte monotonie} doit être citée explicitement ; le TVI seul ne donne que l'existence.
\end{itemize}
\end{corrige}

\begin{corrige}[Exercice 10 — Type Probatoire 2 : trigonométrie et modélisation]
\textbf{1.} $u(0)=10-8e^{0}+2\sin 0=10-8=2$ (V).\\
$u(0{,}1)=10-8e^{-0{,}2}+2\sin(10\pi)=10-8\times 0{,}8187+0\approx 3{,}45$ (V).

\textbf{2.a)} $t\mapsto e^{-2t}$ et $t\mapsto\sin(100\pi t)$ sont continues sur $\mathbb{R}$ (composées de fonctions continues), donc $u$, somme de fonctions continues, est \textbf{continue sur $[0\,;0{,}1]$}.

\textbf{2.b)} $u'(t)=-8\times(-2)e^{-2t}+2\times 100\pi\cos(100\pi t)=16e^{-2t}+200\pi\cos(100\pi t)$.

\textbf{2.c)} Le terme $2\sin(100\pi t)$ est périodique de période $0{,}02$~s. Sur l'intervalle $[0\,;0{,}005]$, $100\pi t\in[0\,;\pi/2]$, donc $\cos(100\pi t)\geq 0$ et $u'(t)>0$ : $u$ est \textbf{strictement croissante} sur $[0\,;0{,}005]$.\\
$u(0)=2$ et $u(0{,}005)=10-8e^{-0{,}01}+2\sin(\pi/2)\approx 10-7{,}92+2=4{,}08$.
\begin{center}\begin{tabular}{|c|ccc|}\hline $t$ & $0$ & & $0{,}005$ \\ \hline $u$ & $2$ & $\nearrow$ & $\approx 4{,}08$ \\ \hline\end{tabular}\end{center}

\textbf{3.a)} La valeur $u_0=2{,}2$ V appartient à l'intervalle image $[2\,;4{,}08]$. $u$ est continue et strictement croissante sur $[0\,;0{,}005]$, donc d'après le théorème de la bijection il existe un \textbf{unique} $t_0\in[0\,;0{,}005]$ tel que $u(t_0)=2{,}2$.

\textbf{3.b)} À la calculatrice (mode \textbf{radian}) : $u(0{,}0003)\approx 2{,}19<2{,}2$ et $u(0{,}0004)\approx 2{,}26>2{,}2$. Donc $\boxed{0{,}0003<t_0<0{,}0004}$ (en secondes).

\textbf{3.c)} Pour le technicien : la tension aux bornes du condensateur atteint la valeur consignée de $2{,}2$~V environ \SI{0,3}{ms} à \SI{0,4}{ms} après la mise sous tension.

\textbf{4.} $\displaystyle\lim_{t\to 0}\dfrac{u(t)-u(0)}{t}=u'(0)=16e^{0}+200\pi\cos 0=16+200\pi\approx 644{,}3$ V/s.\\
Ce nombre est le \textbf{taux d'accroissement initial de la tension} : la « vitesse » de charge du condensateur à l'instant $t=0$.

\textbf{Barème indicatif (sur 10)} : question 1 : 1,5 pt ; question 2 : 3 pts ; question 3 : 3,5 pts ; question 4 : 2 pts.

\textbf{Points de vigilance} :
\begin{itemize}
\item La calculatrice doit être en \textbf{mode radian} pour évaluer $\sin(100\pi t)$.
\item Dans la dérivée, ne pas oublier les coefficients issus de la composition : $(e^{-2t})'=-2e^{-2t}$ et $(\sin(100\pi t))'=100\pi\cos(100\pi t)$.
\item La fonction $u$ n'est \textbf{pas} monotone sur $[0\,;0{,}1]$ à cause de l'oscillation sinusoïdale ; il faut restreindre l'intervalle (ici $[0\,;0{,}005]$) pour appliquer le théorème de la bijection.
\item Au TVI, toujours vérifier et écrire que la valeur cible appartient à l'intervalle image $[u(0)\,;u(0{,}005)]$.
\end{itemize}
\end{corrige}

\begin{corrige}[Exercice 11 — Type Probatoire 3 : optimisation en atelier]
\textbf{1.a)} On découpe un carré de côté $x$ à \emph{chaque} coin : il faut $x\geq 0$ et $2x\leq 40$ (les deux découpes opposées ne doivent pas se chevaucher), donc $x\in[0\,;20]$.

\textbf{1.b)} Après pliage, le bac est un pavé droit de longueur et largeur $40-2x$ et de hauteur $x$ :
\[V(x)=x(40-2x)^2.\]
En développant : $(40-2x)^2=1600-160x+4x^2$, donc
\[V(x)=x(4x^2-160x+1600)=4x^3-160x^2+1600x.\]

\textbf{2.a)} $V$ est une fonction \textbf{polynôme} : elle est continue et dérivable sur $\mathbb{R}$, donc en particulier sur $[0\,;20]$.

\textbf{2.b)} $V'(x)=12x^2-320x+1600$. Vérifions la factorisation :
\[12\left(x-\dfrac{20}{3}\right)(x-20)=4(3x-20)(x-20)=4(3x^2-80x+400)=12x^2-320x+1600.\ \checkmark\]

\textbf{2.c)} $V'(x)$ est un trinôme de racines $x=\dfrac{20}{3}$ et $x=20$, positif à l'extérieur des racines. Sur $[0\,;20]$ : $V'(x)>0$ sur $\left[0\,;\frac{20}{3}\right[$ et $V'(x)<0$ sur $\left]\frac{20}{3}\,;20\right]$. De plus $V(0)=0$, $V(20)=0$ et
\[V\left(\dfrac{20}{3}\right)=\dfrac{20}{3}\left(40-\dfrac{40}{3}\right)^2=\dfrac{20}{3}\times\left(\dfrac{80}{3}\right)^2=\dfrac{128000}{27}\approx 4740{,}7.\]
\begin{center}\begin{tabular}{|c|ccccc|}\hline $x$ & $0$ & & $\frac{20}{3}$ & & $20$ \\ \hline $V'(x)$ & & $+$ & $0$ & $-$ & \\ \hline $V$ & $0$ & $\nearrow$ & $\frac{128000}{27}$ & $\searrow$ & $0$ \\ \hline\end{tabular}\end{center}

\textbf{3.} Le volume est maximal pour $x=\dfrac{20}{3}\approx 6{,}67$ cm, et
\[V_{\max}=\dfrac{128000}{27}\approx \SI{4741}{cm^3}\quad(\text{soit environ } 4{,}74\ \text{litres}).\]

\textbf{4.a)} $V$ est continue et strictement croissante sur $\left[0\,;\frac{20}{3}\right]$, et $1500\in\left[V(0)\,;V\!\left(\frac{20}{3}\right)\right]=[0\,;\approx 4741]$. D'après le théorème des valeurs intermédiaires (version bijection), il existe une \textbf{unique} valeur $x_0\in\left]0\,;\dfrac{20}{3}\right[$ telle que $V(x_0)=1500$.

\textbf{4.b)} $V(1)=1\times 38^2=1444<1500$ et $V(1{,}1)=1{,}1\times 37{,}8^2\approx 1571{,}7>1500$. Donc $\boxed{1{,}0<x_0<1{,}1}$ (en cm).

\textbf{5.} \textbf{Conclusion pour le tôlier} : si l'objectif est le \emph{volume maximal}, il faut découper des carrés de côté $\dfrac{20}{3}\approx 6{,}7$ cm, pour obtenir un bac d'environ \SI{4741}{cm^3}. Si le client exige exactement \SI{1500}{cm^3}, une découpe d'environ \SI{1}{cm} (entre $1{,}0$ et $1{,}1$ cm) convient ; le bac sera alors très plat, donc moins pratique : sauf contrainte impérative du client, la découpe de $6{,}7$ cm est préférable car elle exploite au mieux la feuille.

\textbf{Barème indicatif (sur 10)} : question 1 : 2 pts ; question 2 : 3 pts ; question 3 : 1,5 pt ; question 4 : 2 pts ; question 5 : 1,5 pt.

\textbf{Points de vigilance} :
\begin{itemize}
\item Justifier le domaine $[0\,;20]$ par une contrainte \emph{géométrique} ($2x\leq 40$), pas seulement par le calcul.
\item Pour vérifier une factorisation, développer et comparer — ne pas l'affirmer sans calcul.
\item Au TVI, préciser la continuité \emph{et} la monotonie de $V$ sur l'intervalle choisi, et vérifier que $1500$ appartient à l'intervalle image.
\item Bien distinguer les deux questions : le maximum (question 3) et le volume imposé (question 4) ne donnent pas la même découpe.
\end{itemize}
\end{corrige}

\newpage

\coursec{Fiche bilan}

\begin{popfigure}
\begin{tikzpicture}[mindmap, grow cyclic, every node/.style={concept, font=\small},
  root concept/.append style={concept color=popblue, font=\bfseries\small, minimum size=3.2cm},
  level 1/.append style={level distance=3.6cm, sibling angle=60},
  level 1 concept/.append style={minimum size=2.5cm, font=\scriptsize}]
\node[root concept]{Fonctions numériques d'une variable réelle}
  child[concept color=poporange]{ node {Continuité : pas de saut, $f(a)=\lim\limits_{x\to a}f(x)$} }
  child[concept color=popgreen]{ node {TVI : $f(a)\cdot f(b)<0$ donne une solution} }
  child[concept color=poppurple]{ node {Continue + strictement monotone $\Rightarrow$ bijection et $f^{-1}$} }
  child[concept color=poppink]{ node {Trigo : $\lim\limits_{x\to0}\dfrac{\sin x}{x}=1$ ; $\lim\limits_{x\to0}\dfrac{1-\cos x}{x^2}=\dfrac12$} }
  child[concept color=popgold]{ node {Irrationnelles : rationaliser par le conjugué} }
  child[concept color=popblueL!80!popdark]{ node {Représentation graphique : domaine, limites, variations, asymptotes} };
\end{tikzpicture}
\legende{Carte mentale de la leçon : les six idées forces.}
\end{popfigure}

\begin{bilan}
\textbf{Définitions clés}
\begin{itemize}
\item $f$ est \cle{continue} en $a$ si $\lim\limits_{x\to a}f(x)=f(a)$ ; continue sur un intervalle $I$ si elle est continue en tout point de $I$ (courbe « sans lever le crayon »).
\item $f$ est \cle{strictement monotone} sur $I$ si elle est strictement croissante ou strictement décroissante sur $I$.
\item Si $f$ est continue et strictement monotone sur $I$, elle réalise une \cle{bijection} de $I$ sur $J=f(I)$ ; sa \cle{réciproque} $f^{-1}$ est continue, de même sens de variation, et $f^{-1}(f(x))=x$.
\end{itemize}

\textbf{Théorèmes et formules à connaître par c\oe ur}
\begin{center}
\begin{tabularx}{\linewidth}{|l|X|}
\hline
\rowcolor{popblueL} \textbf{Résultat} & \textbf{Énoncé} \\
\hline
Théorème des valeurs intermédiaires & $f$ continue sur $[a;b]$ : pour tout $k$ entre $f(a)$ et $f(b)$, il existe $c\in[a;b]$ avec $f(c)=k$. \\
\hline
Corollaire (bijection) & $f$ continue et strictement monotone sur $[a;b]$, $k$ entre $f(a)$ et $f(b)$ : l'équation $f(x)=k$ admet une \textbf{unique} solution. \\
\hline
Limites trigonométriques & $\lim\limits_{x\to0}\dfrac{\sin x}{x}=1$ ;\quad $\lim\limits_{x\to0}\dfrac{\tan x}{x}=1$ ;\quad $\lim\limits_{x\to0}\dfrac{1-\cos x}{x^2}=\dfrac12$. \\
\hline
Irrationnelles & $\sqrt{a}-\sqrt{b}=\dfrac{a-b}{\sqrt{a}+\sqrt{b}}$ (multiplier par la quantité conjuguée). \\
\hline
Asymptotes & $\lim\limits_{x\to+\infty}f(x)=b$ : asymptote horizontale $y=b$ ;\quad $\lim\limits_{x\to a}f(x)=\pm\infty$ : asymptote verticale $x=a$. \\
\hline
\end{tabularx}
\end{center}

\textbf{Méthodes express}
\begin{itemize}
\item \cle{Localiser une solution} de $f(x)=0$ : vérifier la continuité, calculer $f(a)$ et $f(b)$ de signes contraires, conclure par le TVI, puis encadrer par balayage (table de valeurs).
\item \cle{Lever une indétermination} « $\infty-\infty$ » avec des racines : multiplier numérateur et dénominateur par l'expression conjuguée.
\item \cle{Lever une indétermination} trigonométrique en $0$ : faire apparaître $\dfrac{\sin u}{u}$ avec $u\to0$, ou encadrer avec $-1\leq\sin x\leq1$.
\item \cle{Étude complète} : domaine $\to$ limites aux bornes $\to$ continuité/dérivée $\to$ tableau de variations $\to$ asymptotes $\to$ points remarquables $\to$ tracé.
\end{itemize}
\end{bilan}

\begin{aretenir}[Checklist avant le Probatoire]
\begin{itemize}
\item[$\square$] Je sais vérifier la continuité d'une fonction en un point (limite = image).
\item[$\square$] Je sais énoncer le théorème des valeurs intermédiaires avec toutes ses hypothèses.
\item[$\square$] Je sais justifier l'existence \emph{et l'unicité} d'une solution de $f(x)=k$ (continuité + stricte monotonie).
\item[$\square$] Je sais encadrer une solution à $10^{-1}$ près par balayage.
\item[$\square$] Je connais par c\oe ur $\lim\limits_{x\to0}\dfrac{\sin x}{x}=1$ et $\lim\limits_{x\to0}\dfrac{1-\cos x}{x^2}=\dfrac12$.
\item[$\square$] Je sais utiliser l'encadrement $-1\leq\cos x\leq1$ pour les limites en l'infini.
\item[$\square$] Je sais rationaliser $\sqrt{a}-\sqrt{b}$ grâce à la quantité conjuguée.
\item[$\square$] Je sais déterminer le domaine d'une fonction avec racine carrée (quantité sous la racine $\geq0$).
\item[$\square$] Je sais construire un tableau de variations complet et cohérent avec les limites.
\item[$\square$] Je sais repérer et tracer les asymptotes horizontales et verticales.
\item[$\square$] Je sais relier les courbes de $f$ et de $f^{-1}$ (symétrie par rapport à $y=x$).
\item[$\square$] Je rédige mes justifications en citant les hypothèses avant de conclure.
\end{itemize}
\end{aretenir}

\findecours

\end{document}
